% Lectura y comentario : libertad de expresión y delitos de prensa ; análisis de publicidad y fake news — Espagnol Tle A — Cours signé OnBuch+
% Programme officiel MINESEC. Compiler avec : tectonic E26-libertad-expresion-delitos-prensa-fake-news.tex
\newcommand{\DOCMATIERE}{Espagnol}
\newcommand{\DOCNIVEAU}{Tle A}
\newcommand{\DOCMODULE}{Módulo 5 : Médias, communication et culture de la paix}
\newcommand{\DOCLECON}{E26}
\newcommand{\DOCTITRE}{Lectura y comentario : libertad de expresión y delitos de prensa ; análisis de publicidad y fake news}
% OnBuch+ — préambule « cours signés » ENRICHI (compilé avec tectonic / XeLaTeX)
% Système visuel « Pop » d'OnBuch+ : fond crème (jamais blanc), bordures encre
% épaisses, ombres « dures » décalées, titres Archivo Black en capitales.
% Version MATHÉMATIQUES : graphes (pgfplots/tikz), boîtes pédagogiques
% (définition, propriété, méthode, attention, le savais-tu, exercice résolu,
% exercice, TP, bilan), page de garde et signature OnBuch+ ancrée en bas.
%
% Macros attendues dans main.tex AVANT \input{preamble} :
%   \DOCMATIERE  \DOCNIVEAU  \DOCMODULE  \DOCLECON (numéro)  \DOCTITRE
\documentclass[11pt]{article}

\usepackage{fontspec}
\usepackage[spanish,french]{babel} % français = langue principale ; l'espagnol via \esp{..} / espagnol
\usepackage[a4paper,margin=2cm,top=3.4cm,bottom=2.8cm,headheight=1.4cm,footskip=1.3cm]{geometry}
\usepackage{amsmath,amssymb,amsfonts}
\usepackage{mathtools} % \xmapsto, \coloneqq... souvent utilisés par les modèles
\usepackage{cancel} % simplifications barrées (bilans de forces)
% \abs{z} (module, valeur absolue) et \norm{u} (norme), utilisés par les modèles
\providecommand{\abs}{}\let\abs\relax\DeclarePairedDelimiter{\abs}{\lvert}{\rvert}
\providecommand{\norm}{}\let\norm\relax\DeclarePairedDelimiter{\norm}{\lVert}{\rVert}
\usepackage[table]{xcolor} % [table] : fournit aussi \rowcolors (alternance de lignes)
\usepackage{pagecolor}
\usepackage{tikz}
\usetikzlibrary{arrows.meta,positioning,calc,shapes.geometric,shapes.misc,shapes.arrows,decorations.pathmorphing,decorations.pathreplacing,decorations.markings,patterns,shadows,mindmap,trees,fit,backgrounds,matrix,intersections,angles,quotes}
\usepackage{pgfplots}
\usepackage[version=4]{mhchem} % équations nucléaires \ce{^{235}_{92}U}
\usepackage[european,siunitx]{circuitikz} % schémas électriques
\pgfplotsset{compat=1.18}
\usepgfplotslibrary{fillbetween,groupplots} % aires entre courbes (calcul intégral)
\usepackage{enumitem}
\usepackage{fancyhdr}
% [most] charge toutes les bibliothèques tcolorbox (listings, minted, raster,
% documentation...) et gonfle inutilement la mémoire TeX sur un document long
% (~300 boîtes) : on ne charge que ce qui est réellement utilisé.
\usepackage[skins,breakable]{tcolorbox}
\usepackage{graphicx}
\usepackage{colortbl}
\usepackage{booktabs}
\usepackage{tabularx}
\usepackage{diagbox} % case d'en-tête diagonale des tableaux à double entrée
% Colonnes extensibles centrée / gauche / droite (C, L, R), souvent utilisées par les modèles
\newcolumntype{C}{>{\centering\arraybackslash}X}
\newcolumntype{L}{>{\raggedright\arraybackslash}X}
\newcolumntype{R}{>{\raggedleft\arraybackslash}X}
\usepackage{array}
\usepackage{multicol}
\usepackage{siunitx}
% parse-numbers=false : accepte aussi \SI{2,5 \times 10^{-3}}{...}
\sisetup{locale=FR,output-decimal-marker={,},group-minimum-digits=4,parse-numbers=false}
\DeclareSIUnit\ohm{\ensuremath{\symup{\Omega}}} % U+2126 absent de la police
\DeclareSIUnit\division{div} % réglages d'oscilloscope (ms/div, V/div)
\DeclareSIUnit\tr{tr} % tours (vitesse de rotation, ex. \SI{1500}{\tr\per\minute})
\DeclareSIUnit\molar{M} % concentration molaire (ex. \SI{3}{\molar}, \SI{1}{\micro\molar})
\DeclareSIUnit\an{an} % année (le modèle confond parfois avec \year, un compteur TeX) — ex. \SI{5}{\kilo\meter\per\an}
\DeclareSIUnit\FCFA{FCFA} % franc CFA (ex. \SI{500}{\FCFA\per\kilo\gram})
\DeclareSIUnit\degreeBrix{\textdegree Brix} % teneur en sucre d'un sirop/jus (réfractométrie)
\DeclareSIUnit\month{mois} % le modèle utilise parfois \month, un compteur TeX, comme unité de durée
\usepackage{qrcode}
% \degree : commande fréquemment utilisée par les modèles pour un angle ou une
% température en dehors de \SI{}{} (non fournie par les paquets chargés).
\providecommand{\degree}{^{\circ}}
% \qed (fin de démonstration), utilisé par les modèles sans amsthm
\providecommand{\qed}{\hfill$\square$}
% \barre : double barre des valeurs interdites dans un tableau de variations (tkz-tab)
\providecommand{\barre}{\ensuremath{\|}}
% environnement proof (amsthm non chargé) utilisé par les modèles
\ifdefined\proof\else\newenvironment{proof}[1][Démonstration]{\par\noindent\textit{#1.}\ }{\qed\par}\fi
\usepackage{lastpage}
\usepackage{hyperref}
\hypersetup{hidelinks}

% ============================================================
% Palette « Pop » OnBuch+ — mêmes hex que la classe P (plus_ui.dart)
% ============================================================
\definecolor{poporange}{HTML}{F59321}
\definecolor{popdark}{HTML}{E07A0C}
\definecolor{popcream}{HTML}{FFF6E8}
\definecolor{popink}{HTML}{1C1714}
\definecolor{poppurple}{HTML}{7A5AE0}
\definecolor{popgreen}{HTML}{2E9E6A}
\definecolor{poppink}{HTML}{E0457B}
\definecolor{popblue}{HTML}{2F7DE1}
\definecolor{popgold}{HTML}{C98A12}
\definecolor{popmuted}{HTML}{8A7F76}
\definecolor{popline}{HTML}{EADFD2}
\definecolor{popgray}{HTML}{8A7F76}
\colorlet{popgrey}{popgray}
% Alias tolérants (le modèle nomme parfois les couleurs différemment)
\colorlet{popwhite}{white}
\colorlet{popblack}{popink}
\colorlet{popred}{poppink}
\colorlet{popyellow}{popgold}
% Teintes claires (fonds de tableaux/encadrés) — le modèle nomme parfois
% les couleurs avec un suffixe L (ex. popblueL) sans les avoir définies.
\colorlet{poporangeL}{poporange!20}
\colorlet{popdarkL}{popdark!20}
\colorlet{poppurpleL}{poppurple!20}
\colorlet{popgreenL}{popgreen!20}
\colorlet{poppinkL}{poppink!20}
\colorlet{popblueL}{popblue!20}
\colorlet{popgoldL}{popgold!20}
\colorlet{popredL}{popred!20}
\colorlet{popgrayL}{popgray!20}
% Teintes claires pour fonds de tableaux / zones
\colorlet{popblueL}{popblue!10!white}
\colorlet{poporangeL}{poporange!14!white}
\colorlet{popgreenL}{popgreen!12!white}
\colorlet{poppurpleL}{poppurple!12!white}
\colorlet{poppinkL}{poppink!12!white}
\colorlet{popgoldL}{popgold!14!white}

\pagecolor{popcream}

% ============================================================
% Typographie
% ============================================================
\defaultfontfeatures{Ligatures=TeX}
\setmainfont{PlusJakartaSans-Medium.ttf}[Path=../../../fonts/,BoldFont=PlusJakartaSans-Bold.ttf]
% Maths en sans-serif (Fira Math) avec chiffres et lettres droites de Plus
% Jakarta, pour que les formules se fondent dans le texte.
\usepackage[math-style=ISO,bold-style=ISO]{unicode-math}
\setmathfont{FiraMath-Regular.otf}[Path=../../../fonts/]
\setmathfont{PlusJakartaSans-Medium.ttf}[Path=../../../fonts/,range={up/num,up/{Latin,latin},\mathup}]
\usepackage{icomma} % virgule décimale sans espace en mode math
% Caractères mathématiques Unicode absents des polices de texte (Plus Jakarta,
% Archivo Black) : rendus par leur équivalent mathématique, en texte comme en
% formule — sinon ils s'affichent en carré vide (ex. « Arithmétique dans ℤ »).
\usepackage{newunicodechar}
\newunicodechar{ℝ}{\ensuremath{\mathbb{R}}}
\newunicodechar{ℤ}{\ensuremath{\mathbb{Z}}}
\newunicodechar{ℂ}{\ensuremath{\mathbb{C}}}
\newunicodechar{ℕ}{\ensuremath{\mathbb{N}}}
\newunicodechar{ℚ}{\ensuremath{\mathbb{Q}}}
\newunicodechar{𝔻}{\ensuremath{\mathbb{D}}}
\newunicodechar{α}{\ensuremath{\alpha}}
\newunicodechar{β}{\ensuremath{\beta}}
\newunicodechar{γ}{\ensuremath{\gamma}}
\newunicodechar{δ}{\ensuremath{\delta}}
\newunicodechar{ε}{\ensuremath{\varepsilon}}
\newunicodechar{θ}{\ensuremath{\theta}}
\newunicodechar{λ}{\ensuremath{\lambda}}
\newunicodechar{μ}{\ensuremath{\mu}}
\newunicodechar{σ}{\ensuremath{\sigma}}
\newunicodechar{τ}{\ensuremath{\tau}}
\newunicodechar{φ}{\ensuremath{\varphi}}
\newunicodechar{ψ}{\ensuremath{\psi}}
\newunicodechar{ω}{\ensuremath{\omega}}
\newunicodechar{Σ}{\ensuremath{\Sigma}}
% majuscules produites par \MakeUppercase dans les titres (α-aminés -> Α)
\newunicodechar{Α}{\ensuremath{\alpha}}
\newunicodechar{Β}{\ensuremath{\beta}}
\newunicodechar{Γ}{\ensuremath{\Gamma}}
\newunicodechar{Δ}{\ensuremath{\Delta}}
\newunicodechar{Ε}{\ensuremath{\varepsilon}}
\newunicodechar{Θ}{\ensuremath{\Theta}}
\newunicodechar{Λ}{\ensuremath{\Lambda}}
\newunicodechar{Μ}{\ensuremath{\mu}}
\newunicodechar{Τ}{\ensuremath{\tau}}
\newunicodechar{Φ}{\ensuremath{\Phi}}
\newunicodechar{Ψ}{\ensuremath{\Psi}}
\newunicodechar{Ω}{\ensuremath{\Omega}}
\newunicodechar{↦}{\ensuremath{\mapsto}}
\newunicodechar{∈}{\ensuremath{\in}}
\newunicodechar{∉}{\ensuremath{\notin}}
\newunicodechar{∧}{\ensuremath{\wedge}}
\newunicodechar{⇔}{\ensuremath{\Leftrightarrow}}
\newunicodechar{⟺}{\ensuremath{\Longleftrightarrow}}
\newunicodechar{⇒}{\ensuremath{\Rightarrow}}
\newunicodechar{⟹}{\ensuremath{\Longrightarrow}}
\newunicodechar{∘}{\ensuremath{\circ}}
\newunicodechar{∪}{\ensuremath{\cup}}
\newunicodechar{∩}{\ensuremath{\cap}}
\newunicodechar{≡}{\ensuremath{\equiv}}
\newunicodechar{⊂}{\ensuremath{\subset}}
\newunicodechar{⊥}{\ensuremath{\perp}}
\newunicodechar{∃}{\ensuremath{\exists}}
\newunicodechar{∀}{\ensuremath{\forall}}
\newunicodechar{∛}{\ensuremath{\sqrt[3]{\,}}}
\newunicodechar{∜}{\ensuremath{\sqrt[4]{\,}}}
\newunicodechar{⃗}{\ensuremath{{}^{\to}}}
\newunicodechar{ⁿ}{\ifmmode^{n}\else\textsuperscript{\ensuremath{n}}\fi}
\newunicodechar{ˣ}{\ifmmode^{x}\else\textsuperscript{\ensuremath{x}}\fi}
\newunicodechar{ᵇ}{\ifmmode^{b}\else\textsuperscript{\ensuremath{b}}\fi}
\newunicodechar{ʳ}{\ifmmode^{r}\else\textsuperscript{\ensuremath{r}}\fi}
\newunicodechar{ᵘ}{\ifmmode^{u}\else\textsuperscript{\ensuremath{u}}\fi}
\newunicodechar{ʸ}{\ifmmode^{y}\else\textsuperscript{\ensuremath{y}}\fi}
\newunicodechar{ᵃ}{\ifmmode^{a}\else\textsuperscript{\ensuremath{a}}\fi}
\newunicodechar{ᵉ}{\ifmmode^{e}\else\textsuperscript{\ensuremath{e}}\fi}
\newunicodechar{ᵏ}{\ifmmode^{k}\else\textsuperscript{\ensuremath{k}}\fi}
\newunicodechar{ᵖ}{\ifmmode^{p}\else\textsuperscript{\ensuremath{p}}\fi}
\newunicodechar{ᴺ}{\ifmmode^{N}\else\textsuperscript{\ensuremath{N}}\fi}
\newunicodechar{ᶜ}{\ifmmode^{c}\else\textsuperscript{\ensuremath{c}}\fi}
\newunicodechar{ʰ}{\ifmmode^{h}\else\textsuperscript{\ensuremath{h}}\fi}
\newunicodechar{ᵠ}{\ifmmode^{q}\else\textsuperscript{\ensuremath{q}}\fi}
\newunicodechar{ₙ}{\ifmmode_{n}\else\textsubscript{\ensuremath{n}}\fi}
\newunicodechar{ₐ}{\ifmmode_{a}\else\textsubscript{\ensuremath{a}}\fi}
\newunicodechar{ᵢ}{\ifmmode_{i}\else\textsubscript{\ensuremath{i}}\fi}
\newunicodechar{ᵣ}{\ifmmode_{r}\else\textsubscript{\ensuremath{r}}\fi}
\newunicodechar{ₖ}{\ifmmode_{k}\else\textsubscript{\ensuremath{k}}\fi}
\newunicodechar{ᵦ}{\ifmmode_{\beta}\else\textsubscript{\ensuremath{\beta}}\fi}
\newunicodechar{ₓ}{\ifmmode_{x}\else\textsubscript{\ensuremath{x}}\fi}
\newfontfamily\popdisplay{ArchivoBlack-Regular.ttf}[Path=../../../fonts/]
\newfontfamily\poplogo{SpaceGrotesk-Bold.ttf}[Path=../../../fonts/]
\newfontfamily\popsemi{PlusJakartaSans-SemiBold.ttf}[Path=../../../fonts/]
\newfontfamily\popxbold{PlusJakartaSans-ExtraBold.ttf}[Path=../../../fonts/]
\newfontfamily\popxboldls{PlusJakartaSans-ExtraBold.ttf}[Path=../../../fonts/,LetterSpace=3.0]

\setlist[enumerate]{leftmargin=1.7em,itemsep=5pt,topsep=4pt}
\setlist[itemize]{leftmargin=1.7em,itemsep=4pt,topsep=4pt,label={\color{poporange}\textbullet}}
\setlength{\parindent}{0pt}
\setlength{\parskip}{5pt}
\renewcommand{\arraystretch}{1.25}

% pgfplots : style Pop par défaut
\pgfplotsset{
  every axis/.append style={axis line style={popink,line width=1pt},tick style={popink},
    grid style={popline,line width=0.5pt},label style={font=\small},tick label style={font=\footnotesize}},
  popcurve/.style={poporange,line width=1.8pt,no markers,smooth},
}
% Même style disponible dans un \draw TikZ simple (hors pgfplots)
\tikzset{popcurve/.style={poporange,line width=1.8pt}}
\tikzset{popgrid/.style={popline,very thin}}
\colorlet{popcurve}{poporange} % le nom du style est parfois utilisé comme couleur

% ============================================================
% Pill / squiggle
% ============================================================
\newcommand{\poppill}[3]{%
  \tikz[baseline=-0.55ex]\node[draw=popink,line width=1.2pt,rounded corners=2.6pt,%
    fill=#2,inner xsep=6.5pt,inner ysep=3pt]{%
    \popxboldls\fontsize{7}{7}\selectfont\color{#3}\MakeUppercase{#1}};%
}
\newcommand{\popsquiggle}[1]{%
  \tikz[baseline=-0.4ex]{
    \draw[#1,line width=2.6pt,line cap=round]
      plot [smooth] coordinates {(0,0.09) (0.34,0.01) (0.68,0.21) (1.02,0.01) (1.36,0.10)};
  }%
}
% Mot-clé surligné (marqueur orange sous le texte)
\newcommand{\cle}[1]{{\popxbold\color{popdark}#1}}

% ============================================================
% En-tête / pied
% ============================================================
\pagestyle{fancy}
\fancyhf{}
\renewcommand{\headrulewidth}{0pt}
\renewcommand{\footrulewidth}{0pt}
\lhead{\raisebox{-0.15ex}{\poplogo\fontsize{13}{13}\selectfont\color{popink}OnBuch\color{poporange}+}}
\rhead{\poppill{\DOCMATIERE\ \textbullet\ \DOCNIVEAU\ \textbullet\ Leçon \DOCLECON}{white}{popink}}
\lfoot{\popxbold\fontsize{7.6}{7.6}\selectfont\color{popmuted}\MakeUppercase{Cours signé OnBuch+}\ \textbullet\ {\popsemi\DOCTITRE}}
\rfoot{\tikz[baseline=-0.6ex]\node[rounded corners=8.5pt,fill=popink,minimum height=17pt,minimum width=17pt,inner xsep=5pt,inner sep=0]{\popxbold\fontsize{8}{8}\selectfont\color{popcream}\thepage\,/\,\pageref*{LastPage}};}

% ============================================================
% Titres
% ============================================================
\newcommand{\chaptitre}[2]{%
  \begin{tcolorbox}[enhanced,colback=poporange,colframe=popink,boxrule=2.6pt,arc=11pt,
      drop shadow=popink,left=15pt,right=15pt,top=12pt,bottom=11pt,before skip=3pt,after skip=18pt]
    \poppill{#2}{popink}{popcream}\par\vspace{9pt}
    {\raggedright\hyphenpenalty=10000\exhyphenpenalty=10000\popdisplay\fontsize{22}{24}\selectfont\color{popcream}\MakeUppercase{#1}\par}
  \end{tcolorbox}
}
% Section numérotée : pastille orange avec le numéro + titre Archivo
\newcounter{popsec}
\newcommand{\coursec}[1]{%
  \stepcounter{popsec}\setcounter{popsub}{0}%
  \par\vspace{16pt}\noindent
  \begin{minipage}{\linewidth}
  \tikz[baseline=-0.7ex]\node[draw=popink,line width=1.6pt,fill=poporange,rounded corners=4pt,minimum size=22pt,inner sep=0]{\popdisplay\fontsize{12}{12}\selectfont\color{popcream}\thepopsec};%
  \hspace{8pt}{\popdisplay\fontsize{15}{17}\selectfont\color{popink}\MakeUppercase{#1}}\par
  \vspace{4pt}\noindent{\color{poporange}\rule{36pt}{3.6pt}}
  \end{minipage}\par\nopagebreak\vspace{8pt}
}
\newcounter{popsub}
\newcommand{\courssub}[1]{%
  \stepcounter{popsub}%
  \par\vspace{9pt}\noindent{\popxbold\fontsize{11.5}{13}\selectfont\color{popink}{\color{poporange}\thepopsec.\thepopsub}\hspace{6pt}#1}\par\nopagebreak\vspace{4pt}
}

% ============================================================
% Boîtes pédagogiques — même recette : fond blanc, bordure épaisse colorée,
% ombre dure encre, badge plein. Titre optionnel [#1].
% ============================================================
\tcbset{popbase/.style={breakable,enhanced,colback=white,boxrule=2.3pt,arc=9pt,
  shadow={3pt}{-3pt}{0pt}{popink},left=14pt,right=14pt,top=11pt,bottom=11pt,before skip=11pt,after skip=15pt,
  fontupper=\popsemi\fontsize{10.6}{14.5}\selectfont\color{popink},
  fontlower=\popsemi\fontsize{10.6}{14.5}\selectfont\color{popink}}}
% Coche dessinée (le glyphe ✓ n'existe pas dans les polices du document)
\newcommand{\popcheck}[1]{\tikz[baseline=-0.5ex]\draw[#1,line width=1.6pt,line cap=round,line join=round](-0.13,0.0)--(-0.04,-0.09)--(0.13,0.1);}
\providecommand{\coche}{\popcheck{popgreen}}
\providecommand{\resultat}[1]{\par\smallskip\noindent\fcolorbox{popgreen}{popgreenL}{\parbox{\dimexpr\linewidth-2\fboxsep-2\fboxrule\relax}{\textbf{Résultat.} #1}}\par\smallskip}
\newcommand{\popboxhead}[3]{\poppill{#1}{#2}{white}\ifx\relax#3\relax\else\hspace{6pt}{\popxbold\fontsize{10.6}{12}\selectfont\color{popink}#3}\fi\par\vspace{7pt}}

\newtcolorbox{aretenir}[1][]{popbase,colframe=popgreen,before upper={\popboxhead{À retenir}{popgreen}{#1}}}
% Texte en espagnol (document de lecture, dialogue, poème, extrait) : cadre
% d'étude. Un titre optionnel (source, auteur) s'affiche dans l'en-tête.
\newtcolorbox{textebox}[1][]{popbase,colframe=popblue,before upper={\popboxhead{Texto}{popblue}{#1}\begin{otherlanguage}{spanish}},after upper={\end{otherlanguage}}}
% Marques ✓ / ✗ (forme correcte / fautive) souvent utilisées par les modèles
\providecommand{\cross}{\textcolor{poppink}{\ensuremath{\times}}}
\providecommand{\xmark}{\cross}\providecommand{\crossmark}{\cross}\providecommand{\cmark}{\ensuremath{\checkmark}}
% Ligne de réponse / justification à compléter (inventée par les modèles)
\providecommand{\justification}[1]{\par\noindent\textit{Justification~:} \dotfill\par}
% \textipa (package tipa non chargé) : la police ne couvre pas l'API -> simple passage
\providecommand{\textipa}[1]{#1}
% Mot ou phrase en espagnol à l'intérieur d'un paragraphe français
\newcommand{\esp}[1]{\ifmmode\text{\textit{#1}}\else\begingroup\selectlanguage{spanish}\itshape #1\endgroup\fi}
\newtcolorbox{exemplebox}[1][]{popbase,colframe=poporange,before upper={\popboxhead{Exemple}{poporange}{#1}}}
\newtcolorbox{definition}[1][]{popbase,colframe=popblue,before upper={\popboxhead{Définition}{popblue}{#1}}}
\newtcolorbox{propriete}[1][]{popbase,colframe=poppurple,before upper={\popboxhead{Propriété}{poppurple}{#1}}}
\newtcolorbox{methode}[1][]{popbase,colframe=popgold,before upper={\popboxhead{Méthode}{popgold}{#1}}}
\newtcolorbox{attention}[1][]{popbase,colframe=poppink,before upper={\popboxhead{Attention}{poppink}{#1}}}
\newtcolorbox{experience}[1][]{popbase,colframe=popblue,colback=popblueL,before upper={\popboxhead{Expérience}{popblue}{#1}}}
% « Le savais-tu ? » : carte violette pleine (contexte, culture, Cameroun)
\newtcolorbox{savaistu}[1][]{popbase,colback=poppurple,colframe=popink,
  fontupper=\popsemi\fontsize{10.4}{14.2}\selectfont\color{white},
  before upper={\poppill{Le savais-tu\,?}{white}{poppurple}\ifx\relax#1\relax\else\hspace{6pt}{\popxbold\color{white}#1}\fi\par\vspace{7pt}}}
% Exercice résolu : énoncé en haut, \tcblower puis la solution sur fond crème
\newcounter{popexo}\newcounter{popexr}
\newtcolorbox{exoresolu}[1][]{popbase,colframe=poporange,colbacklower=popcream,
  segmentation style={popink,line width=1.2pt,dashed},
  before upper={\stepcounter{popexr}\popboxhead{Exercice résolu \thepopexr}{poporange}{#1}},
  before lower={\poppill{Solution}{popink}{popcream}\par\vspace{6pt}}}
% Exercice (non résolu en place) — la correction est dans la partie Corrigés
\newtcolorbox{exercice}[1][]{popbase,colframe=popink,
  before upper={\stepcounter{popexo}\popboxhead{Exercice \thepopexo}{popink}{#1}}}
% Corrigé
\newtcolorbox{corrige}[1][]{popbase,colframe=popgreen,colback=popgreenL,
  before upper={\popboxhead{Corrigé}{popgreen}{#1}}}
% Objectifs / prérequis / bilan
\newtcolorbox{objectifs}{popbase,colframe=popink,colback=poporangeL,
  before upper={\popboxhead{Objectifs de la leçon}{popink}{}}}
\newtcolorbox{prerequis}{popbase,colframe=popink,colback=popblueL,
  before upper={\popboxhead{Prérequis}{popblue}{}}}
\newtcolorbox{bilan}{popbase,colframe=popink,colback=white,boxrule=2.8pt,
  before upper={\popboxhead{Fiche bilan}{poporange}{L'essentiel à connaître pour l'examen}}}
% Figure encadrée avec légende
\newtcolorbox{popfigure}[1][]{enhanced,breakable=false,colback=white,colframe=popline,boxrule=1.4pt,arc=8pt,
  left=10pt,right=10pt,top=10pt,bottom=8pt,before skip=10pt,after skip=12pt,halign=center,
  lower separated=false,halign lower=center,
  fontlower=\popsemi\fontsize{9}{11.5}\selectfont\color{popmuted},
  #1}
\newcommand{\legende}[1]{\par\vspace{6pt}{\popsemi\fontsize{9}{11.5}\selectfont\color{popmuted}#1}}

% ============================================================
% Signature OnBuch+
% ============================================================
\newcommand{\poplogoinline}[1]{{\poplogo\fontsize{#1}{#1}\selectfont\color{popink}OnBuch\color{poporange}+}}
\newcommand{\signatureonbuch}{%
  \begin{tcolorbox}[enhanced,colback=popink,colframe=popink,boxrule=2.2pt,arc=10pt,
      drop shadow=poporange,left=14pt,right=14pt,top=10pt,bottom=10pt,before skip=0pt,after skip=0pt]
    \tikz[baseline=-0.7ex]\node[circle,fill=popgreen,draw=popcream,line width=1.3pt,minimum size=22pt,inner sep=0]{\popcheck{white}};%
    \hspace{9pt}\begin{minipage}[c]{0.62\linewidth}
      {\popxbold\fontsize{12}{13}\selectfont\color{popcream}Signé\ \textbullet\ {\poplogo OnBuch\color{poporange}+}}\par\vspace{2pt}
      {\raggedright\popsemi\fontsize{8}{10.5}\selectfont\color{popline}\MakeUppercase{Équipe pédagogique OnBuch+}\\\MakeUppercase{Conforme au programme officiel MINESEC}\par}
    \end{minipage}\hfill
    {\poplogo\fontsize{22}{22}\selectfont\color{popcream}OnBuch\color{poporange}+}
  \end{tcolorbox}%
}

% ============================================================
% Page de garde
% #1 titre · #2 matière · #3 niveau · #4 module · #5 n° leçon · #6 durée ·
% #7 description / objectifs (texte ou itemize)
% ============================================================
\newcommand{\pagegarde}[7]{%
  \thispagestyle{empty}
  \newgeometry{a4paper,margin=1.7cm,top=1.6cm,bottom=1.4cm}%
  \begin{tikzpicture}[remember picture,overlay]
    \fill[poporange!18] ([xshift=-3.2cm,yshift=-3.6cm]current page.north east) circle (4.6cm);
    \fill[poppurple!14] ([xshift=2.4cm,yshift=7.6cm]current page.south west) circle (3.8cm);
  \end{tikzpicture}%
  {\poplogo\fontsize{30}{30}\selectfont\color{popink}OnBuch\color{poporange}+}\hfill
  \raisebox{0.6ex}{\poppill{Cours signé \textbullet\ Premium}{popink}{popcream}}\par
  \vspace{1pt}\popsquiggle{poppurple}\par
  \vspace{8pt}
  \begin{tcolorbox}[enhanced,colback=poporange,colframe=popink,boxrule=2.8pt,arc=13pt,
      drop shadow=popink,left=18pt,right=18pt,top=12pt,bottom=13pt,after skip=10pt]
    \poppill{#2 \textbullet\ #3}{popink}{popcream}\hspace{5pt}\poppill{#4}{white}{popink}\par\vspace{8pt}
    {\popdisplay\fontsize{14}{14}\selectfont\color{popink}LEÇON #5}\par\vspace{4pt}
    {\raggedright\hyphenpenalty=10000\exhyphenpenalty=10000\popdisplay\fontsize{25}{27}\selectfont\color{popcream}\MakeUppercase{#1}\par}
    \vspace{8pt}
    {\popxbold\fontsize{9.5}{9.5}\selectfont\color{popink}\MakeUppercase{Durée conseillée : #6}}
  \end{tcolorbox}
  \begin{tcolorbox}[enhanced,colback=white,colframe=popink,boxrule=2.2pt,arc=10pt,
      drop shadow=popink,left=16pt,right=16pt,top=10pt,bottom=10pt,after skip=8pt]
    \poppill{Dans cette leçon}{poppurple}{white}\par\vspace{5pt}
    \popsemi\fontsize{9.3}{12.6}\selectfont\color{popink}#7
  \end{tcolorbox}
  \noindent\begin{minipage}[t]{0.487\linewidth}
    \begin{tcolorbox}[enhanced,colback=popgreen,colframe=popink,boxrule=2.2pt,arc=10pt,
        drop shadow=popink,left=11pt,right=11pt,top=10pt,bottom=10pt,height=3.1cm,valign=center]
      \tcbox[colback=white,colframe=popink,boxrule=1.3pt,arc=5pt,boxsep=0pt,nobeforeafter,
        left=3pt,right=3pt,top=3pt,bottom=3pt,valign=center]{\qrcode[height=1.9cm]{https://play.google.com/store/apps/details?id=com.luvvix.onbuch&hl=fr}}%
      \hspace{8pt}\begin{minipage}[c]{0.5\linewidth}
        {\popdisplay\fontsize{11}{12.5}\selectfont\color{white}\MakeUppercase{Télécharge OnBuch}}\par\vspace{4pt}
        {\raggedright\popsemi\fontsize{8.4}{11}\selectfont\color{white}Gratuit \textbullet\ Google Play\\Tuteur IA, annales, concours, résultats\par}
      \end{minipage}
    \end{tcolorbox}
  \end{minipage}\hfill
  \begin{minipage}[t]{0.487\linewidth}
    \begin{tcolorbox}[enhanced,colback=poppurple,colframe=popink,boxrule=2.2pt,arc=10pt,
        drop shadow=popink,left=11pt,right=11pt,top=10pt,bottom=10pt,height=3.1cm,valign=center]
      \tikz[baseline=-0.7ex]\node[circle,fill=white,draw=popink,line width=1.3pt,minimum size=1.6cm,inner sep=0]{\popdisplay\fontsize{24}{24}\selectfont\color{poppurple}+};%
      \hspace{8pt}\begin{minipage}[c]{0.6\linewidth}
        {\popdisplay\fontsize{11}{12.5}\selectfont\color{white}\MakeUppercase{Passe à OnBuch+}}\par\vspace{4pt}
        {\raggedright\popsemi\fontsize{8.4}{11}\selectfont\color{white}Tous les cours signés, Tuteur IA illimité, fiches PDF — onglet OnBuch+ de l'app\par}
      \end{minipage}
    \end{tcolorbox}
  \end{minipage}
  \vspace{8pt}
  \popcontenu
  \vfill
  \signatureonbuch
  \restoregeometry
  \newpage
}

% Rangée « Ce cours contient » (page de garde)
\newcommand{\poptile}[3]{% #1 couleur #2 gros texte #3 légende
  \begin{tcolorbox}[enhanced,colback=#1,colframe=popink,boxrule=2pt,arc=9pt,shadow={2.5pt}{-2.5pt}{0pt}{popink},
      left=6pt,right=6pt,top=9pt,bottom=9pt,halign=center,valign=center,height=1.75cm,width=\linewidth,nobeforeafter]
    {\popdisplay\fontsize{12.5}{13}\selectfont\color{white}\MakeUppercase{#2}}\par\vspace{3pt}
    {\centering\popsemi\fontsize{7.8}{9.5}\selectfont\color{white}#3\par}
  \end{tcolorbox}}
\newcommand{\popcontenu}{%
  \par\noindent{\popxboldls\fontsize{7.5}{7.5}\selectfont\color{popmuted}CE COURS SIGNÉ CONTIENT}\par\vspace{7pt}
  \noindent\begin{minipage}[t]{0.235\linewidth}\poptile{popblue}{Cours}{complet, illustré, pas à pas}\end{minipage}\hfill
  \begin{minipage}[t]{0.235\linewidth}\poptile{poporange}{Méthodes}{et exercices résolus}\end{minipage}\hfill
  \begin{minipage}[t]{0.235\linewidth}\poptile{poppink}{Exercices}{type examen + corrigés}\end{minipage}\hfill
  \begin{minipage}[t]{0.235\linewidth}\poptile{popgreen}{Fiche bilan}{l'essentiel à retenir}\end{minipage}\par}

% Sommaire
\newtcolorbox{sommaire}{enhanced,colback=white,colframe=popink,boxrule=2.2pt,arc=10pt,
  drop shadow=popink,left=18pt,right=18pt,top=13pt,bottom=13pt,before skip=6pt,after skip=20pt,
  before upper={\poppill{Sommaire}{poppurple}{white}\par\vspace{9pt}\popsemi\fontsize{10.6}{14.5}\selectfont\color{popink}}}

% Signature de fin de cours — poussée tout en bas de la dernière page
\newcommand{\findecours}{%
  \par\vfill\nopagebreak
  \begin{center}\popsquiggle{poporange}\end{center}\vspace{4pt}
  \signatureonbuch
}
\AtBeginDocument{\def\llbracket{\mathopen{[\mkern-3.5mu[}}\def\rrbracket{\mathclose{]\mkern-3.5mu]}}} % intervalles d'entiers (glyphe ⟦ absent de la police)
% Glyphes absents de Fira Math : redéfinis à partir de glyphes disponibles
\AtBeginDocument{%
  \def\setminus{\mathbin{\backslash}}\let\smallsetminus\setminus
  \def\vdots{\mathord{\vbox{\baselineskip3.2pt\lineskiplimit0pt\kern4pt\hbox{.}\hbox{.}\hbox{.}}}}%
  \def\ddots{\mathinner{\mkern1mu\raise6.5pt\hbox{.}\mkern2mu\raise3.6pt\hbox{.}\mkern2mu\raise0.7pt\hbox{.}\mkern1mu}}%
  \def\triangle{\mathord{\scalebox{0.9}{$\symup{\Delta}$}}}%
  \def\nabla{\mathord{\raisebox{\depth}{\scalebox{1}[-1]{$\symup{\Delta}$}}}}%
  \def\checkmark{\popcheck{popdark}}%
  \def\star{\mathbin{\ast}}\def\bigstar{\mathord{\ast}}%
  \def\diamond{\mathbin{\scalebox{0.6}{$\Diamond$}}}%
  \def\bigcup{\mathop{\mathchoice{\raisebox{-0.3em}{\scalebox{1.7}{$\cup$}}}{\scalebox{1.25}{$\cup$}}{\cup}{\cup}}\displaylimits}%
  \def\bigcap{\mathop{\mathchoice{\raisebox{-0.3em}{\scalebox{1.7}{$\cap$}}}{\scalebox{1.25}{$\cap$}}{\cap}{\cap}}\displaylimits}%
  \def\lesseqqgtr{\mathrel{\lessgtr}}\def\bigoplus{\mathop{\oplus}\displaylimits}%
}
\providecommand{\ssi}{si et seulement si} % abréviation fréquente
\AtBeginDocument{\providecommand{\wideparen}{\overparen}} % arc de cercle

\begin{document}

\pagegarde{Lectura y comentario : libertad de expresión y delitos de prensa ; análisis de publicidad y fake news}{Espagnol}{Tle A}{Módulo 5 : Médias, communication et culture de la paix}{E26}{2 h}{Cette leçon de lecture-commentaire propose un texte original sur la liberté d'expression et les délits de presse, suivi d'un travail méthodique sur l'analyse d'une publicité et la détection des fake news. Elle prépare directement à l'épreuve de commentaire du baccalauréat et à l'esprit critique face aux médias.
\vspace{4pt}
\begin{multicols}{2}\small
\begin{itemize}
\item À la fin de cette leçon, je sais lire et comprendre un texte argumentatif espagnol sur la liberté d'expression
\item je sais définir et employer le lexique des délits de presse (difamación, calumnia, censura)
\item je sais définir et employer le lexique de la désinformation (bulo, desinformación, manipulación)
\item je sais distinguer liberté d'expression, opinion et délit de presse
\item je sais analyser les procédés d'une publicité (eslogan, imagen, destinatario, mensaje)
\item je sais repérer une fausse information grâce à une méthode de vérification en cinq étapes
\end{itemize}\end{multicols}}

\begin{sommaire}
\begin{multicols}{2}
\begin{enumerate}
\item Texte support : « Libertad de expresión: ¿un derecho sin límites? »
\item Compréhension du texte
\item Lexique thématique par champs
\item Analyse d'une publicité
\item Détecter les fake news
\item Méthode du commentaire et commentaire modèle
\item Activité d'intégration
\item Méthodes et exercices résolus
\item Exercices
\item Corrigés des exercices
\item Fiche bilan
\end{enumerate}
\end{multicols}
\end{sommaire}

\begin{objectifs}
\begin{itemize}
\item À la fin de cette leçon, je sais lire et comprendre un texte argumentatif espagnol sur la liberté d'expression
\item je sais définir et employer le lexique des délits de presse (difamación, calumnia, censura)
\item je sais définir et employer le lexique de la désinformation (bulo, desinformación, manipulación)
\item je sais distinguer liberté d'expression, opinion et délit de presse
\item je sais analyser les procédés d'une publicité (eslogan, imagen, destinatario, mensaje)
\item je sais repérer une fausse information grâce à une méthode de vérification en cinq étapes
\item je sais construire un commentaire de texte structuré (introduction, développement, conclusion)
\item je sais rédiger un court article ou commentaire argumenté en espagnol sur les médias
\end{itemize}
\end{objectifs}

\begin{prerequis}
\begin{itemize}
\item Lexique des médias vu en Première (prensa, periódico, noticia, red social)
\item Connecteurs logiques de l'argumentation (sin embargo, por lo tanto, además, en cambio)
\item Expression de l'opinion (creo que, me parece que, en mi opinión) et subjonctif après verbes d'opinion négatifs
\item Méthode de la compréhension de texte (repérage du thème, de la thèse, des arguments)
\item Passé simple et imparfait pour rapporter des faits
\end{itemize}
\end{prerequis}

\newpage

\coursec{Texte support : « Libertad de expresión: ¿un derecho sin límites? »}

\courssub{Situation d'accroche et problématique}

Pendant les élections au Cameroun, les réseaux sociaux débordent de messages : certains informent, d'autres manipulent. Un lycéen de Yaoundé reçoit sur WhatsApp une photo soi-disant prise à Douala... qui date en réalité de dix ans et vient du Venezuela. En Espagne comme en Amérique latine, les médias parlent de \esp{bulos} et de \esp{desinformación}, et les tribunaux jugent des \esp{delitos de prensa} comme la \esp{difamación}. Où s'arrête la liberté d'expression et où commence le délit ? C'est la question que pose notre texte d'aujourd'hui.

\begin{aretenir}[Problématique de la leçon]
\esp{¿Es la libertad de expresión un derecho sin límites?} — La liberté d'expression est-elle un droit sans limites ? Tout le commentaire de texte consistera à suivre comment l'auteur répond à cette question : d'abord un constat, puis une thèse nuancée, enfin une proposition.
\end{aretenir}

\courssub{Présentation et lecture globale du texte}

Avant de lire, observons les \cle{indices paratextuels} : le titre, la source, la signature. Notre texte est un \cle{article d'opinion} (\esp{artículo de opinión}) publié dans la revue fictive \emph{El Debate} (Madrid). Le registre est \cle{journalistique} : phrases construites, connecteurs logiques, exemples d'actualité. La thèse est \cle{nuancée} : l'auteur défend la liberté d'expression tout en reconnaissant la nécessité de limites légales.

\begin{methode}[Première lecture : les quatre questions à se poser]
\begin{enumerate}
\item \textbf{Type de texte} : s'agit-il d'un récit, d'une description, d'un texte argumentatif ? Ici : \esp{artículo de opinión} (texte argumentatif).
\item \textbf{Source} : d'où vient le texte ? Ici : revue \emph{El Debate}, Madrid.
\item \textbf{Thème} : de quoi parle-t-on ? Ici : \esp{la libertad de expresión y los delitos de prensa}.
\item \textbf{Problématique} : quelle question le texte pose-t-il ? Ici : \esp{¿un derecho sin límites?}
\end{enumerate}
\end{methode}

\courssub{Le texte}

\begin{textebox}[« Libertad de expresión: ¿un derecho sin límites? » — El Debate, Madrid]
En los últimos años, las redes sociales han convertido a cada ciudadano en un periodista potencial. Basta con un teléfono móvil para publicar una noticia, una opinión o una denuncia que puede llegar a millones de personas en pocos minutos. Sin embargo, esta libertad tiene límites: la difamación y la calumnia son delitos de prensa castigados por la ley en la mayoría de los países democráticos.

El caso de la prensa española es ilustrativo. Varios periódicos han sido demandados por publicar informaciones falsas sobre la vida privada de ciertos políticos. Los tribunales han recordado que un titular mentiroso no es una opinión, sino un delito. No obstante, la censura previa sigue prohibida: ningún gobierno puede impedir la publicación de un artículo antes de su aparición.

En América latina, el problema es diferente. Allí, los bulos se difunden más rápido que las noticias verificadas, sobre todo durante las campañas electorales. Una fuente anónima puede inventar un escándalo y destruir la reputación de un candidato en veinticuatro horas. En varios países africanos, los periodistas denuncian leyes que, con el pretexto de luchar contra la desinformación, sirven en realidad para silenciar a la oposición.

¿Cuál es la solución? Por eso proponemos una verdadera educación en medios desde la escuela: aprender a verificar una fuente, a distinguir un hecho de una opinión y a reconocer la manipulación. En definitiva, la libertad de expresión no es un derecho sin límites, pero sus límites deben proteger la dignidad de las personas sin ahogar nunca el debate democrático.
\end{textebox}

\textbf{Glossaire (10 mots difficiles)}

\begin{center}
\begin{tabularx}{\linewidth}{|l|X|l|X|}
\hline
\rowcolor{popblueL} \textbf{Español} & \textbf{Français} & \textbf{Español} & \textbf{Français} \\
\hline
difamación & diffamation & fuente & source \\
\hline
calumnia & calomnie & demandar & poursuivre en justice \\
\hline
bulo & fausse rumeur, intox & difundir & diffuser, répandre \\
\hline
censura & censure & verificar & vérifier \\
\hline
titular & titre (d'article), gros titre & límites & limites \\
\hline
\end{tabularx}
\end{center}

\begin{attention}[Faux amis et pièges de traduction]
\begin{itemize}
\item \esp{prensa} signifie « la presse, les journaux » : ne pas confondre avec le portugais \emph{imprensa} ni traduire \esp{delitos de prensa} par « délits de presse » sans comprendre qu'il s'agit d'infractions commises \emph{par une publication}.
\item \esp{demandar} = « poursuivre en justice », et non « demander » (qui se dit \esp{pedir} ou \esp{preguntar}). \esp{El periódico fue demandado.} « Le journal a été poursuivi en justice. »
\item \esp{un titular} = « un titre de journal », et non « un titulaire ». \esp{Los titulares de hoy} « les gros titres d'aujourd'hui ».
\item \esp{una fuente} = « une source » (d'information), pas « une fontaine » dans ce contexte.
\item \esp{informaciones} : en espagnol, \esp{información} peut prendre le pluriel quand on parle de plusieurs « informations » distinctes ; en français, « information » est le plus souvent singulier.
\end{itemize}
\end{attention}

\begin{definition}[Delito de prensa]
Un \cle{delito de prensa} est une infraction commise par le biais d'une publication (presse écrite, radio, télévision, réseaux sociaux), comme la \esp{difamación} (atteinte à l'honneur par des faits), la \esp{calumnia} (accusation d'un délit que l'on sait fausse) ou l'\esp{injuria} (injure grave). Exemple : \esp{Publicar una acusación falsa en un periódico es un delito de prensa.} « Publier une accusation fausse dans un journal est un délit de presse. »
\end{definition}

\begin{savaistu}[L'article 20 de la Constitution espagnole]
En Espagne, l'article 20 de la Constitution de 1978 garantit la \esp{libertad de expresión} et interdit la \esp{censura previa} (censure préalable). Mais le code pénal punit la \esp{calumnia} et l'\esp{injuria}. La formule juridique est célèbre : ces droits sont limités par « \esp{el respeto a los derechos reconocidos en este título} », c'est-à-dire le respect de l'honneur, de la vie privée et de l'image des personnes. Au Cameroun, la législation sur la communication sociale poursuit un objectif comparable : protéger la réputation sans étouffer le débat public.
\end{savaistu}

\courssub{Structure du texte}

Le texte suit le plan classique de l'article d'opinion : \cle{introduction} (constat), \cle{développement} (arguments illustrés d'exemples), \cle{conclusion} (proposition). Observons-le paragraphe par paragraphe :

\begin{itemize}
\item \textbf{Introduction (§1)} : constat sur les réseaux sociaux — « \esp{cada ciudadano es un periodista potencial} » — et annonce de la thèse : la liberté a des limites.
\item \textbf{Argument 1 (§2)} : le droit contre les abus — exemple de la \cle{presse espagnole} poursuivie pour informations fausses ; mais rappel : la censure préalable reste interdite.
\item \textbf{Argument 2 (§3)} : les bulos — exemple des \cle{réseaux sociaux en Amérique latine} pendant les campagnes électorales.
\item \textbf{Argument 3 (§3, fin)} : le danger inverse — en \cle{Afrique}, des lois contre la désinformation peuvent servir à faire taire l'opposition.
\item \textbf{Conclusion (§4)} : proposition — l'\cle{éducation aux médias} à l'école ; formule finale nuancée : des limites qui protègent la dignité « \esp{sin ahogar nunca el debate democrático} ».
\end{itemize}

\begin{popfigure}
\begin{center}
\begin{tikzpicture}[
  node distance=0.55cm,
  box/.style={draw=popblue, fill=popblueL, rounded corners, text width=6.2cm, align=center, font=\small, thick},
  arr/.style={-{Stealth[length=2.5mm]}, thick, popdark}
]
\node[box, align=center] (intro){\textbf{Introducción (§1)}\\ Constat : las redes sociales, cada ciudadano periodista potencial};
\node[box, below=of intro, align=center] (tesis){\textbf{Tesis}\\ La libertad de expresión tiene límites};
\node[box, below=of tesis, align=center] (arg1){\textbf{Argumento 1 (§2)}\\ Difamación y calumnia : delitos (prensa española)};
\node[box, below=of arg1, align=center] (arg2){\textbf{Argumento 2 (§3)}\\ Los bulos se difunden rápido (América latina)};
\node[box, below=of arg2, align=center] (arg3){\textbf{Argumento 3 (§3)}\\ Leyes abusivas contra la oposición (África)};
\node[box, below=of arg3, fill=popgreenL, draw=popgreen, align=center] (concl){\textbf{Conclusión (§4)}\\ Educación en medios : verificar, distinguir, reconocer};
\draw[arr] (intro) -- (tesis);
\draw[arr] (tesis) -- (arg1);
\draw[arr] (arg1) -- (arg2);
\draw[arr] (arg2) -- (arg3);
\draw[arr] (arg3) -- (concl);
\end{tikzpicture}
\end{center}
\legende{Organigramme de la structure du texte : introduction, thèse, trois arguments, conclusion.}
\end{popfigure}

\courssub{Les marqueurs argumentatifs du texte}

Pour suivre la démonstration, repérons les \cle{conectores} (connecteurs logiques). Ce sont eux qui signalent les étapes du raisonnement.

\begin{center}
\begin{tabularx}{\linewidth}{|l|X|X|}
\hline
\rowcolor{popblueL} \textbf{Conector} & \textbf{Valeur logique} & \textbf{Exemple du texte} \\
\hline
\esp{sin embargo} & opposition, restriction & \esp{Sin embargo, esta libertad tiene límites.} \\
\hline
\esp{no obstante} & concession, opposition forte & \esp{No obstante, la censura previa sigue prohibida.} \\
\hline
\esp{por eso} & conséquence, conclusion partielle & \esp{Por eso proponemos una verdadera educación en medios.} \\
\hline
\esp{en definitiva} & conclusion générale, résumé & \esp{En definitiva, la libertad de expresión no es un derecho sin límites.} \\
\hline
\esp{sobre todo} & mise en relief, insistance & \esp{Sobre todo durante las campañas electorales.} \\
\hline
\end{tabularx}
\end{center}

\begin{attention}[Erreurs fréquentes des francophones]
\begin{itemize}
\item \esp{sin embargo} se place souvent après une virgule, comme « cependant » : \esp{La ley es clara; sin embargo, no se aplica.} Ne pas écrire \esp{sin de embargo}.
\item \esp{por eso} = « c'est pourquoi » ; ne pas confondre avec \esp{porque} (parce que). \esp{No verificó la fuente; por eso difundió un bulo.} « Il n'a pas vérifié la source ; c'est pourquoi il a diffusé une intox. »
\item \esp{en definitiva} = « en fin de compte » ; le français « définitivement » se dit plutôt \esp{definitivamente}.
\item \esp{basta con} + nom ou infinitif = « il suffit de » : \esp{Basta con un teléfono móvil.} « Il suffit d'un téléphone portable. » Ne pas oublier la préposition \esp{con}.
\end{itemize}
\end{attention}

\courssub{Carte mentale lexicale autour du thème}

\begin{popfigure}
\begin{center}
\begin{tikzpicture}[mindmap, grow cyclic, every node/.style={concept, align=center, font=\small},
  root concept/.append style={concept color=popblue, text=white, font=\small\bfseries},
  level 1/.append style={sibling angle=90, level distance=3.1cm},
  level 2/.append style={sibling angle=60, level distance=2.3cm, font=\scriptsize}]
\node[root concept]{libertad de expresión}
  child[concept color=popgreen] { node {derechos}
    child { node {opinar} }
    child { node {informar} }
    child { node {denunciar} } }
  child[concept color=poporange] { node {límites}
    child { node {honor} }
    child { node {vida privada} }
    child { node {dignidad} } }
  child[concept color=poppink] { node {delitos}
    child { node {difamación} }
    child { node {calumnia} }
    child { node {injuria} } }
  child[concept color=popgold] { node {medios}
    child { node {prensa} }
    child { node {redes sociales} }
    child { node {bulos} } };
\end{tikzpicture}
\end{center}
\legende{Carte mentale : le champ lexical de la liberté d'expression (droits, limites, délits, médias).}
\end{popfigure}

\courssub{Exercice résolu de repérage}

\begin{exoresolu}[Repérage dans le texte]
\textbf{Énoncé.} 1) Quel connecteur introduit la thèse de l'auteur au premier paragraphe ? 2) Citez les trois exemples géographiques utilisés pour illustrer le développement. 3) Quelle solution l'auteur propose-t-il dans la conclusion ? Traduisez la phrase qui la contient. \tcblower
\textbf{Solution.} 1) Le connecteur est \esp{sin embargo} : \esp{Sin embargo, esta libertad tiene límites} — il introduit la restriction qui constitue la thèse. 2) Les trois exemples sont : la presse espagnole (journaux poursuivis pour informations fausses), les réseaux sociaux en Amérique latine (bulos pendant les campagnes électorales) et plusieurs pays africains (lois utilisées pour réduire l'opposition au silence). 3) La solution est l'éducation aux médias : \esp{Por eso proponemos una verdadera educación en medios desde la escuela: aprender a verificar una fuente, a distinguir un hecho de una opinión y a reconocer la manipulación.} Traduction : « C'est pourquoi nous proposons une véritable éducation aux médias dès l'école : apprendre à vérifier une source, à distinguer un fait d'une opinion et à reconnaître la manipulation. »
\end{exoresolu}

\begin{exercice}[À vous de jouer]
\begin{enumerate}
\item Relevez dans le texte deux verbes du champ lexical de la vérification de l'information et traduisez-les.
\item Trouvez dans le texte l'expression qui signifie « faire taire l'opposition ».
\item Complétez avec \esp{sin embargo}, \esp{por eso} ou \esp{en definitiva} : \esp{El periodista verificó sus fuentes; \ldots, publicó la noticia sin miedo.}
\end{enumerate}
\end{exercice}

\begin{corrige}[Exercice « À vous de jouer »]
\begin{enumerate}
\item \esp{verificar} (vérifier) et \esp{distinguir} (distinguer) ; on accepte aussi \esp{reconocer} (reconnaître).
\item \esp{silenciar a la oposición}.
\item \esp{El periodista verificó sus fuentes; por eso publicó la noticia sin miedo.} « Le journaliste a vérifié ses sources ; c'est pourquoi il a publié la nouvelle sans crainte. » (conséquence logique : \esp{por eso}).
\end{enumerate}
\end{corrige}

\begin{aretenir}[Ce qu'il faut retenir de cette section]
\begin{itemize}
\item Le texte est un \cle{article d'opinion} à thèse nuancée : la liberté d'expression est un droit fondamental, mais elle a des limites légales (\esp{difamación}, \esp{calumnia}, \esp{injuria}).
\item Sa structure est un modèle pour vos propres essais : constat → thèse → arguments illustrés d'exemples (Espagne, Amérique latine, Afrique) → proposition (éducation aux médias).
\item Les connecteurs \esp{sin embargo}, \esp{no obstante}, \esp{por eso}, \esp{en definitiva} balisent le raisonnement : apprenez à les repérer et à les réutiliser.
\end{itemize}
\end{aretenir}

\coursec{Lexique thématique par champs}

\courssub{Le champ des droits et des limites}

\begin{center}
\begin{tabularx}{\linewidth}{|X|X|X|}
\hline
\rowcolor{popblueL} \textbf{Español} & \textbf{Français} & \textbf{Exemple} \\
\hline
\esp{la libertad de expresión} & la liberté d'expression & \esp{La libertad de expresión es un derecho fundamental.} \\
\hline
\esp{el derecho a la información} & le droit à l'information & \esp{Los ciudadanos tienen derecho a la información.} \\
\hline
\esp{el honor} & l'honneur & \esp{La ley protege el honor de las personas.} \\
\hline
\esp{la vida privada} & la vie privée & \esp{No se puede publicar nada sobre su vida privada.} \\
\hline
\esp{la dignidad} & la dignité & \esp{Los límites protegen la dignidad.} \\
\hline
\esp{la censura (previa)} & la censure (préalable) & \esp{La censura previa está prohibida.} \\
\hline
\end{tabularx}
\end{center}

\courssub{Le champ des délits et de la justice}

\begin{center}
\begin{tabularx}{\linewidth}{|X|X|X|}
\hline
\rowcolor{popblueL} \textbf{Español} & \textbf{Français} & \textbf{Exemple} \\
\hline
\esp{la difamación} & la diffamation & \esp{La difamación es un delito de prensa.} \\
\hline
\esp{la calumnia} & la calomnie & \esp{Fue condenado por calumnia.} \\
\hline
\esp{la injuria} & l'injure (grave) & \esp{Las injurias están castigadas por la ley.} \\
\hline
\esp{el delito de prensa} & le délit de presse & \esp{Publicar mentiras es un delito de prensa.} \\
\hline
\esp{demandar (a alguien)} & poursuivre (qqn) en justice & \esp{El político demandó al periódico.} \\
\hline
\esp{el tribunal} & le tribunal & \esp{El tribunal recordó la ley.} \\
\hline
\esp{condenar} & condamner & \esp{Condenaron al periodista a pagar una multa.} \\
\hline
\end{tabularx}
\end{center}

\courssub{Le champ de l'information et de la désinformation}

\begin{center}
\begin{tabularx}{\linewidth}{|X|X|X|}
\hline
\rowcolor{popblueL} \textbf{Español} & \textbf{Français} & \textbf{Exemple} \\
\hline
\esp{el bulo} & la fausse rumeur, l'intox & \esp{El bulo se difundió por WhatsApp.} \\
\hline
\esp{la desinformación} & la désinformation & \esp{Hay que luchar contra la desinformación.} \\
\hline
\esp{la manipulación} & la manipulation & \esp{Aprendimos a reconocer la manipulación.} \\
\hline
\esp{la fuente (anónima)} & la source (anonyme) & \esp{Una fuente anónima inventó el escándalo.} \\
\hline
\esp{verificar / comprobar} & vérifier & \esp{Verifica la fuente antes de compartir.} \\
\hline
\esp{el titular} & le titre, le gros titre & \esp{El titular era mentiroso.} \\
\hline
\esp{las noticias verificadas} & les informations vérifiées & \esp{Los bulos van más rápido que las noticias verificadas.} \\
\hline
\esp{compartir / difundir} & partager / diffuser & \esp{No difundas noticias falsas.} \\
\hline
\end{tabularx}
\end{center}

\begin{attention}[Genre et orthographe à surveiller]
\begin{itemize}
\item \esp{la difamación}, \esp{la calumnia}, \esp{la censura}, \esp{la fuente} : tous féminins. Mais \esp{el bulo}, \esp{el titular}, \esp{el delito}, \esp{el tribunal} : masculins.
\item Accentuation : \esp{difamación}, \esp{información}, \esp{opinión}, \esp{manipulación} prennent un accent sur la dernière syllabe (mots terminés par \esp{-ción}). Au pluriel, l'accent disparaît : \esp{informaciones}, \esp{opiniones}.
\item \esp{noticia} = « une nouvelle, une information » ; « les nouvelles » (le journal télévisé) se dit \esp{las noticias}.
\end{itemize}
\end{attention}

\coursec{Analyse d'une publicité}

\courssub{Observer une publicité : le texte support}

\begin{textebox}[Publicité fictive : « Refresco Tropical Sol »]
¡Descubre el sabor de África! Refresco Tropical Sol, la bebida número uno de los jóvenes cameruneses. Con fruta natural de la región del Litoral, te da energía para todo el día. Nueve de cada diez jóvenes lo prefieren. ¡Pruébalo ya y sé libre! Tropical Sol: el sabor de tu éxito.
\end{textebox}

\textbf{Traduction.} « Découvrez le goût de l'Afrique ! Le soda Tropical Sol, la boisson numéro un des jeunes Camerounais. Avec des fruits naturels de la région du Littoral, il te donne de l'énergie pour toute la journée. Neuf jeunes sur dix le préfèrent. Goûte-le tout de suite et sois libre ! Tropical Sol : le goût de ton succès. »

\courssub{La méthode d'analyse}

\begin{methode}[Analyser une publicité en cinq étapes]
\begin{enumerate}
\item \textbf{Le produit et la cible} : que vend-on et à qui ? Ici : un soda, destiné aux \esp{jóvenes} (jeunes).
\item \textbf{Le slogan} (\esp{eslogan}) : phrase courte et mémorable. Ici : \esp{El sabor de tu éxito} — il associe la boisson à la réussite.
\item \textbf{Les arguments} : rationnels (\esp{fruta natural}) ou émotionnels (\esp{libertad}, \esp{éxito}, \esp{energía}).
\item \textbf{Les procédés de manipulation} : chiffres invérifiables (\esp{nueve de cada diez}), superlatifs (\esp{número uno}), impératifs (\esp{¡Pruébalo ya!}), tutoiement complice (\esp{tu éxito}).
\item \textbf{Le jugement critique} : la publicité informe-t-elle ou manipule-t-elle ? Ici, aucune preuve n'est donnée : c'est de la \esp{persuasión}, pas de l'information.
\end{enumerate}
\end{methode}

\begin{propriete}[L'impératif publicitaire]
La publicité utilise massivement l'\cle{impératif} pour provoquer l'achat. Formes de tutoiement : \esp{¡Descubre!} (découvre), \esp{¡Pruébalo!} (goûte-le), \esp{¡Sé libre!} (sois libre — impératif irrégulier de \esp{ser}). Attention au pronom accolé : \esp{pruébalo} = \esp{prueba} + \esp{lo}, avec accent d'intensité conservé. À l'écrit, n'oubliez pas les signes d'exclamation encadrants : \esp{¡...!}
\end{propriete}

\begin{exemplebox}[Procédés publicitaires : exemples traduits]
\begin{itemize}
\item \esp{La bebida número uno.} « La boisson numéro un. » → superlatif invérifiable.
\item \esp{Nueve de cada diez jóvenes lo prefieren.} « Neuf jeunes sur dix le préfèrent. » → faux chiffre : quelle étude le prouve ?
\item \esp{Te da energía para todo el día.} « Il te donne de l'énergie pour toute la journée. » → promesse exagérée.
\item \esp{¡Pruébalo ya y sé libre!} « Goûte-le tout de suite et sois libre ! » → impératifs + valeur affective (la liberté) sans rapport avec le produit.
\end{itemize}
\end{exemplebox}

\begin{exoresolu}[Analyse guidée]
\textbf{Énoncé.} Relevez dans la publicité « Tropical Sol » deux procédés de manipulation et expliquez-les en une phrase chacun. \tcblower
\textbf{Solution.} 1) Le chiffre \esp{nueve de cada diez jóvenes} : statistique invérifiable qui crée un effet de masse (« tout le monde le boit »). 2) L'impératif \esp{¡Pruébalo ya!} avec l'adverbe \esp{ya} (« tout de suite ») : il crée l'urgence et court-circuite la réflexion du consommateur. On peut ajouter le slogan \esp{El sabor de tu éxito}, qui associe par simple juxtaposition la boisson et la réussite sociale.
\end{exoresolu}

\coursec{Détecter les fake news}

\courssub{Qu'est-ce qu'un « bulo » ?}

\begin{definition}[Bulo et desinformación]
Un \cle{bulo} est une fausse information fabriquée pour tromper, qui se répand vite sur les réseaux sociaux. La \cle{desinformación} est l'ensemble des techniques de manipulation de l'information : bulos, photos sorties de leur contexte, titres trompeurs (\esp{titulares engañosos}), fausses sources. En français courant : « fake news », « intox ».
\end{definition}

\courssub{La grille de vérification}

\begin{methode}[Cinq questions pour détecter une fausse information]
\begin{enumerate}
\item \textbf{La source} (\esp{la fuente}) : qui publie ? Un site connu ou un compte anonyme ? \esp{¿Quién publica la noticia?}
\item \textbf{La date} (\esp{la fecha}) : l'information est-elle récente ou ressortie d'un vieux contexte ? \esp{¿De cuándo es la foto?}
\item \textbf{Le titre} (\esp{el titular}) : est-il exagéré, choquant, écrit en majuscules ? Les bulos jouent sur l'émotion.
\item \textbf{Les preuves} (\esp{las pruebas}) : y a-t-il des faits vérifiables, des témoins nommés, d'autres médias qui confirment ?
\item \textbf{Le recoupement} (\esp{la verificación}) : chercher la même information dans deux ou trois médias sérieux avant de la partager. \esp{Verifica antes de compartir.}
\end{enumerate}
\end{methode}

\begin{exemplebox}[Étude de cas : le message WhatsApp du lycéen]
Message reçu : \esp{¡URGENTE! ¡COMPARTE! Manifestación violenta ahora mismo en Douala. La policía dispara contra los estudiantes. Foto impactante.} « URGENT ! PARTAGE ! Manifestation violente en ce moment même à Douala. La police tire sur les étudiants. Photo choquante. »
Signaux d'alerte : majuscules et impératifs (\esp{¡COMPARTE!}), aucune source nommée, vocabulaire émotionnel (\esp{impactante}, \esp{urgente}), photo non datée. Vérification : la photo, prise en réalité il y a dix ans dans un autre pays, ne montre pas Douala. Conclusion : c'est un \esp{bulo}.
\end{exemplebox}

\begin{attention}[Vocabulaire : ne pas calquer le français]
\begin{itemize}
\item « Une fausse nouvelle » se dit \esp{una noticia falsa} ou \esp{un bulo}, jamais \esp{una falsa noticia} dans le sens journalistique (l'adjectif épithète change de sens selon sa place : \esp{falsa} avant le nom est littéraire).
\item « Partager » (sur les réseaux) = \esp{compartir} ; « diffuser » = \esp{difundir}. \esp{No compartas bulos.} « Ne partage pas d'intox. »
\item « Vérifier » = \esp{verificar} ou \esp{comprobar} (irrégulier : \esp{compruebo}, \esp{comprueba}...). \esp{Comprueba la fuente.} « Vérifie la source. »
\end{itemize}
\end{attention}

\begin{experience}[Atelier : chasseurs de bulos]
Par groupes de quatre, les élèves reçoivent trois messages écrits au tableau par le professeur (deux vrais, un faux). Chaque groupe applique la grille des cinq questions, puis présente son verdict en espagnol en deux minutes : \esp{En nuestra opinión, es un bulo porque no tiene fuente y el titular es exagerado.} « À notre avis, c'est une intox parce qu'il n'a pas de source et que le titre est exagéré. » La classe vote ensuite pour le groupe le plus convaincant.
\end{experience}

\coursec{Méthode du commentaire de texte et commentaire modèle}

\courssub{La méthode pas à pas}

\begin{methode}[Commenter un texte argumentatif au baccalauréat]
\begin{enumerate}
\item \textbf{Lire deux fois} le texte : une fois pour le sens global, une fois en soulignant connecteurs, thèse et exemples.
\item \textbf{Présenter le texte} (introduction du commentaire) : type, source, thème, problématique. \esp{Este texto es un artículo de opinión publicado en la revista El Debate. Trata de la libertad de expresión y sus límites.}
\item \textbf{Dégager la structure} : annoncer le plan en deux ou trois parties (\esp{primero... luego... finalmente...}).
\item \textbf{Analyser} chaque partie : relever les arguments, les exemples, les procédés (connecteurs, questions rhétoriques, chiffres) et citer le texte entre guillemets.
\item \textbf{Conclure} : bilan de la thèse de l'auteur + ouverture personnelle argumentée (\esp{En mi opinión...}, \esp{A mi juicio...}).
\end{enumerate}
\end{methode}

\begin{propriete}[Les expressions d'opinion et le subjonctif]
Après \esp{creo que}, \esp{pienso que}, \esp{me parece que} affirmatifs : \textbf{indicatif}. \esp{Creo que la libertad tiene límites.} « Je pense que la liberté a des limites. »
Après les formes \textbf{négatives} : \textbf{subjonctif}. \esp{No creo que la libertad sea un derecho sin límites.} « Je ne pense pas que la liberté soit un droit sans limites. »
Rappel du subjonctif présent de \esp{ser} : \esp{sea, seas, sea, seamos, seáis, sean}.
\end{propriete}

\courssub{Commentaire modèle (rédigé en espagnol)}

\begin{textebox}[Comentario modelo]
Este texto es un artículo de opinión publicado en la revista El Debate, de Madrid. Trata de la libertad de expresión y de los delitos de prensa, y plantea una pregunta central: ¿es este derecho un derecho sin límites?

En primer lugar, el autor constata que las redes sociales han convertido a cada ciudadano en « un periodista potencial ». Sin embargo, anuncia enseguida su tesis: esta libertad tiene límites, porque la difamación y la calumnia son delitos castigados por la ley.

Luego, el desarrollo presenta tres argumentos ilustrados con ejemplos geográficos. Primero, el caso de la prensa española muestra que los tribunales castigan las informaciones falsas, aunque la censura previa sigue prohibida. Después, el ejemplo de América latina demuestra que los bulos se difunden más rápido que las noticias verificadas durante las campañas electorales. Finalmente, el autor denuncia las leyes africanas que, con el pretexto de luchar contra la desinformación, sirven para silenciar a la oposición.

En conclusión, el autor propone una solución educativa: « una verdadera educación en medios desde la escuela ». En mi opinión, esta propuesta es muy acertada, porque un ciudadano que sabe verificar una fuente no necesita censura. No creo que la libertad de expresión sea un derecho sin límites, pero pienso que sus límites deben proteger la dignidad sin ahogar el debate democrático.
\end{textebox}

\textbf{Traduction du commentaire modèle.} « Ce texte est un article d'opinion publié dans la revue El Debate, de Madrid. Il traite de la liberté d'expression et des délits de presse, et pose une question centrale : ce droit est-il un droit sans limites ? / Premièrement, l'auteur constate que les réseaux sociaux ont fait de chaque citoyen "un journaliste potentiel". Cependant, il annonce aussitôt sa thèse : cette liberté a des limites, car la diffamation et la calomnie sont des délits punis par la loi. / Ensuite, le développement présente trois arguments illustrés d'exemples géographiques. D'abord, le cas de la presse espagnole montre que les tribunaux punissent les informations fausses, bien que la censure préalable reste interdite. Puis, l'exemple de l'Amérique latine démontre que les intox se diffusent plus vite que les informations vérifiées pendant les campagnes électorales. Enfin, l'auteur dénonce les lois africaines qui, sous prétexte de lutter contre la désinformation, servent à faire taire l'opposition. / En conclusion, l'auteur propose une solution éducative : "une véritable éducation aux médias dès l'école". À mon avis, cette proposition est très judicieuse, car un citoyen qui sait vérifier une source n'a pas besoin de censure. Je ne pense pas que la liberté d'expression soit un droit sans limites, mais je pense que ses limites doivent protéger la dignité sans étouffer le débat démocratique. »

\begin{attention}[Les erreurs à éviter dans votre commentaire]
\begin{itemize}
\item Ne pas résumer sans analyser : chaque citation doit être \textbf{commentée} (procédé + effet).
\item \esp{tratar de} = « traiter de » : \esp{El texto trata de la libertad de expresión.} Ne pas confondre avec \esp{tratar de} + infinitif (« essayer de »).
\item \esp{acertado} = « judicieux, juste » (d'une idée) ; ne pas dire \esp{correcto} pour une opinion.
\item Attention à l'accord : \esp{una pregunta central}, \esp{leyes africanas}, \esp{noticias verificadas}.
\end{itemize}
\end{attention}

\begin{exercice}[Entraînement au commentaire]
Rédigez en espagnol un commentaire de 15 à 20 lignes du texte « Libertad de expresión: ¿un derecho sin límites? » en suivant le plan du commentaire modèle : présentation, structure, analyse d'un procédé au choix (connecteurs, exemples géographiques ou question rhétorique), conclusion personnelle avec \esp{no creo que} + subjonctif.
\end{exercice}

\begin{corrige}[Exercice « Entraînement au commentaire » — éléments attendus]
\begin{itemize}
\item \textbf{Présentation} : type (\esp{artículo de opinión}), source (\esp{El Debate}, Madrid), thème, problématique sous forme de question.
\item \textbf{Structure} : introduction (constat), développement (trois arguments : Espagne, Amérique latine, Afrique), conclusion (proposition éducative).
\item \textbf{Analyse du procédé} : par exemple, la question rhétorique \esp{¿Cuál es la solución?} relance le raisonnement et prépare la conclusion ; ou les connecteurs \esp{sin embargo} / \esp{en definitiva} qui balisent la thèse nuancée.
\item \textbf{Conclusion personnelle} : avis argumenté avec subjonctif correct, par exemple \esp{No creo que la censura sea la solución; pienso que la educación en medios es más eficaz.} « Je ne pense pas que la censure soit la solution ; je pense que l'éducation aux médias est plus efficace. »
\end{itemize}
\end{corrige}

\begin{aretenir}[Bilan de la leçon]
\begin{itemize}
\item \textbf{Lexique} : \esp{difamación}, \esp{calumnia}, \esp{injuria}, \esp{censura}, \esp{bulo}, \esp{desinformación}, \esp{publicidad}, \esp{eslogan}, \esp{manipulación}, \esp{fuente}, \esp{titular}, \esp{verificar}.
\item \textbf{Méthode} : un commentaire = présentation + structure + analyse de procédés cités + conclusion personnelle.
\item \textbf{Grammaire mobilisée} : connecteurs argumentatifs, impératif publicitaire, subjonctif après \esp{no creo que}.
\item \textbf{Esprit critique} : avant de partager une information, toujours vérifier la source, la date, le titre et les preuves : \esp{Verifica antes de compartir.}
\end{itemize}
\end{aretenir}

\coursec{Compréhension du texte}

Dans la section précédente, nous avons lu et observé le texte support \esp{« Libertad de expresión: ¿un derecho sin límites? »}. Nous avons identifié son thème général et dégagé le champ lexical des médias. Il s'agit maintenant de passer à l'étape décisive de toute épreuve de lecture au baccalauréat : la \cle{compréhension du texte}, c'est-à-dire la capacité à répondre, en espagnol et en phrases complètes, à des questions de compréhension globale, fine et d'interprétation, en s'appuyant toujours sur des citations précises. Pour travailler rigoureusement, relisons d'abord le texte avec ses lignes numérotées.

\courssub{Le texte support relu, lignes numérotées}

\begin{textebox}[Libertad de expresión: ¿un derecho sin límites? (artículo de opinión)]
\textbf{(1)} En toda sociedad democrática, la libertad de expresión es un derecho fundamental: \textbf{(2)} los ciudadanos pueden criticar al gobierno, informar y opinar sin miedo. Sin \textbf{(3)} embargo, este derecho no es absoluto. Cuando una mentira destruye la \textbf{(4)} reputación de una persona, ya no hablamos de libertad, sino de difamación. \textbf{(5)} La calumnia es acusar a alguien de un delito sabiendo que es falso, y la \textbf{(6)} ley la castiga en España, en América Latina y en Guinea Ecuatorial.

\textbf{(7)} Hoy, el problema se agrava con las redes sociales. Los ciudadanos asisten \textbf{(8)} a la difusión de bulos a una velocidad impresionante: una noticia falsa \textbf{(9)} puede recorrer el mundo antes de que la verdad se ponga los zapatos. La \textbf{(10)} desinformación no es un simple error: es una estrategia de manipulación que \textbf{(11)} busca engañar al público, influir en las elecciones o dañar a un adversario.

\textbf{(12)} Frente a este fenómeno, algunos reclaman más censura. Pero la censura \textbf{(13)} es un arma peligrosa: ayer silenció a los periodistas honestos, hoy podría \textbf{(14)} silenciar a cualquier ciudadano. La solución no es prohibir, sino educar. \textbf{(15)} Enseñar a verificar las fuentes, a comparar las informaciones y a dudar \textbf{(16)} de los titulares sensacionalistas es la mejor defensa de la democracia.

\textbf{(17)} La libertad de expresión es, en definitiva, un derecho acompañado de una \textbf{(18)} responsabilidad. Quien habla debe respetar la verdad y la dignidad de los \textbf{(19)} demás. Solo así la palabra seguirá siendo un instrumento de paz y no un \textbf{(20)} arma de destrucción.
\end{textebox}

\textbf{Glossaire (mots difficiles) :}

\begin{center}
\begin{tabularx}{\linewidth}{|X|X|}
\hline
\rowcolor{popblueL} \textbf{Español} & \textbf{Français} \\
\hline
el derecho fundamental & le droit fondamental \\
\hline
sin embargo & cependant, pourtant \\
\hline
la reputación & la réputation \\
\hline
castigar & punir \\
\hline
agravarse & s'aggraver \\
\hline
el bulo & la rumeur (fausse), l'intox \\
\hline
engañar & tromper, duper \\
\hline
reclamar & réclamer, exiger \\
\hline
silenciar & faire taire, réduire au silence \\
\hline
la fuente & la source \\
\hline
el titular sensacionalista & le titre sensationnaliste (putaclic) \\
\hline
la dignidad & la dignité \\
\hline
\end{tabularx}
\end{center}

\begin{attention}[Faux ami : \esp{asistir}]
Le verbe \esp{asistir} est un piège classique pour les francophones. \esp{Asistir a} signifie « assister à » au sens d'\emph{être présent, être témoin de}, et non « aider » (qui se dit \esp{ayudar}). Ainsi, à la ligne 7 : \esp{« los ciudadanos asisten a la difusión de bulos »} = « les citoyens \emph{assistent à / sont témoins de} la diffusion de rumeurs ». Ne traduisez jamais : « les citoyens aident la diffusion des rumeurs ».\par
Autre faux ami du texte : \esp{reclamar} (ligne 12) signifie « réclamer, exiger », et non « réclamer » au sens de « protester contre » ; pour « se plaindre », l'espagnol utilise \esp{quejarse}.
\end{attention}

\begin{propriete}[Deux tournures à observer dans le texte]
\begin{itemize}
\item \esp{« la ley la castiga »} (ligne 6) : le pronom \esp{la} reprend \esp{la calumnia} (complément d'objet direct féminin singulier placé avant le verbe). Traduction : « et la loi \emph{la} punit ».
\item \esp{« antes de que la verdad se ponga los zapatos »} (ligne 9) : après \esp{antes de que}, l'espagnol exige le \cle{subjonctif} (\esp{se ponga}, subjonctif présent de \esp{ponerse}), même quand le français utilise l'indicatif : « avant que la vérité \emph{mette} ses chaussures ». C'est une image proverbiale : la rumeur court plus vite que la vérité.
\end{itemize}
\end{propriete}

\courssub{Questions de compréhension globale}

Les questions globales portent sur le texte dans son ensemble : thème, genre, thèse, ton. La stratégie consiste à relire le \cle{titre}, le \cle{premier paragraphe} et la \cle{conclusion}.

\begin{exemplebox}[Questions globales et réponses modèles]
\textbf{1. ¿Cuál es el tema del texto?}\par
\esp{El tema del texto es la libertad de expresión y sus límites en la era de las redes sociales.}\par
« Le thème du texte est la liberté d'expression et ses limites à l'ère des réseaux sociaux. »

\textbf{2. ¿Qué tipo de texto es?}\par
\esp{Es un artículo de opinión, porque el autor defiende una tesis personal y utiliza argumentos para convencer al lector.}\par
« C'est un article d'opinion, car l'auteur défend une thèse personnelle et utilise des arguments pour convaincre le lecteur. »

\textbf{3. ¿Cuál es la tesis del autor?}\par
\esp{El autor defiende que la libertad de expresión es un derecho fundamental, pero no absoluto: debe ir acompañada de responsabilidad, como dice en la línea 17: « un derecho acompañado de una responsabilidad ».}\par
« L'auteur soutient que la liberté d'expression est un droit fondamental mais non absolu : elle doit s'accompagner de responsabilité. »

\textbf{4. ¿Cuál es el tono del texto?}\par
\esp{El tono es crítico pero matizado: el autor critica la desinformación y la manipulación (líneas 10-11), pero también rechaza la censura, a la que llama « un arma peligrosa » (línea 13).}\par
« Le ton est critique mais nuancé : l'auteur critique la désinformation, mais refuse aussi la censure. »
\end{exemplebox}

Remarquez la méthode : chaque réponse est une \cle{phrase complète} en espagnol, qui reprend les mots de la question et s'appuie, quand c'est possible, sur une \cle{citation entre guillemets} avec le numéro de ligne. Notez aussi la tournure de la réponse 4 : \esp{« a la que llama »} = « qu'il appelle », avec la préposition \esp{a} reprise devant le relatif \esp{la que} (on dit \esp{llamar algo a alguien}).

\courssub{Questions de compréhension fine : le vrai/faux justifié}

Au baccalauréat, l'exercice de \esp{verdadero o falso} exige deux choses : la réponse (\esp{Verdadero} / \esp{Falso}) et la \cle{justification} par une citation exacte du texte. Une réponse non justifiée ne vaut rien.

\begin{exoresolu}[Verdadero o falso : justifica con citas del texto]
Énoncé. Indique si les affirmations suivantes sont vraies ou fausses et justifie chaque réponse par une citation du texte.\par
a) \esp{La libertad de expresión no tiene ningún límite.}\par
b) \esp{La calumnia es castigada por la ley.}\par
c) \esp{El autor piensa que la censura es la mejor solución.}\par
d) \esp{La desinformación es siempre un error involuntario.}
\tcblower
\textbf{Solution détaillée.}\par
a) \esp{Falso. El texto dice en la línea 3: « este derecho no es absoluto ».}\par
(Faux : le texte affirme explicitement que ce droit n'est pas absolu.)\par
b) \esp{Verdadero. En las líneas 5-6 se lee: « la ley la castiga en España, en América Latina y en Guinea Ecuatorial ».}\par
(Vrai : le texte dit que la loi punit la calomnie.)\par
c) \esp{Falso. El autor afirma en la línea 14: « La solución no es prohibir, sino educar », y califica la censura de « arma peligrosa » (línea 13).}\par
(Faux : l'auteur rejette la censure et propose l'éducation.)\par
d) \esp{Falso. Según la línea 10, « La desinformación no es un simple error: es una estrategia de manipulación ».}\par
(Faux : la désinformation est présentée comme une stratégie délibérée, non comme une erreur.)
\end{exoresolu}

\begin{attention}[Pièges du vrai/faux]
\begin{itemize}
\item Ne recopiez pas trois lignes entières : citez uniquement le fragment qui prouve votre réponse.
\item Attention aux \cle{paraphrases piégeuses} de l'énoncé : « la desinformación es un error involuntario » contredit le texte, qui dit \esp{« no es un simple error »}.
\item Écrivez \esp{Verdadero} ou \esp{Falso} en toutes lettres, jamais une simple croix.
\item Méfiez-vous des mots \emph{absolus} glissés dans l'énoncé (\esp{siempre}, \esp{nunca}, \esp{ningún}, \esp{todos}) : ils rendent souvent l'affirmation fausse par rapport au texte.
\end{itemize}
\end{attention}

\courssub{Vocabulaire en contexte : trouver un synonyme dans le texte}

Cette question teste votre capacité à mobiliser le \cle{champ lexical} du texte. Il faut chercher un mot du texte qui a le même sens que le mot proposé.

\begin{exemplebox}[Synonymes à retrouver dans le texte]
\textbf{1. Busca en el texto un sinónimo de « mentira ».}\par
\esp{Un sinónimo de « mentira » es « bulo » (línea 8), que designa una noticia falsa difundida para engañar.}\par
« Un synonyme de "mensonge" est "bulo" (rumeur, intox). »

\textbf{2. Busca un sinónimo de « noticia falsa ».}\par
\esp{Un sinónimo de « noticia falsa » es « bulo »; también se puede hablar de « desinformación » (línea 10) cuando la falsedad es deliberada.}

\textbf{3. ¿Qué palabra del texto significa « hacer callar »?}\par
\esp{La palabra es « silenciar » (líneas 13-14): « ayer silenció a los periodistas honestos ».}
\end{exemplebox}

\begin{exemplebox}[Exemple de définition en contexte, traduite]
\esp{« La calumnia es acusar a alguien de un delito sabiendo que es falso »} (ligne 5) = « La calomnie, c'est accuser quelqu'un d'un délit en sachant que c'est faux. »\par
Remarquez la construction espagnole \esp{acusar a alguien de un delito} : la préposition \esp{de} introduit l'accusation, comme le « de » français dans « accuser quelqu'un \emph{de} quelque chose ». De même : \esp{« Lo acusaron de robar »} = « On l'a accusé d'avoir volé. »
\end{exemplebox}

\begin{center}
\begin{tabularx}{\linewidth}{|X|X|X|}
\hline
\rowcolor{popblueL} \textbf{Mot du texte} & \textbf{Français} & \textbf{Synonyme / famille} \\
\hline
el bulo & la rumeur, l'intox & la mentira, la noticia falsa \\
\hline
la difamación & la diffamation & dañar la reputación \\
\hline
la calumnia & la calomnie & acusar falsamente \\
\hline
la censura & la censure & prohibir, silenciar \\
\hline
la desinformación & la désinformation & la manipulación \\
\hline
engañar & tromper & mentir, manipular \\
\hline
\end{tabularx}
\end{center}

\begin{attention}[Ne confondez pas \esp{difamación} et \esp{calumnia}]
Les deux mots désignent des \cle{délits de presse}, mais ils ne sont pas synonymes :\par
\esp{la difamación} = dire du mal de quelqu'un, porter atteinte à sa réputation (par des propos ou des écrits) ;\par
\esp{la calumnia} = accuser faussement quelqu'un d'un \emph{délit} précis, en sachant que c'est faux.\par
Toute calomnie est une diffamation, mais toute diffamation n'est pas une calomnie. En version, traduire \esp{calumnia} par « diffamation » est une imprécision sanctionnée.
\end{attention}

\courssub{Questions d'interprétation}

Les questions d'interprétation demandent d'aller \emph{au-delà} de la citation : il faut expliquer les intentions de l'auteur avec ses propres mots.

\begin{exemplebox}[Interprétation : réponse modèle]
\textbf{¿Por qué opone el autor « derecho » y « responsabilidad »?}\par
\esp{El autor opone « derecho » y « responsabilidad » para mostrar que la libertad de expresión tiene dos caras: es un derecho que nos protege, pero también un deber hacia los demás. En las líneas 17-18 escribe: « un derecho acompañado de una responsabilidad. Quien habla debe respetar la verdad y la dignidad de los demás ». Así, quien usa su libertad para difamar o manipular traiciona ese derecho.}\par
« L'auteur oppose "droit" et "responsabilité" pour montrer que la liberté d'expression a deux visages : c'est un droit qui nous protège, mais aussi un devoir envers les autres. Celui qui utilise sa liberté pour diffamer ou manipuler trahit ce droit. »
\end{exemplebox}

\courssub{Question personnelle type Bac}

La dernière question est souvent \cle{personnelle} : on vous demande votre avis, justifié, en espagnol. Il n'y a pas de « bonne » opinion, mais il y a une bonne \emph{forme} : opinion claire + argument + exemple.

\begin{exemplebox}[¿Crees que las redes sociales deben ser censuradas? Justifica tu respuesta.]
\esp{En mi opinión, las redes sociales no deben ser censuradas, porque la censura puede convertirse en un instrumento contra la libertad, como advierte el texto: « la censura es un arma peligrosa ». Sin embargo, creo que las plataformas deben verificar las informaciones y eliminar los bulos que dañan a las personas. En mi país, por ejemplo, durante las elecciones circulan muchas noticias falsas por WhatsApp. Por eso pienso que la mejor solución es educar a los ciudadanos para que comprueben las fuentes antes de compartir.}\par
« À mon avis, les réseaux sociaux ne doivent pas être censurés, car la censure peut devenir un instrument contre la liberté. Cependant, je pense que les plateformes doivent vérifier les informations et supprimer les rumeurs qui nuisent aux personnes. Dans mon pays, par exemple, pendant les élections, beaucoup de fausses nouvelles circulent sur WhatsApp. C'est pourquoi je pense que la meilleure solution est d'éduquer les citoyens pour qu'ils vérifient les sources avant de partager. »
\end{exemplebox}

Observez la dernière phrase : \esp{« educar a los ciudadanos para que comprueben las fuentes »}. Après \esp{para que} (« pour que », but + sujet différent), l'espagnol impose le \cle{subjonctif} : \esp{comprueben} (subjonctif présent de \esp{comprobar}, verbe à changement radical o $\rightarrow$ ue). Ne dites jamais \esp{para que comprueban}.

\courssub{Méthode générale de réponse}

\begin{methode}[Répondre à une question de compréhension au Bac]
\begin{enumerate}
\item \textbf{Reformuler la question} dans le début de votre réponse : \esp{¿Cuál es la tesis del autor?} $\rightarrow$ \esp{La tesis del autor es que...}
\item \textbf{Citer le texte entre guillemets} avec le numéro de ligne : \esp{como dice en la línea 5: « ... »}.
\item \textbf{Expliquer avec ses propres mots} ce que la citation prouve : \esp{Esto significa que... / Es decir,...}
\end{enumerate}
Règle d'or : répondez toujours en \cle{phrases complètes en espagnol}, jamais par un seul mot, jamais en français.
\end{methode}

\begin{popfigure}
\begin{center}
\begin{tikzpicture}
\node[draw=popblue, fill=popblueL, rounded corners, minimum width=5.2cm, minimum height=1cm, align=center, font=\bfseries] (t1) at (0,0) {Tipo de pregunta};
\node[draw=poporange, fill=poporangeL, rounded corners, minimum width=7.6cm, minimum height=1cm, align=center, font=\bfseries] (t2) at (6.6,0) {Estrategia de respuesta};
\node[draw=popblue, fill=white, rounded corners, minimum width=5.2cm, minimum height=1.1cm, align=center] (a1) at (0,-1.2) {Global (tema, tesis, tono)};
\node[draw=poporange, fill=white, rounded corners, minimum width=7.6cm, minimum height=1.1cm, align=center] (a2) at (6.6,-1.2) {Relire le titre, le premier\\ paragraphe et la conclusion};
\node[draw=popblue, fill=white, rounded corners, minimum width=5.2cm, minimum height=1.1cm, align=center] (b1) at (0,-2.5) {Fina (verdadero/falso)};
\node[draw=poporange, fill=white, rounded corners, minimum width=7.6cm, minimum height=1.1cm, align=center] (b2) at (6.6,-2.5) {Citer la ligne exacte\\ entre guillemets};
\node[draw=popblue, fill=white, rounded corners, minimum width=5.2cm, minimum height=1.1cm, align=center] (c1) at (0,-3.8) {Vocabulario en contexto};
\node[draw=poporange, fill=white, rounded corners, minimum width=7.6cm, minimum height=1.1cm, align=center] (c2) at (6.6,-3.8) {Chercher le synonyme\\ dans le champ lexical du texte};
\node[draw=popblue, fill=white, rounded corners, minimum width=5.2cm, minimum height=1.1cm, align=center] (d1) at (0,-5.1) {Interpretación y opinión};
\node[draw=poporange, fill=white, rounded corners, minimum width=7.6cm, minimum height=1.1cm, align=center] (d2) at (6.6,-5.1) {Citer puis expliquer\\ avec ses propres mots};
\end{tikzpicture}
\end{center}
\legende{Pregunta / Estrategia de respuesta : la bonne méthode pour chaque type de question.}
\end{popfigure}

\begin{exercice}[À vous de jouer]
Répondez en espagnol, en phrases complètes, avec une citation et son numéro de ligne.\par
1. ¿Qué significa la expresión « antes de que la verdad se ponga los zapatos » (línea 9)?\par
2. Verdadero o falso : \esp{El autor propone prohibir las redes sociales.}\par
3. Busca en el texto un sinónimo de « engañar ».\par
4. ¿Estás de acuerdo con la tesis del autor? Justifica tu respuesta en tres frases.
\end{exercice}

\begin{corrige}[Exercice « À vous de jouer »]
1. \esp{Significa que las noticias falsas se difunden mucho más rápido que la verdad: un bulo « puede recorrer el mundo » (línea 9) antes de que se conozca la realidad.}\par
2. \esp{Falso. El autor dice en la línea 14: « La solución no es prohibir, sino educar ».}\par
3. \esp{Un sinónimo de « engañar » es « manipular » (línea 10: « una estrategia de manipulación que busca engañar al público »).}\par
4. Réponse libre, par exemple : \esp{Sí, estoy de acuerdo, porque la libertad sin responsabilidad puede destruir reputaciones. Además, la educación es más eficaz que la censura. Finalmente, verificar las fuentes protege la democracia.}
\end{corrige}

\begin{aretenir}[L'essentiel de la section]
\begin{itemize}
\item Réponse = \cle{phrase complète en espagnol} + \cle{citation entre guillemets} + numéro de ligne.
\item Question globale $\rightarrow$ titre, début, conclusion ; question fine $\rightarrow$ citation exacte.
\item \esp{Verdadero/Falso} non justifié = zéro point.
\item \esp{Asistir a} = être témoin de ; \esp{bulo} = rumeur, intox ; \esp{calumnia} = calomnie (accusation fausse d'un délit), à ne pas confondre avec \esp{difamación}.
\item Après \esp{antes de que} et \esp{para que} : subjonctif obligatoire.
\item Question personnelle : opinion + argument + exemple, en trois ou quatre phrases.
\end{itemize}
\end{aretenir}

\coursec{Lexique thématique par champs}

La lecture du texte « Libertad de expresión: ¿un derecho sin límites? » nous a permis de comprendre les grandes idées du débat : un droit fondamental, mais des abus punis par la loi. Pour commenter ce texte — et tout texte de presse — avec précision, il faut maintenant maîtriser son \cle{vocabulaire}. Un bon commentaire ne dit jamais \esp{el texto habla de cosas de los periódicos} : il dit \esp{el texto denuncia la difamación y los bulos}. Organisons donc ce lexique en quatre champs sémantiques, comme les quatre branches d'un même arbre.

\begin{popfigure}
\begin{tikzpicture}[mindmap, grow cyclic, every node/.style={concept, text=white, font=\small\bfseries},
root concept/.append style={concept color=popink, font=\bfseries},
level 1/.append style={level distance=3.4cm, sibling angle=90},
level 2/.append style={level distance=2.1cm, sibling angle=45, font=\scriptsize}]
\node[root concept]{Medios y libertad}
  child[concept color=popblue]{ node{Derechos}
    child{ node{opinión} }
    child{ node{prensa libre} }
    child{ node{fuente} }
  }
  child[concept color=poporange]{ node{Delitos}
    child{ node{difamación} }
    child{ node{calumnia} }
    child{ node{censura} }
  }
  child[concept color=poppurple]{ node{Desinformación}
    child{ node{bulo} }
    child{ node{rumor} }
    child{ node{verificar} }
  }
  child[concept color=popgreen]{ node{Publicidad}
    child{ node{anuncio} }
    child{ node{eslogan} }
    child{ node{campaña} }
  };
\end{tikzpicture}
\legende{Carte mentale du lexique : quatre champs sémantiques autour des médias et de la liberté.}
\end{popfigure}

\courssub{Champ 1 : la liberté et les droits}

Observons d'abord quelques phrases tirées de notre texte et de la presse hispanophone :

\esp{La libertad de expresión es un derecho fundamental.} « La liberté d'expression est un droit fondamental. »

\esp{Una prensa libre es esencial en una democracia.} « Une presse libre est essentielle dans une démocratie. »

\esp{El periodista no quiso revelar su fuente.} « Le journaliste n'a pas voulu révéler sa source. »

\begin{definition}[Le champ des droits]
\begin{itemize}
\item \esp{la libertad de expresión} : la liberté d'expression (droit de dire et d'écrire ce que l'on pense).
\item \esp{el derecho} : le droit. Attention au genre : on dit \esp{un derecho humano}, « un droit humain ».
\item \esp{la opinión} : l'opinion, l'avis. \esp{En mi opinión...} « À mon avis... »
\item \esp{la prensa libre} : la presse libre (non contrôlée par le pouvoir).
\item \esp{el/la periodista} : le/la journaliste. Même forme aux deux genres, seul l'article change.
\item \esp{la fuente} : la source (d'une information). \esp{fuente fiable} : source fiable.
\end{itemize}
\end{definition}

\begin{attention}[Derecho ou derecha ?]
Les francophones confondent souvent \esp{el derecho} (le droit) et \esp{la derecha} (la droite, direction ou tendance politique). \esp{Tengo derecho a opinar} signifie « J'ai le droit de donner mon avis », et non « j'ai la droite de donner mon avis ». De même, \esp{tener derecho a} + infinitif = « avoir le droit de » : \esp{Todos tienen derecho a informarse.} « Tous ont le droit de s'informer. »
\end{attention}

\courssub{Champ 2 : les délits de presse}

Quand la liberté d'expression nuit à autrui, la loi intervient. Voici le vocabulaire juridique indispensable :

\begin{definition}[Les délits de presse]
\begin{itemize}
\item \esp{la difamación} : la diffamation (accuser quelqu'un de faits faux qui nuisent à sa réputation).
\item \esp{la calumnia} : la calomnie (accuser faussement quelqu'un d'un délit ou d'un crime ; c'est une diffamation aggravée).
\item \esp{la injuria} : l'injure, l'outrage (parole offensante contre l'honneur d'une personne).
\item \esp{la censura} : la censure (interdiction de publier, exercée par une autorité).
\item \esp{demandar} : poursuivre en justice, intenter un procès. \esp{Demandaron al periódico.} « Ils ont poursuivi le journal en justice. »
\item \esp{condenar} : condamner (jugement). \esp{Fue condenado a pagar una multa.} « Il a été condamné à payer une amende. »
\item \esp{la multa} : l'amende.
\end{itemize}
\end{definition}

\begin{exemplebox}[Exemples traduits]
\esp{El periódico fue condenado por difamación.} « Le journal a été condamné pour diffamation. »

\esp{El político demandó al semanario por calumnia.} « L'homme politique a poursuivi l'hebdomadaire pour calomnie. »

\esp{La censura impidió la publicación del artículo.} « La censure a empêché la publication de l'article. »
\end{exemplebox}

\begin{attention}[Faux amis juridiques]
\begin{itemize}
\item \esp{demandar} ne signifie pas « demander » mais « poursuivre en justice ». « Demander » se dit \esp{pedir} (réclamer) ou \esp{preguntar} (poser une question).
\item \esp{condenar} = « condamner » au tribunal ; il s'emploie aussi au figuré pour « condamner » une idée, c'est-à-dire la désapprouver fermement : \esp{El autor condena la manipulación.} « L'auteur condamne la manipulation. »
\item \esp{la injuria} n'est pas « l'injure » au sens de blessure : une blessure se dit \esp{una herida}.
\end{itemize}
\end{attention}

\begin{propriete}[Censura previa et responsabilidad posterior]
Distinction capitale pour le commentaire : la \esp{censura previa} (censure préalable) interdit une publication \emph{avant} qu'elle paraisse ; la \esp{responsabilidad posterior} (responsabilité ultérieure) laisse publier librement, mais l'auteur répond \emph{après} de ses écrits devant les tribunaux. Les démocraties modernes refusent en général la censure préalable mais appliquent la responsabilité postérieure : \esp{Se puede publicar libremente, pero se responde ante la ley.} « On peut publier librement, mais on répond devant la loi. »
\end{propriete}

\courssub{Champ 3 : la désinformation}

\begin{definition}[Bulo]
Un \esp{bulo} est une fausse information créée et diffusée \emph{volontairement} pour tromper le public. C'est l'équivalent espagnol de « fake news » ou de « hoax ». Exemple : \esp{El bulo sobre las vacunas circuló por las redes sociales.} « La fausse information sur les vaccins a circulé sur les réseaux sociaux. »
\end{definition}

\begin{savaistu}[D'où vient le mot « bulo » ?]
Le mot \esp{bulo} désignait autrefois une ruse, un leurre. En Espagne, il s'est imposé dans le langage journalistique pour nommer les « fake news », alors qu'en Amérique latine on entend souvent aussi \esp{noticia falsa} ou l'anglicisme \esp{fake news}. La Real Academia Española l'a intégré à son dictionnaire avec le sens de « noticia falsa difundida con intención de engañar » (« fausse information diffusée dans l'intention de tromper »).
\end{savaistu}

\begin{definition}[Le champ de la désinformation]
\begin{itemize}
\item \esp{la desinformación} : la désinformation (ensemble de fausses informations diffusées pour tromper).
\item \esp{la noticia falsa} : la fausse information, la « fake news ».
\item \esp{la manipulación} : la manipulation (déformer les faits pour influencer l'opinion).
\item \esp{el rumor} : la rumeur (information non vérifiée qui circule de bouche à oreille).
\item \esp{el titular engañoso} : le titre trompeur (titre qui exagère ou déforme le contenu de l'article).
\item \esp{verificar} : vérifier. \esp{Hay que verificar las fuentes.} « Il faut vérifier les sources. »
\end{itemize}
\end{definition}

\begin{propriete}[Bulo ou error : la nuance d'intention]
La différence essentielle est l'\emph{intention}. Un \esp{error} est une faute involontaire : le journaliste s'est trompé de date ou de chiffre et publiera un \esp{fe de erratas} (un rectificatif). Un \esp{bulo} est fabriqué \emph{à dessein} pour tromper, nuire ou attirer des clics. Comparez :

\esp{El diario cometió un error en la cifra.} « Le journal a commis une erreur dans le chiffre. » (involontaire)

\esp{El bulo fue creado para dañar la imagen del candidato.} « La fausse information a été créée pour nuire à l'image du candidat. » (volontaire)
\end{propriete}

\begin{attention}[« Una noticia », pas « una nueva »]
Piège classique du francophone : \esp{una noticia} = « une nouvelle, une information » ; \esp{las noticias} = « les nouvelles » (le journal télévisé : \esp{ver las noticias} = « regarder les informations »). Le mot \esp{nueva} est l'adjectif « nouvelle » et ne peut pas désigner une information : on ne dit jamais « una nueva » pour « une nouvelle ». De même, « une nouvelle » littéraire (court récit) se dit \esp{un cuento} ou \esp{una novela corta}.
\end{attention}

\courssub{Champ 4 : la publicité}

\begin{definition}[Eslogan]
Un \esp{eslogan} est une phrase courte et frappante d'une publicité, facile à retenir, qui résume l'image d'une marque. Exemple célèbre : « Just do it ». En espagnol : \esp{El eslogan de la campaña es muy pegadizo.} « Le slogan de la campagne est très accrocheur. »
\end{definition}

\begin{definition}[Le champ de la publicité]
\begin{itemize}
\item \esp{la publicidad} : la publicité (l'activité, le secteur). Attention : « une publicité » (une annonce précise) se dit \esp{un anuncio}.
\item \esp{el anuncio} : l'annonce publicitaire, la « pub ». \esp{un anuncio de televisión} : « une publicité télévisée ».
\item \esp{el cartel} : l'affiche, le panneau publicitaire.
\item \esp{el/la consumidor(a)} : le/la consommateur(-trice).
\item \esp{el/la destinatario(a)} : le/la destinataire (la cible du message publicitaire).
\item \esp{la campaña publicitaria} : la campagne publicitaire.
\end{itemize}
\end{definition}

\begin{exemplebox}[Exemples traduits]
\esp{La campaña publicitaria está dirigida a los jóvenes.} « La campagne publicitaire s'adresse aux jeunes. »

\esp{El cartel anuncia un nuevo refresco.} « L'affiche fait la publicité d'un nouveau soda. »

\esp{Los consumidores deben leer bien las etiquetas.} « Les consommateurs doivent bien lire les étiquettes. »
\end{exemplebox}

\courssub{Les verbes clés et les expressions du commentaire}

\begin{center}
\begin{tabularx}{\linewidth}{|X|X|X|}\hline
\rowcolor{popblueL} Español & Français & Exemple traduit \\ \hline
difundir & diffuser, répandre & \esp{Difundir un bulo} : diffuser une fausse information \\ \hline
publicar & publier & \esp{Publicar un artículo} : publier un article \\ \hline
censurar & censurer & \esp{Censurar un reportaje} : censurer un reportage \\ \hline
manipular & manipuler & \esp{Manipular la información} : manipuler l'information \\ \hline
denunciar & dénoncer & \esp{Denunciar un abuso} : dénoncer un abus \\ \hline
informar & informer & \esp{Informar con objetividad} : informer avec objectivité \\ \hline
desmentir & démentir & \esp{Desmentir un rumor} : démentir une rumeur \\ \hline
\end{tabularx}
\end{center}

\begin{propriete}[Conjugaison de « desmentir », verbe irrégulier]
\esp{Desmentir} suit la conjugaison de \esp{mentir} (e $\rightarrow$ ie au présent, e $\rightarrow$ i au prétérit et au gérondif) : présent \esp{desmiento, desmientes, desmiente, desmentimos, desmentís, desmienten} ; prétérit \esp{desmentí, desmentiste, desmintió, desmentimos, desmentisteis, desmintieron}. \esp{El gobierno desmintió la noticia.} « Le gouvernement a démenti l'information. »
\end{propriete}

\begin{aretenir}[Expressions utiles pour le commentaire]
\begin{itemize}
\item \esp{según el autor} : selon l'auteur. \esp{Según el autor, la censura es inaceptable.} « Selon l'auteur, la censure est inacceptable. »
\item \esp{el texto denuncia...} : le texte dénonce... \esp{El texto denuncia la manipulación de los medios.} « Le texte dénonce la manipulation des médias. »
\item \esp{se trata de...} : il s'agit de... \esp{Se trata de un artículo de opinión.} « Il s'agit d'un article d'opinion. »
\item \esp{en definitiva} : en définitive, en somme. \esp{En definitiva, el derecho tiene límites.} « En définitive, le droit a des limites. »
\end{itemize}
\end{aretenir}

\courssub{Texte d'application : le lexique en contexte}

\begin{textebox}[« El bulo que recorrió Yaoundé »]
La semana pasada, un bulo recorrió las redes sociales de Yaoundé: un mensaje anunciaba que un famoso jugador de fútbol regalaría dinero a quienes compartieran la noticia. En pocas horas, miles de personas difundieron la información sin verificar la fuente. Un periodista de un diario local independiente investigó el caso y descubrió que el titular era engañoso: la página original no existía y el anuncio imitaba el eslogan de una conocida campaña publicitaria.

El periodista publicó un artículo para desmentir el rumor y denunciar la manipulación. Según el autor, se trata de un ejemplo claro de desinformación: los creadores del bulo querían obtener datos personales de los consumidores. En definitiva, el texto recuerda que la libertad de expresión es un derecho, pero que difundir noticias falsas puede ser un delito. Si el bulo daña la reputación de una persona, su autor puede ser demandado por difamación y condenado a pagar una multa.

La lección es sencilla: antes de compartir, hay que verificar. Una sociedad informada es la mejor defensa contra la mentira.
\end{textebox}

\textbf{Glossaire.} \esp{regalar} : offrir ; \esp{compartieran} : partagent (subjonctif imparfait, après \esp{quienes} à valeur indéfinie) ; \esp{descubrió} : a découvert ; \esp{obtener datos personales} : obtenir des données personnelles ; \esp{dañar} : nuire à ; \esp{compartir} : partager ; \esp{la mentira} : le mensonge.

\textbf{Questions de compréhension.}
\begin{enumerate}
\item ¿Qué anunciaba el bulo que circuló por Yaoundé ?
\item ¿Cómo descubrió el periodista que la noticia era falsa ?
\item ¿Qué delito puede cometer el autor de un bulo que daña la reputación de alguien ?
\item ¿Cuál es, en definitiva, la lección del texto ?
\end{enumerate}

\begin{exoresolu}[Traduction et vocabulaire]
Énoncé. Traduisez en espagnol : « Le journal a été condamné pour diffamation. Il faut vérifier les sources avant de diffuser une information. » \tcblower
Solution. Première phrase : « condamné » se dit \esp{condenado} ; « pour diffamation » se traduit par \esp{por difamación} (la préposition \esp{por} exprime ici la cause). On obtient : \esp{El periódico fue condenado por difamación.} Deuxième phrase : « il faut » se rend par \esp{hay que} + infinitif ; « vérifier les sources » : \esp{verificar las fuentes} ; « avant de » + infinitif : \esp{antes de} + infinitif ; « diffuser une information » : \esp{difundir una noticia} (et non « una nueva » !). On obtient : \esp{Hay que verificar las fuentes antes de difundir una noticia.}
\end{exoresolu}

\begin{exercice}[À vous de jouer]
\begin{enumerate}
\item Complétez avec le mot du champ lexical qui convient : \esp{difamación, bulo, eslogan, fuente, multa, censura}.
\begin{enumerate}
\item El \_\_\_\_\_\_\_\_ de la marca « Just do it » es conocido en todo el mundo.
\item El periodista protege siempre el nombre de su \_\_\_\_\_\_\_\_.
\item El juez impuso una \_\_\_\_\_\_\_\_ de dos millones de francos al semanario.
\item El mensaje sobre el jugador de fútbol era un \_\_\_\_\_\_\_\_.
\item La revista fue condenada por \_\_\_\_\_\_\_\_: había acusado al ministro sin pruebas.
\item Durante la dictadura, la \_\_\_\_\_\_\_\_ impedía publicar cualquier crítica.
\end{enumerate}
\item Traduisez en français : \esp{El texto denuncia la manipulación de los consumidores por la publicidad engañosa.}
\item Traduisez en espagnol : « Selon l'auteur, il s'agit d'une fausse information créée pour tromper les consommateurs. »
\end{enumerate}
\end{exercice}

\begin{corrige}[Exercice « À vous de jouer »]
1. a) \esp{eslogan} ; b) \esp{fuente} ; c) \esp{multa} ; d) \esp{bulo} ; e) \esp{difamación} ; f) \esp{censura}. 2. « Le texte dénonce la manipulation des consommateurs par la publicité trompeuse. » 3. \esp{Según el autor, se trata de una noticia falsa (un bulo) creada para engañar a los consumidores.} Remarques : « tromper » se dit \esp{engañar} ; « pour » + infinitif de but se traduit par \esp{para} + infinitif ; devant « les consommateurs », complément d'objet direct de personnes, il faut la préposition \esp{a} personnelle : \esp{a los consumidores}.
\end{corrige}

Ce lexique maîtrisé, nous sommes prêts à l'utiliser concrètement : la prochaine étape consistera à analyser une publicité réelle, en repérant son \esp{eslogan}, son \esp{destinatario} et ses procédés de persuasion, avant d'apprendre à démasquer les \esp{bulos}.

\coursec{Analyse d'une publicité}

Dans la section précédente, nous avons constitué notre \cle{lexique thématique} (\esp{difamación}, \esp{bulo}, \esp{publicidad}, \esp{eslogan}, \esp{manipulación}...). Ces mots ne sont pas de simples définitions à mémoriser : ce sont des \textbf{outils d'analyse}. Nous allons maintenant les mettre en action sur l'objet médiatique le plus présent dans notre quotidien : la \cle{publicité}. À la télévision de Yaoundé, sur les murs de Douala, à la radio, sur les réseaux sociaux, les annonces nous sollicitent sans cesse. Savoir les \emph{décortiquer}, c'est refuser d'être un consommateur passif : c'est exactement la compétence que l'examen attend de vous.

\courssub{Observer d'abord : deux publicités sous nos yeux}

Avant toute règle, lisons deux textes publicitaires originaux, l'un espagnol, l'autre camerounais. Observez leur construction : nous les analyserons ensuite méthodiquement.

\begin{textebox}[Anuncio 1 : « SolAndaluz », refresco español]
¿Sed? ¿Calor? ¿Rutina? SolAndaluz tiene la respuesta.\\
Bebe frescura, vive el verano.\\
Nuestro refresco de naranja, con burbujas naturales y sin colorantes artificiales, te devuelve la energía del sur en cada sorbo. 9 de cada 10 jóvenes andaluces lo eligen para sus tardes con amigos.\\
María, 19 años, estudiante de Sevilla, lo confirma : « Desde que descubrí SolAndaluz, mis veranos tienen otro sabor. No puedo imaginar una fiesta sin él. »\\
SolAndaluz. Vive tu vida.\\
Disponible en todos los supermercados. Compra ahora y participa en el sorteo de un viaje a la costa.
\end{textebox}

\begin{textebox}[Anuncio 2 : « KamerFone », operador camerunés]
¡Camerún entero está conectado!\\
Con KamerFone, llama a tu familia de Bafoussam a Maroua sin preocuparte por el saldo. Nuestra red cubre las diez regiones del país, desde las playas de Kribi hasta los mercados de Bamenda.\\
Más de 8 millones de cameruneses ya confían en nosotros. ¿Y tú, qué esperas?\\
Ngono, comerciante del mercado Mokolo de Yaundé, testifica : « Con KamerFone, hablo cada día con mi hijo que estudia en Duala. La red nunca falla. »\\
Activa ya tu tarjeta SIM KamerFone : 100 \% de bonificación en tu primera recarga del mes.\\
KamerFone. Juntos, más cerca.
\end{textebox}

\textbf{Glossaire :} \esp{refresco} = boisson rafraîchissante ; \esp{burbujas} = bulles ; \esp{sorbo} = gorgée ; \esp{sorteo} = tirage au sort ; \esp{saldo} = crédit téléphonique ; \esp{recarga} = recharge ; \esp{bonificación} = bonus, avantage supplémentaire ; \esp{tarjeta SIM} = carte SIM ; \esp{testifica} = témoigne ; \esp{la red} = le réseau.

\textbf{Questions d'observation :} Qui parle dans chaque texte ? À qui s'adresse-t-il ? Que promet-il ? Quels procédés reconnaissez-vous (chiffres, témoignage, impératifs) ?

\courssub{La grille d'analyse : les cinq composantes de l'annonce}

Toute publicité, qu'elle soit espagnole ou camerounaise, peut être décrite avec la même grille. Retenez cinq éléments : l'\cle{émetteur} (\esp{emisor}), le \cle{destinataire} (\esp{destinatario}), le \cle{message} (\esp{mensaje}), le \cle{support} (\esp{soporte}) et l'\cle{eslogan} avec l'\cle{imagen}.

\begin{definition}[Les cinq composantes]
\begin{itemize}
\item \esp{Emisor} : la marque ou l'entreprise qui paie et qui parle (\esp{SolAndaluz}, \esp{KamerFone}).
\item \esp{Destinatario} : la cible visée (\esp{jóvenes}, \esp{familias}, \esp{comerciantes}...).
\item \esp{Mensaje} : ce que dit l'annonce. On distingue le \textbf{message explicite} (ce qui est dit clairement : « sans colorants », « 100 \% de bonificación ») et le \textbf{message implicite} (ce qui est suggéré : « si tu bois SolAndaluz, tu auras des amis et des fêtes »).
\item \esp{Soporte} : le canal de diffusion : \esp{prensa} (presse écrite), \esp{televisión}, \esp{radio}, \esp{redes sociales}, \esp{carteles} (affiches).
\item \esp{Eslogan e imagen} : la phrase-choc et le travail visuel qui rendent le message mémorable.
\end{itemize}
\end{definition}

\begin{popfigure}
\begin{tikzpicture}[mindmap, grow cyclic, every node/.style={concept, align=center}, concept color=popblueL, text=popdark,
  root concept/.append style={concept color=poporange, text=white, font=\Large\bfseries},
  level 1/.append style={level distance=3.2cm, sibling angle=72},
  level 1 concept/.append style={concept color=popblue!40, font=\small\bfseries}]
\node[root concept]{ANUNCIO}
  child { node[align=center]{Emisor\\ \emph{la marca}} }
  child { node[align=center]{Destinatario\\ \emph{la cible}} }
  child { node[align=center]{Mensaje\\ \emph{explícito / implícito}} }
  child { node[align=center]{Eslogan\\ \emph{la phrase-choc}} }
  child { node[align=center]{Imagen\\ \emph{colores, miradas}} };
\end{tikzpicture}
\legende{Carte mentale : les cinq branches de l'analyse d'une publicité.}
\end{popfigure}

\begin{methode}[Analyser une publicité en quatre questions]
Face à une publicité à l'examen ou dans la vie, posez-vous toujours, dans l'ordre :
\begin{enumerate}
\item \textbf{¿Quién habla?} Identifiez l'émetteur : quelle marque, quel produit, quel intérêt économique ?
\item \textbf{¿A quién?} Déterminez la cible : âge, milieu social, besoins, désirs. Une boisson vise les jeunes ; un opérateur téléphonique vise les familles et les commerçants.
\item \textbf{¿Qué promete?} Distinguez la promesse explicite (le produit, le prix, le bonus) de la promesse implicite (le bonheur, la réussite, l'amitié, la modernité).
\item \textbf{¿Cómo convence?} Repérez les procédés : eslogan, impératifs, chiffres, témoignage, couleurs, émotions. C'est ici que se joue la \esp{persuasión}.
\end{enumerate}
Rédigez ensuite votre analyse en suivant ce même ordre : c'est le plan tout fait de votre paragraphe de commentaire.
\end{methode}

\courssub{L'eslogan : la phrase qui reste dans la tête}

L'\cle{eslogan} est la formule courte et frappante qui résume la promesse de la marque. Observez nos deux exemples : \esp{Bebe frescura, vive el verano} et \esp{Juntos, más cerca}. Qu'ont-ils en commun ? Ils sont brefs, rythmés, et ils commandent ou suggèrent.

\begin{propriete}[Les procédés typiques de l'eslogan]
\begin{itemize}
\item \textbf{L'impératif} : il crée un rapport direct et dynamique avec le consommateur. \esp{Compra ahora} (Achète maintenant), \esp{Vive tu vida} (Vis ta vie), \esp{Bebe frescura} (Bois la fraîcheur).
\item \textbf{Le parallélisme} : deux structures identiques donnent du rythme. \esp{Bebe frescura, vive el verano} : impératif + nom, impératif + nom.
\item \textbf{La rime et l'assonance} : elles rendent la phrase mémorable. \esp{SolAndaluz, energía y luz} (rime en \emph{-uz}).
\item \textbf{Le jeu de mots} (\esp{juego de palabras}) : il amuse et flatte l'intelligence du lecteur. \esp{KamerFone : la red que nunca se desconecta de ti} (double sens : le réseau technique et le lien affectif).
\item \textbf{Le tutoiement} : en espagnol, l'impératif de \esp{tú} (\esp{bebe}, \esp{vive}, \esp{compra}) crée la proximité ; le vouvoiement (\esp{compre usted}) est réservé aux produits de luxe ou aux services bancaires.
\end{itemize}
\end{propriete}

\begin{propriete}[Rappel de grammaire : former l'impératif affirmatif]
L'impératif est le mode roi de la publicité. Deux formes à maîtriser :
\begin{itemize}
\item \textbf{Tú} (tutoiement) : on reprend la 3\up{e} personne du singulier du présent de l'indicatif : \esp{beber} $\rightarrow$ \esp{bebe}, \esp{vivir} $\rightarrow$ \esp{vive}, \esp{comprar} $\rightarrow$ \esp{compra}, \esp{activar} $\rightarrow$ \esp{activa}, \esp{llamar} $\rightarrow$ \esp{llama}.
\item \textbf{Usted} (vouvoiement) : on utilise le présent du subjonctif : \esp{compre usted}, \esp{visite usted}, \esp{llame usted}.
\item \textbf{Exceptions} à l'impératif \esp{tú} : \esp{decir} $\rightarrow$ \esp{di}, \esp{hacer} $\rightarrow$ \esp{haz}, \esp{ir} $\rightarrow$ \esp{ve}, \esp{poner} $\rightarrow$ \esp{pon}, \esp{salir} $\rightarrow$ \esp{sal}, \esp{ser} $\rightarrow$ \esp{sé}, \esp{tener} $\rightarrow$ \esp{ten}, \esp{venir} $\rightarrow$ \esp{ven}.
\end{itemize}
Exemple publicitaire : \esp{Sé tú mismo, ven y compra.} = « Sois toi-même, viens et achète. »
\end{propriete}

\begin{exemplebox}[Eslogan décortiqué]
\esp{Bebe frescura, vive el verano.} = « Bois la fraîcheur, vis l'été. »\\
Analyse : deux \textbf{impératifs} (\esp{bebe}, \esp{vive}) en \textbf{parallélisme} parfait ; la boisson est assimilée à une sensation (\esp{frescura}) puis à une saison entière (\esp{el verano}) : c'est le passage typique du produit réel au bonheur promis. Le message implicite : boire ce soda, c'est vivre pleinement.\\
Autre exemple : \esp{KamerFone. Juntos, más cerca.} = « Ensemble, plus proches. » : deux adverbes, aucun verbe, une promesse affective (la proximité familiale) plutôt que technique.
\end{exemplebox}

\begin{attention}[Ne confondez pas « publicidad » et « propaganda »]
\esp{Publicidad} = la \textbf{publicité}, à visée commerciale (vendre un produit). \esp{Propaganda} = la \textbf{propagande}, à visée politique ou idéologique (faire adhérer à des idées). Un francophone qui traduit \esp{propaganda} par « publicité » commet un faux sens. À l'inverse, ne dites jamais \esp{propaganda de un refresco} pour « la publicité d'un soda » : dites \esp{publicidad de un refresco} ou \esp{un anuncio de un refresco}. Autre piège : \esp{un anuncio} = une annonce/un spot publicitaire, jamais « une annonciation ».
\end{attention}

\courssub{L'image : couleurs, regards et stéréotypes}

Même sans voir l'image, le texte publicitaire la décrit ou la suggère. Trois questions suffisent :

\begin{itemize}
\item \textbf{Les couleurs} (\esp{los colores}) : l'orange et le jaune évoquent le soleil, l'énergie, l'été (SolAndaluz) ; le vert et le jaune peuvent rappeler les couleurs nationales (une marque camerounaise joue sur la fierté du pays).
\item \textbf{Les regards} (\esp{las miradas}) : un personnage qui regarde le lecteur droit dans les yeux crée de la complicité ; un personnage qui regarde l'horizon invite au rêve.
\item \textbf{La mise en scène et les stéréotypes} (\esp{los estereotipos}) : la jeunesse souriante en fête pour les boissons, la mère de famille attentionnée pour les produits alimentaires, l'homme d'affaires pressé pour les téléphones. Repérer ces clichés, c'est déjà faire preuve d'esprit critique : la publicité montre rarement la vie réelle, elle montre la vie \emph{rêvée}.
\end{itemize}

\begin{center}
\begin{tabularx}{\linewidth}{|X|X|X|}
\hline
\rowcolor{popblueL} \textbf{Élément} & \textbf{SolAndaluz (Espagne)} & \textbf{KamerFone (Cameroun)} \\
\hline
Emisor & Marque de boisson & Opérateur téléphonique \\
\hline
Destinatario & Jeunes (15-25 ans) & Familles, commerçants, tout le pays \\
\hline
Promesa explícita & Sans colorants, sorteo & Couverture nationale, bonificación del 100 \% \\
\hline
Promesa implícita & Amitié, fêtes, été idéal & Unité nationale, lien familial \\
\hline
Cifras & « 9 de cada 10 jóvenes » & « Más de 8 millones » \\
\hline
Testimonio & María, 19 años, Sevilla & Ngono, mercado Mokolo \\
\hline
Emoción & Joie, évasion & Fierté, solidarité familiale \\
\hline
Eslogan & « Vive tu vida » & « Juntos, más cerca » \\
\hline
\end{tabularx}
\end{center}

Ce tableau comparatif montre une vérité essentielle : \textbf{les procédés sont universels}, seuls les décors culturels changent. La plage andalouse et le marché Mokolo jouent exactement le même rôle : ancrer le produit dans un quotidien que la cible reconnaît.

\courssub{Les procédés de persuasion : comment l'annonce nous convainc}

\begin{propriete}[Les cinq grands procédés de persuasion]
\begin{enumerate}
\item \textbf{Le témoignage} (\esp{el testimonio}) : un consommateur « ordinaire » affirme que le produit a changé sa vie (\esp{María}, \esp{Ngono}). Il crée l'identification : « si ça marche pour elle, ça marchera pour moi ».
\item \textbf{Les chiffres} (\esp{las cifras}) : \esp{9 de cada 10 jóvenes}, \esp{más de 8 millones}. Ils donnent une apparence scientifique et créent l'effet de masse : « tout le monde le fait ». Attention : la source de ces chiffres est rarement indiquée.
\item \textbf{La peur} (\esp{el miedo}) : « sans ce produit, vous serez en retard, isolé, en danger ». Très utilisée pour les assurances et la sécurité.
\item \textbf{L'humour} (\esp{el humor}) : il désamorce la méfiance et rend la marque sympathique.
\item \textbf{L'appel aux émotions} (\esp{la apelación a las emociones}) : famille, amour, fierté nationale, nostalgie. C'est le procédé le plus puissant, car on achète souvent avec le cœur, pas avec la raison.
\end{enumerate}
\end{propriete}

\begin{exoresolu}[Repérer les procédés dans un eslogan]
Énoncé : dans l'annonce KamerFone, relevez trois procédés de persuasion et expliquez leur effet sur le lecteur, en rédigeant quatre à six lignes en espagnol.
\tcblower
Solution détaillée :\\
\esp{En primer lugar, el anuncio utiliza las cifras : « más de 8 millones de cameruneses ». Este dato crea un efecto de masa : el lector piensa que debe confiar porque todos confían ya. En segundo lugar, encontramos el testimonio de Ngono, comerciante del mercado Mokolo : este personaje cercano y realista permite la identificación del público camerunés. Por último, el eslogan « Juntos, más cerca » apela a las emociones, en particular al amor familiar, más que a las cualidades técnicas del producto. Así, la publicidad no vende un teléfono : vende la proximidad con la familia.}\\
Traduction : « Premièrement, l'annonce utilise les chiffres : "plus de 8 millions de Camerounais". Cette donnée crée un effet de masse : le lecteur pense qu'il doit faire confiance parce que tous font déjà confiance. Deuxièmement, on trouve le témoignage de Ngono, commerçante du marché Mokolo : ce personnage proche et réaliste permet l'identification du public camerounais. Enfin, le slogan "Ensemble, plus proches" fait appel aux émotions, en particulier à l'amour familial, plutôt qu'aux qualités techniques du produit. Ainsi, la publicité ne vend pas un téléphone : elle vend la proximité avec la famille. »
\end{exoresolu}

\courssub{Esprit critique : publicité, information et « publicidad engañosa »}

Dernière étape, la plus importante pour le commentaire : distinguer \cle{publicité} et \cle{information}. L'\textbf{information} cherche à décrire la réalité de manière vérifiable ; la \textbf{publicité} cherche à vendre, et elle a le droit d'exagérer, de sélectionner, d'embellir. Mais il existe une limite : la \cle{publicidad engañosa}, la publicité mensongère, qui affirme des faits faux (un produit « 100 \% natural » qui ne l'est pas, un bonus qui n'existe pas, des chiffres inventés).

\begin{savaistu}[Qui surveille la publicité ?]
En Espagne, l'association \textbf{Autocontrol} réunit les annonceurs et veille à ce que la publicité soit « veraz » (véridique) : un consommateur peut porter plainte contre une \esp{publicidad engañosa}. Au Cameroun aussi, le consommateur n'est pas désarmé : il peut dénoncer les fausses promesses publicitaires auprès des structures de défense des consommateurs et des autorités de régulation de la communication. Retenez ce vocabulaire : \esp{denunciar} (dénoncer), \esp{una queja} (une plainte), \esp{engañar} (tromper), \esp{veraz} (véridique).
\end{savaistu}

\begin{attention}[Pièges de traduction et de jugement]
\begin{itemize}
\item \esp{Engañoso} = trompeur, mensonger (faux ami partiel : ce n'est pas « engageant » !). \esp{Una publicidad engañosa} = une publicité mensongère.
\item Ne confondez pas \esp{el anunciante} (l'annonceur, celui qui paie la pub) et \esp{el anuncio} (l'annonce elle-même).
\item Dans votre commentaire, évitez le jugement naïf « cette publicité est bien ». Dites plutôt : \esp{Esta publicidad es eficaz porque utiliza el testimonio y las cifras, pero el mensaje implícito exagera la felicidad que promete el producto.} (« Cette publicité est efficace parce qu'elle utilise le témoignage et les chiffres, mais le message implicite exagère le bonheur que le produit promet. »)
\end{itemize}
\end{attention}

\begin{exercice}[À vous de jouer]
Rédigez en espagnol (10 à 15 lignes) l'analyse de l'annonce « SolAndaluz » en suivant les quatre questions de la méthode : ¿Quién habla? ¿A quién? ¿Qué promete? ¿Cómo convence? Terminez par une phrase d'esprit critique sur la frontière entre persuasion et \esp{publicidad engañosa}.
\end{exercice}

\begin{corrige}[Exercice : analyse de « SolAndaluz »]
Éléments de réponse attendus : l'émetteur est la marque SolAndaluz, qui vend un soda à l'orange ; la cible est la jeunesse (\esp{jóvenes andaluces}) ; la promesse explicite est un produit naturel (\esp{sin colorantes}) et un jeu-concours (\esp{sorteo}) ; la promesse implicite est l'amitié, la fête et l'été parfait ; les procédés sont l'impératif et le parallélisme de l'eslogan (\esp{Bebe frescura, vive el verano}), le chiffre (\esp{9 de cada 10}), le témoignage de María et l'appel aux émotions. Phrase critique possible : \esp{El anuncio no miente abiertamente, pero sugiere que una bebida puede dar la felicidad : ahí empieza la exageración publicitaria.} (« L'annonce ne ment pas ouvertement, mais elle suggère qu'une boisson peut donner le bonheur : c'est là que commence l'exagération publicitaire. »)
\end{corrige}

\begin{aretenir}[L'essentiel de la section]
Analyser une publicité, c'est répondre à quatre questions : \esp{¿Quién habla? ¿A quién? ¿Qué promete? ¿Cómo convence?} L'eslogan repose sur l'impératif, le rythme et le parallélisme ; l'image sur les couleurs, les regards et les stéréotypes ; la persuasion sur le témoignage, les chiffres, la peur, l'humour et les émotions. Enfin, l'esprit critique consiste à distinguer information et publicité, et à savoir nommer la limite : la \esp{publicidad engañosa}. Ces outils vous serviront directement dans la section suivante, consacrée à la détection des fausses informations.
\end{aretenir}

\coursec{Détecter les fake news}

Après avoir appris à décrypter les techniques de la publicité (l'\esp{eslogan}, l'image, l'argument d'autorité), nous disposons déjà d'un œil critique. Mais la publicité, elle, assume son intention : vendre. Il existe un danger plus insidieux sur nos téléphones et nos réseaux sociaux : les fausses informations, qui se déguisent en vraies nouvelles pour nous tromper. À Yaoundé comme à Madrid, un message WhatsApp peut faire le tour d'un quartier en dix minutes. Apprendre à détecter les \esp{fake news} est donc devenu une compétence de citoyen — et une compétence exigée au baccalauréat, où l'on vous demandera de distinguer le fait de l'opinion et de juger la fiabilité d'un texte.

\courssub{Qu'est-ce qu'une fake news ?}

\begin{textebox}[Un mensaje que circula por WhatsApp]
¡URGENTE! ¡COMPARTE CON TODOS TUS CONTACTOS!\\
El Gobierno va a cerrar todos los mercados del país durante tres meses a partir del lunes. Un médico famoso lo ha confirmado: hay una epidemia terrible y los alimentos están contaminados. Mi primo trabaja en el ministerio y me lo ha dicho. Compra arroz y aceite AHORA antes de que se acabe todo. Los periodistas no dicen nada porque tienen miedo. ¡No esperes a que sea demasiado tarde! Reenvía este mensaje a tu familia. ¡Dios te bendiga!
\end{textebox}

\textbf{Glossaire} : \esp{comparte} = partage (impératif) ; \esp{alimentos} = aliments ; \esp{contaminados} = contaminés ; \esp{se acabe} = qu'il n'y en ait plus (subjonctif) ; \esp{reenvía} = transfère (impératif) ; \esp{tienen miedo} = ils ont peur (idiome \esp{tener miedo}).

\textbf{Questions d'observation} :
\begin{enumerate}
\item Ce message cite-t-il une source vérifiable (un journal, un site officiel, un nom complet) ?
\item Quelle est sa date de publication ?
\item Quels mots cherchent à provoquer la peur et l'urgence ?
\item Que vous demande-t-on de faire immédiatement ?
\end{enumerate}

Vous l'avez deviné : ce message est un \esp{bulo} typique. Observons maintenant le vocabulaire précis.

\begin{definition}[Fake news, bulo, desinformación]
En espagnol, on emploie trois mots complémentaires :
\begin{itemize}
\item \esp{las \cle{fake news}} (anglicisme très courant) : les fausses informations ;
\item \esp{el \cle{bulo}} : le mot espagnol traditionnel pour « canular, fausse rumeur » ;
\item \esp{la \cle{desinformación}} : la désinformation, c'est-à-dire l'action de diffuser de fausses informations \textbf{volontairement}, dans l'intention de tromper.
\end{itemize}
Trois caractéristiques définissent une fake news : elle est \textbf{fausse}, elle est créée avec l'\textbf{intention de tromper} (ou de gagner de l'argent grâce aux clics), et elle se diffuse à une \textbf{vitesse extraordinaire} sur les réseaux sociaux, beaucoup plus vite qu'une vraie nouvelle vérifiée.
\end{definition}

\begin{attention}[Desinformar ou estar mal informado ?]
Ne confondez pas les deux :
\begin{itemize}
\item \esp{desinformar} = désinformer \textbf{volontairement}, mentir pour manipuler. \esp{Esa página web desinforma a la gente.} « Ce site web désinforme les gens. »
\item \esp{estar mal informado} = être mal informé, sans mauvaise intention. \esp{Mi vecino está mal informado, pero no miente.} « Mon voisin est mal informé, mais il ne ment pas. »
\end{itemize}
La nuance d'\textbf{intention} est essentielle : celui qui partage un \esp{bulo} par naïveté est mal informé ; celui qui le fabrique désinforme.
\end{attention}

\courssub{Les signaux d'alerte}

Comment reconnaître une fausse information au premier coup d'œil ? Voici les cinq signaux d'alerte classiques, à connaître en espagnol :

\begin{center}
\begin{tabularx}{\linewidth}{|X|X|X|}\hline
\rowcolor{popblueL} Señal de alerta & Français & Exemple \\ \hline
\esp{titular exagerado} & titre exagéré, sensationnaliste & \esp{¡INCREÍBLE! ¡NO LO VAS A CREER!} \\ \hline
\esp{falta de fuente} & absence de source identifiable & \esp{« un médico famoso », « mi primo »} \\ \hline
\esp{fecha antigua o ausente} & date ancienne ou absente & une vieille nouvelle ressortie comme neuve \\ \hline
\esp{foto trucada o fuera de contexto} & photo truquée ou sortie de son contexte & \esp{La foto era de otro país.} \\ \hline
\esp{errores de ortografía} & fautes d'orthographe et de grammaire & majuscules partout, accents oubliés \\ \hline
\end{tabularx}
\end{center}

Reprenons notre message WhatsApp : il cumule \textbf{tous} les signaux — titre en majuscules (\esp{¡URGENTE!}), source anonyme (\esp{mi primo trabaja en el ministerio}), aucune date, appel à la peur, et injonction de partager. Un vrai article de presse, lui, donne un titre mesuré, le nom du journaliste, la date, le lieu et des sources citées.

\begin{exemplebox}[Exemples traduits]
\esp{No compartas una noticia sin verificarla.} « Ne partage pas une information sans l'avoir vérifiée. » (Remarque : impératif négatif + subjonctif, \esp{compartas}.)\\
\esp{La foto era de otro país.} « La photo venait d'un autre pays. »\\
\esp{Ese titular es exagerado para llamar la atención.} « Ce titre est exagéré pour attirer l'attention. »\\
\esp{El artículo no cita ninguna fuente fiable.} « L'article ne cite aucune source fiable. »
\end{exemplebox}

\courssub{La méthode de vérification en cinq étapes}

Face à une information douteuse, le journaliste comme le citoyen suivent une démarche en cinq étapes :

\begin{enumerate}
\item \textbf{Lire au-delà du titre} (\esp{leer más allá del titular}) : le titre est fait pour choquer ; le contenu dit souvent autre chose, ou rien du tout.
\item \textbf{Vérifier la source et l'auteur} (\esp{verificar la fuente y el autor}) : qui publie ? Un média connu, un site officiel, ou une page anonyme ? L'auteur existe-t-il vraiment ?
\item \textbf{Vérifier la date} (\esp{verificar la fecha}) : une vraie nouvelle de 2015 ressortie en 2025 devient une fausse nouvelle.
\item \textbf{Recouper avec d'autres médias} (\esp{contrastar con otros medios}) : si une seule page parle d'un « événement majeur » et qu'aucun grand média n'en parle, méfiance.
\item \textbf{Vérifier l'image} (\esp{verificar la imagen}) : grâce à la \esp{búsqueda inversa de imágenes} (recherche inversée), on découvre souvent que la photo vient d'un autre pays ou d'un autre événement.
\end{enumerate}

\begin{methode}[La règle des 5 V]
Retenez la règle des \textbf{5 V}, facile à mémoriser en espagnol :
\begin{enumerate}
\item \textbf{V}er el titular (regarder le titre : est-il exagéré ?)
\item \textbf{V}erificar la fuente (vérifier la source et l'auteur)
\item \textbf{V}er la fecha (regarder la date)
\item \textbf{V}erificar con otros medios (recouper avec d'autres médias)
\item \textbf{V}erificar la imagen (vérifier l'image par recherche inversée)
\end{enumerate}
Si l'information échoue à une seule de ces étapes, la règle d'or s'applique : \esp{no compartir} — ne pas partager.
\end{methode}

\begin{popfigure}
\begin{tikzpicture}[
  node distance=0.55cm,
  etapa/.style={rectangle, rounded corners, draw=popblue, fill=popblueL, text width=6.5cm, align=center, font=\small, inner sep=5pt},
  decision/.style={diamond, draw=poporange, fill=poporangeL, aspect=2.2, align=center, font=\small\bfseries, inner sep=2pt},
  final/.style={rectangle, rounded corners, draw=popgreen, fill=popgreenL, align=center, font=\small\bfseries, inner sep=5pt},
  finalno/.style={rectangle, rounded corners, draw=poppink, fill=poppinkL, align=center, font=\small\bfseries, inner sep=5pt},
  flecha/.style={-{Stealth[length=2.5mm]}, thick, popdark}
]
\node[decision, align=center] (q){¿Es verdad\\esta noticia?};
\node[etapa, below=of q] (e1) {1. Leer más allá del titular};
\node[etapa, below=of e1] (e2) {2. Verificar la fuente y el autor};
\node[etapa, below=of e2] (e3) {3. Verificar la fecha};
\node[etapa, below=of e3] (e4) {4. Contrastar con otros medios};
\node[etapa, below=of e4] (e5) {5. Verificar la imagen};
\node[final, below left=0.7cm and 1.5cm of e5, align=center] (si){Todo correcto :\\COMPARTIR};
\node[finalno, below right=0.7cm and 1.5cm of e5, align=center] (no){Una duda :\\NO COMPARTIR};
\draw[flecha] (q) -- (e1);
\draw[flecha] (e1) -- (e2);
\draw[flecha] (e2) -- (e3);
\draw[flecha] (e3) -- (e4);
\draw[flecha] (e4) -- (e5);
\draw[flecha] (e5) -- (si) node[midway, left, font=\scriptsize]{sí};
\draw[flecha] (e5) -- (no) node[midway, right, font=\scriptsize]{no};
\end{tikzpicture}
\legende{Organigramme de vérification : les 5 étapes avant de partager une nouvelle.}
\end{popfigure}

\courssub{Exemple commenté : anatomía de un bulo}

Analysons pas à pas notre message WhatsApp du début, en appliquant les 5 V :

\begin{enumerate}
\item \textbf{Ver el titular} : \esp{¡URGENTE!} en majuscules, points d'exclamation multiples → titre exagéré, conçu pour paniquer. Signal d'alerte numéro 1.
\item \textbf{Verificar la fuente} : les seules « sources » sont \esp{un médico famoso} sans nom et \esp{mi primo} : aucune source identifiable. Signal d'alerte numéro 2. De plus, l'argument \esp{los periodistas no dicen nada porque tienen miedo} est classique : il vise à vous couper des vrais médias.
\item \textbf{Ver la fecha} : aucune date précise, seulement \esp{a partir del lunes} — quel lundi ? Signal d'alerte numéro 3.
\item \textbf{Verificar con otros medios} : une fermeture de tous les marchés du pays pendant trois mois serait une nouvelle énorme ; si aucun journal, aucune radio, aucun site officiel n'en parle, c'est un \esp{bulo}. Signal d'alerte numéro 4.
\item \textbf{Verificar la imagen} : si une photo de « rayons vides » accompagne le message, une recherche inversée montrera probablement qu'elle vient d'un autre pays ou d'une autre année. Signal d'alerte numéro 5.
\end{enumerate}

Conclusion : cinq signaux sur cinq. La bonne réaction est de ne pas transférer le message et, mieux, de le \esp{desmentir} (démentir) auprès de ceux qui l'ont reçu : \esp{Es un bulo, no lo compartas.} « C'est un canular, ne le partage pas. »

\courssub{Le vocabulaire de la vérification}

\begin{center}
\begin{tabularx}{\linewidth}{|X|X|}\hline
\rowcolor{popblueL} Español & Français \\ \hline
\esp{verificar} & vérifier \\ \hline
\esp{comprobar (ue)} & vérifier, contrôler \\ \hline
\esp{desmentir (ie, i)} & démentir \\ \hline
\esp{una fuente fiable} & une source fiable \\ \hline
\esp{contrastar la información} & recouper l'information \\ \hline
\esp{un \cle{hecho}} & un fait (vérifiable) \\ \hline
\esp{una \cle{opinión}} & une opinion (jugement personnel) \\ \hline
\esp{la manipulación} & la manipulation \\ \hline
\esp{la búsqueda inversa de imágenes} & la recherche inversée d'images \\ \hline
\esp{desmentir un bulo} & démentir un canular \\ \hline
\end{tabularx}
\end{center}

\begin{attention}[Faux amis et pièges de traduction]
\begin{itemize}
\item \esp{verificar} et \esp{comprobar} se traduisent tous deux par « vérifier » ; \esp{comprobar} insiste sur le contrôle concret (\esp{comprobé la fecha} « j'ai vérifié la date »), \esp{verificar} sur l'authenticité (\esp{verificar una noticia}).
\item Attention à l'orthographe : \esp{desmentir} s'écrit avec \textbf{s} (pas de « c » : *\esp{descmentir} n'existe pas). Sa conjugaison (changement de radical \textbf{e $\rightarrow$ ie} au présent, \textbf{e $\rightarrow$ i} au prétérit) : \esp{desmiento, desmientes, desmiente, desmentimos, desmentís, desmienten} ; prétérit : \esp{desmintió, desmintieron}.
\item Ne traduisez pas « une fausse nouvelle » par *\esp{una noticia falsa} dans le sens technique : dites \esp{un bulo} ou \esp{una fake news}.
\end{itemize}
\end{attention}

\courssub{Complément pour le Bac : distinguer \cle{hecho} y \cle{opinión}}

Au commentaire de texte, l'examinateur attend que vous distinguiez ce qui est \textbf{fait} et ce qui est \textbf{opinion}. C'est aussi le cœur de la lutte contre la désinformation.

\begin{propriete}[Hecho vs opinión]
\begin{itemize}
\item Un \esp{hecho} est un fait \textbf{vérifiable} : on peut le prouver, le dater, le mesurer. \esp{El mercado cerró el lunes.} « Le marché a fermé lundi. »
\item Una \esp{opinión} est un jugement \textbf{personnel} : elle exprime ce que quelqu'un pense ou ressent. \esp{El mercado es demasiado ruidoso.} « Le marché est trop bruyant. »
\end{itemize}
Indices linguistiques d'une opinion : les verbes \esp{creer que, pensar que, opinar que, me parece que} (suivis de l'indicatif pour \esp{creer/pienser/opinar} affirmatif), les adjectifs évaluatifs (\esp{terrible, maravilloso, exagerado}), et les adverbes comme \esp{desgraciadamente, por suerte}. Le fait, lui, utilise des verbes d'action, des dates et des chiffres vérifiables.
\end{propriete}

\begin{exemplebox}[Entraînement rapide]
\esp{El río Wouri atraviesa Douala.} → \esp{hecho} (vérifiable sur une carte).\\
\esp{Douala es la ciudad más bonita de África.} → \esp{opinión} (jugement personnel).\\
\esp{Creo que las redes sociales son peligrosas.} → \esp{opinión} (signalée par \esp{creo que} + indicatif).\\
\esp{El artículo fue publicado el 3 de mayo.} → \esp{hecho} (daté, vérifiable).
\end{exemplebox}

\begin{exoresolu}[¿Hecho u opinión? ¿Bulo o noticia?]
\textbf{Énoncé.} Pour chaque phrase, dites s'il s'agit d'un \esp{hecho} ou d'une \esp{opinión}, puis indiquez si le message suivant est probablement un \esp{bulo} : \esp{« ¡ESCÁNDALO! Todos los exámenes del BAC están anulados. Lo sé por un amigo. Reenvía antes de que borren el mensaje. »}
\tcblower
\textbf{Solution détaillée.}
\begin{enumerate}
\item \esp{El examen de español dura dos horas.} → \esp{hecho} : durée vérifiable dans le programme officiel.
\item \esp{El examen es demasiado difícil.} → \esp{opinión} : adjectif évaluatif, jugement personnel.
\item \esp{Según el ministerio, los resultados saldrán en julio.} → \esp{hecho} rapporté avec source citée (\esp{según} + source identifiable).
\item Le message sur l'annulation du BAC est très probablement un \esp{bulo} : titre sensationnaliste (\esp{¡ESCÁNDALO!}), source anonyme (\esp{un amigo}), absence de date officielle, pression à partager (\esp{antes de que borren el mensaje} — remarquez le subjonctif après \esp{antes de que}). La vérification correcte : consulter le site officiel du ministère et les grands médias, puis \esp{no compartir}.
\end{enumerate}
\end{exoresolu}

\begin{savaistu}[Los cazadores de bulos]
En Espagne, des plateformes entières sont spécialisées dans la vérification des fausses informations en langue espagnole : \textbf{Maldita.es} et \textbf{Newtral}, par exemple, emploient des journalistes-\esp{verificadores} qui démontent chaque jour les \esp{bulos} qui circulent sur WhatsApp et les réseaux sociaux. Leur métier s'appelle la \esp{verificación de hechos} (fact-checking). Au Cameroun aussi, des initiatives de vérification traquent les rumeurs, notamment en période électorale : savoir vérifier une information est devenu une compétence démocratique dans tout le monde hispanophone et au-delà.
\end{savaistu}

\begin{aretenir}[L'essentiel de la section]
\begin{itemize}
\item Une fake news (\esp{\cle{bulo}}, \esp{\cle{desinformación}}) est une information fausse, créée pour tromper, qui se diffuse très vite.
\item Cinq signaux d'alerte : \esp{titular exagerado, falta de fuente, fecha antigua, foto trucada, errores de ortografía}.
\item La règle des 5 V : \esp{Ver el titular, Verificar la fuente, Ver la fecha, Verificar con otros medios, Verificar la imagen}.
\item En cas de doute : \esp{No compartas una noticia sin verificarla.}
\item Au Bac, distinguez toujours \esp{\cle{hecho}} (fait vérifiable) et \esp{\cle{opinión}} (jugement personnel).
\end{itemize}
\end{aretenir}

\begin{exercice}[À vous de jouer]
Votre grand-mère reçoit sur WhatsApp une photo d'inondations accompagnée du message : \esp{« ¡AHORA MISMO en Douala! ¡Comparte! »}. Rédigez en espagnol, en 5 à 8 lignes, la réponse que vous lui envoyez pour lui expliquer comment vérifier ce message avant de le partager. Utilisez au moins quatre mots du vocabulaire de la vérification (\esp{verificar, fuente fiable, fecha, bulo, comprobar}).
\end{exercice}

\begin{corrige}[Exercice « À vous de jouer »]
Modèle de réponse possible : \esp{Abuela, ten cuidado con ese mensaje. Antes de compartirlo, hay que verificar la información. Primero, mira la fecha: la foto puede ser antigua. Segundo, comprueba la fuente: ¿quién ha enviado la foto? ¿Es una fuente fiable? Tercero, busca la noticia en otros medios: si es verdad, los periódicos hablarán de las inundaciones. Por último, haz una búsqueda inversa de la imagen: quizá sea de otro país. Si no puedes comprobar nada, es probablemente un bulo. No compartas una noticia sin verificarla.} « Grand-mère, fais attention à ce message. Avant de le partager, il faut vérifier l'information. D'abord, regarde la date : la photo est peut-être ancienne. Ensuite, vérifie la source : qui a envoyé la photo ? Est-ce une source fiable ? Troisièmement, cherche la nouvelle dans d'autres médias : si c'est vrai, les journaux parleront des inondations. Enfin, fais une recherche inversée de l'image : elle vient peut-être d'un autre pays. Si tu ne peux rien vérifier, c'est probablement un canular. Ne partage pas une information sans l'avoir vérifiée. »
\end{corrige}

\coursec{Méthode du commentaire et commentaire modèle}

Dans la section précédente, nous avons appris à détecter les \esp{bulos} et la \esp{desinformación} grâce à une méthode de vérification rigoureuse. Mais au baccalauréat, il ne suffit pas de comprendre un texte : il faut savoir en rendre compte par écrit, de manière organisée et personnelle. C'est l'objet du \cle{comentario de texto}, l'exercice roi de l'épreuve d'espagnol en Terminale A. Cette dernière section vous donne la méthode complète, un modèle rédigé et les critères d'évaluation officiels.

\courssub{Rappel de la méthode : les trois étapes du commentaire}

\begin{definition}[Le comentario de texto]
Le \cle{comentario de texto} est une rédaction en espagnol qui présente, analyse et apprécie un texte. Il comporte toujours trois parties : une \cle{introducción} (présentation du texte + problématique), un \cle{desarrollo} (développement en 2 ou 3 parties) et une \cle{conclusión} (bilan + ouverture). Il ne s'agit ni d'un résumé ni d'une paraphrase : il faut \emph{analyser}, c'est-à-dire dégager les idées de l'auteur, les illustrer par de courtes citations et donner son avis argumenté.
\end{definition}

\begin{methode}[Rédiger un commentaire en cinq étapes]
\begin{enumerate}
\item \textbf{Lire deux fois le texte.} À la première lecture, repérez le thème général ; à la deuxième, soulignez les mots-clés (\esp{libertad de expresión, difamación, censura, bulo}...).
\item \textbf{Formuler la problématique} sous forme de question : \esp{¿Dónde acaba la libertad de expresión y empieza el delito?} « Où s'arrête la liberté d'expression et où commence le délit ? »
\item \textbf{Construire le plan} en 2 ou 3 parties qui répondent à la problématique. Annoncez-le dans l'introduction.
\item \textbf{Rédiger} en citant le texte entre guillemets (3 à 5 mots maximum) et en reliant les idées avec des connecteurs.
\item \textbf{Conclure} par un bilan (\esp{En conclusión...}) et une ouverture (une question, un lien avec l'actualité camerounaise ou mondiale).
\end{enumerate}
\end{methode}

\begin{methode}[Phrases toutes faites pour le commentaire]
Voici des \cle{fórmulas} prêtes à l'emploi, à mémoriser pour l'examen :
\begin{itemize}
\item Présenter le texte : \esp{El texto trata de la libertad de expresión y sus límites.} « Le texte traite de la liberté d'expression et de ses limites. »
\item Présenter la thèse : \esp{El autor defiende la idea de que la libertad exige responsabilidad.} « L'auteur défend l'idée que la liberté exige de la responsabilité. »
\item Introduire une citation : \esp{Como ejemplo, cita los delitos de prensa.} « Comme exemple, il cite les délits de presse. »
\item Donner son avis : \esp{En mi opinión, el autor tiene razón porque...} « À mon avis, l'auteur a raison parce que... »
\item Conclure : \esp{En conclusión, el texto nos invita a reflexionar.} « En conclusion, le texte nous invite à réfléchir. »
\end{itemize}
\end{methode}

\begin{popfigure}
\begin{center}
\begin{tikzpicture}[font=\small]
\fill[popblueL] (-4.5,3.4) -- (4.5,3.4) -- (3.6,2.2) -- (-3.6,2.2) -- cycle;
\node[align=center] at (0,2.8) {\textbf{Introducción} : presentación + problemática\\ \emph{(environ 15 minutes)}};
\fill[poporangeL] (-3.6,2.0) -- (3.6,2.0) -- (2.7,0.8) -- (-2.7,0.8) -- cycle;
\node[align=center] at (0,1.4) {\textbf{Desarrollo I} : la libertad, derecho fundamental};
\fill[popgreenL] (-2.7,0.6) -- (2.7,0.6) -- (1.8,-0.6) -- (-1.8,-0.6) -- cycle;
\node[align=center] at (0,0) {\textbf{Desarrollo II} : los límites : difamación, calumnia};
\fill[poppurpleL] (-1.8,-0.8) -- (1.8,-0.8) -- (0.9,-2.0) -- (-0.9,-2.0) -- cycle;
\node[align=center] at (0,-1.4) {\textbf{Desarrollo III} : la desinformación};
\fill[popgoldL] (-0.9,-2.2) -- (0.9,-2.2) -- (0,-3.2) -- cycle;
\node[align=center] at (0,-2.6) {\textbf{Conclusión}\\ \emph{(10 min)}};
\node[align=center, text=popmuted] at (5.6,0.2) {Le développement\\ occupe environ\\ 30 minutes :\\ c'est le cœur\\ du devoir.};
\end{tikzpicture}
\end{center}
\legende{La pyramide du commentaire : de l'introduction (large) à la conclusion (brève), avec les durées indicatives pour une épreuve de 2 heures.}
\end{popfigure}

\courssub{Problématique et plan modèles}

Pour notre texte \esp{« Libertad de expresión: ¿un derecho sin límites? »}, voici la problématique et le plan que nous allons suivre :

\begin{aretenir}[Problématique et plan modèles]
\textbf{Problématique :} \esp{¿Dónde acaba la libertad de expresión y empieza el delito?}\par
\textbf{Plan en trois parties :}
\begin{enumerate}
\item \esp{La libertad de expresión, un derecho fundamental.} « La liberté d'expression, un droit fondamental. »
\item \esp{Sus límites: los delitos de prensa (difamación, calumnia, injuria).} « Ses limites : les délits de presse. »
\item \esp{La desinformación, un reto actual.} « La désinformation, un défi actuel. »
\end{enumerate}
\end{aretenir}

\courssub{Commentaire modèle rédigé}

Lisez d'abord cet extrait du commentaire modèle : observez comment les connecteurs (soulignés) articulent les idées et comment les citations sont brèves.

\begin{textebox}[Extrait du commentaire modèle]
\underline{En primer lugar}, el autor recuerda que la libertad de expresión es un derecho fundamental, « base de toda democracia », que permite a los ciudadanos informarse y criticar. \underline{Sin embargo}, subraya que este derecho tiene límites: la difamación y la calumnia no son opiniones, sino delitos de prensa que destruyen la reputación de las personas. \underline{Además}, el texto denuncia la desinformación: los bulos se propagan más rápido que la verdad, \underline{sobre todo} en las redes sociales. \underline{En definitiva}, el texto nos invita a verificar antes de compartir.
\end{textebox}

\textbf{Traduction :} « Premièrement, l'auteur rappelle que la liberté d'expression est un droit fondamental, "base de toute démocratie", qui permet aux citoyens de s'informer et de critiquer. Cependant, il souligne que ce droit a des limites : la diffamation et la calomnie ne sont pas des opinions, mais des délits de presse qui détruisent la réputation des personnes. De plus, le texte dénonce la désinformation : les intox se propagent plus vite que la vérité, surtout sur les réseaux sociaux. En définitive, le texte nous invite à vérifier avant de partager. »

Voici maintenant le commentaire complet (environ 18 lignes), tel qu'un excellent élève de Terminale A pourrait le rédiger :

\begin{textebox}[Comentario modelo completo]
El texto « Libertad de expresión: ¿un derecho sin límites? » trata de un tema muy actual: los límites de lo que podemos decir o publicar. El autor defiende la idea de que la libertad exige responsabilidad. Podemos preguntarnos: ¿dónde acaba la libertad de expresión y empieza el delito?

En primer lugar, el autor recuerda que la libertad de expresión es un derecho fundamental, « base de toda democracia ». Gracias a ella, los periodistas informan y los ciudadanos pueden criticar a los poderosos. Sin esta libertad, no hay democracia posible.

Sin embargo, este derecho no es absoluto. El texto explica que la difamación y la calumnia son delitos de prensa: acusar falsamente a alguien no es una opinión, es un crimen contra su honor. Como ejemplo, cita el caso de un rumor falso difundido contra un comerciante, que destruyó su negocio.

Además, el autor denuncia la desinformación: los bulos circulan por WhatsApp y Facebook más rápido que la verdad. Según él, « la libertad sin responsabilidad se convierte en manipulación ».

En conclusión, el texto nos invita a verificar antes de compartir. En mi opinión, esta reflexión es muy importante en el Camerún, donde las noticias falsas pueden dividir a las comunidades. ¿Seremos capaces de educar a los jóvenes en el uso crítico de los medios?
\end{textebox}

\textbf{Glossaire :} \esp{acabar} = s'arrêter, finir ; \esp{los poderosos} = les puissants ; \esp{el honor} = l'honneur, la réputation ; \esp{un rumor} = une rumeur ; \esp{un comerciante} = un commerçant ; \esp{compartir} = partager ; \esp{seremos capaces} = nous serons capables.

\begin{exemplebox}[Exemples traduits : phrases d'analyse]
\begin{itemize}
\item \esp{Según el autor, la libertad sin responsabilidad se convierte en manipulación.} « Selon l'auteur, la liberté sans responsabilité se transforme en manipulation. »
\item \esp{El texto denuncia la censura de los medios independientes.} « Le texte dénonce la censure des médias indépendants. »
\item \esp{Como ejemplo, cita un bulo que circuló durante las elecciones.} « Comme exemple, il cite une intox qui a circulé pendant les élections. »
\item \esp{El autor subraya que verificar es un deber del ciudadano.} « L'auteur souligne que vérifier est un devoir du citoyen. »
\end{itemize}
\end{exemplebox}

\begin{center}
\begin{tabularx}{\linewidth}{|X|X|X|}
\hline
\rowcolor{popblueL} \textbf{Fonction} & \textbf{Connecteurs espagnols} & \textbf{Équivalents français} \\
\hline
Ordonner & \esp{en primer lugar, además, por último} & premièrement, de plus, enfin \\
\hline
Opposer & \esp{sin embargo, no obstante, en cambio} & cependant, néanmoins, en revanche \\
\hline
Expliquer & \esp{es decir, porque, ya que} & c'est-à-dire, parce que, puisque \\
\hline
Donner un avis & \esp{en mi opinión, a mi juicio, creo que} & à mon avis, selon moi, je pense que \\
\hline
Conclure & \esp{en conclusión, en definitiva, por tanto} & en conclusion, en définitive, donc \\
\hline
\end{tabularx}
\end{center}

\begin{attention}[Piège du Bac : la citation trop longue]
Au baccalauréat, \textbf{ne recopiez jamais de longs passages} du texte : le correcteur le sanctionne comme de la paraphrase. Citez \cle{3 à 5 mots} entre guillemets, puis commentez-les avec vos propres mots. Autre piège : le calque du français. On dit \esp{según el autor} (« selon l'auteur »), jamais \esp{*según de el autor} ; et \esp{en mi opinión}, pas \esp{*en mi opinión mía}. Enfin, n'oubliez pas les accents : \esp{opinión, información, conclusión}.
\end{attention}

\begin{exoresolu}[Construire une introduction de commentaire]
\textbf{Énoncé :} À partir du texte \esp{« Libertad de expresión: ¿un derecho sin límites? »}, rédigez une introduction de 4 lignes (présentation + problématique + annonce du plan).
\tcblower
\textbf{Solution détaillée :}\par
\esp{El texto « Libertad de expresión: ¿un derecho sin límites? » trata de la libertad de expresión y de sus límites en la sociedad actual. El autor defiende la idea de que este derecho fundamental exige responsabilidad. Podemos preguntarnos: ¿dónde acaba la libertad y empieza el delito? Para responder, veremos primero la importancia de este derecho, después sus límites — los delitos de prensa — y por último el reto de la desinformación.}\par
Traduction : « Le texte traite de la liberté d'expression et de ses limites dans la société actuelle. L'auteur défend l'idée que ce droit fondamental exige de la responsabilité. Nous pouvons nous demander : où s'arrête la liberté et où commence le délit ? Pour répondre, nous verrons d'abord l'importance de ce droit, ensuite ses limites — les délits de presse — et enfin le défi de la désinformation. »\par
\textbf{Méthode :} phrase 1 = présentation (\esp{El texto trata de...}) ; phrase 2 = thèse (\esp{El autor defiende la idea de que...}) ; phrase 3 = problématique ; phrase 4 = annonce du plan (\esp{primero... después... por último...}).
\end{exoresolu}

\courssub{Production guidée : luchar contra los bulos en el Camerún}

\begin{experience}[Production écrite guidée]
Rédigez en espagnol un paragraphe de \textbf{10 à 15 lignes} répondant à la question : \esp{¿Cómo podemos luchar contra los bulos en el Camerún?} « Comment lutter contre les intox au Cameroun ? »\par
\textbf{Contraintes :} utilisez obligatoirement ces quatre mots du lexique : \esp{bulo}, \esp{desinformación}, \esp{verificar}, \esp{manipulación} ; employez au moins deux connecteurs (\esp{en primer lugar, sin embargo, además, en conclusión}) ; donnez un exemple camerounais concret (réseaux sociaux, élections, écoles de Yaoundé ou Douala).\par
\textbf{Aide au démarrage :} \esp{En primer lugar, para luchar contra los bulos en el Camerún, los ciudadanos deben verificar la información antes de compartirla...}
\end{experience}

\courssub{Critères d'évaluation type Bac}

\begin{aretenir}[Comment le correcteur note votre commentaire]
\begin{itemize}
\item \textbf{Respect du plan} (introduction, développement en 2-3 parties, conclusion) : le squelette du devoir doit être visible.
\item \textbf{Citations} : au moins deux citations courtes (3-5 mots) entre guillemets, commentées et non recopiées.
\item \textbf{Lexique} : emploi correct du vocabulaire du thème (\esp{difamación, calumnia, censura, bulo, desinformación, publicidad, eslogan, manipulación}).
\item \textbf{Correction grammaticale} : accords, conjugaisons, accents, ponctuation espagnole (¿...? ¡...!), connecteurs variés.
\item \textbf{Prise de position} : un avis personnel argumenté (\esp{en mi opinión...}) valorise toujours la copie.
\end{itemize}
\end{aretenir}

\begin{exercice}[Pour s'entraîner]
\begin{enumerate}
\item Remettez dans l'ordre les étapes du commentaire : \esp{conclusión / introducción / desarrollo / lectura del texto / problemática}.
\item Traduisez en espagnol : « Selon l'auteur, la désinformation est un danger pour la démocratie ; en conclusion, il faut vérifier avant de partager. »
\item Trouvez la citation la plus courte possible du commentaire modèle pour illustrer la phrase : \esp{La libertad exige responsabilidad.}
\end{enumerate}
\end{exercice}

\begin{corrige}[Exercice]
\begin{enumerate}
\item \esp{lectura del texto → problemática → introducción → desarrollo → conclusión.}
\item \esp{Según el autor, la desinformación es un peligro para la democracia; en conclusión, hay que verificar antes de compartir.} (Attention à l'accent de \esp{desinformación} et à la construction \esp{hay que} + infinitif.)
\item \esp{« la libertad sin responsabilidad se convierte en manipulación »} — on peut même réduire à \esp{« se convierte en manipulación »} (4 mots).
\end{enumerate}
\end{corrige}

\begin{savaistu}[Le commentaire, une tradition hispanique]
Le \esp{comentario de texto} est l'exercice central du bachillerato espagnol et de la \esp{selectividad} (l'équivalent espagnol du baccalauréat) : les élèves d'Espagne commentent chaque année des textes de presse ou de littérature. Maîtriser cette méthode en Terminale A au Cameroun, c'est donc acquérir exactement la même compétence qu'un lycéen de Madrid, de Mexico ou de Malabo — un atout précieux pour des études supérieures dans tout l'espace hispanophone.
\end{savaistu}

\newpage

\coursec{Activité d'intégration}

\coursec{Lección E26 : Lectura y comentario : libertad de expresión y delitos de prensa ; análisis de publicidad y fake news}

\courssub{Objectifs de l'activité}

À l'issue de cette activité d'intégration, l'élève sera capable de :
\begin{itemize}
\item mobiliser le \cle{lexique des médias} étudié dans la leçon (\esp{bulo}, \esp{fuente}, \esp{verificar}, \esp{manipulación}, \esp{difamación}, \esp{publicidad encubierta}, \esp{desinformación}, \esp{calumnia}, \esp{censura}, \esp{eslogan}) dans une production authentique ;
\item appliquer la \cle{méthode des 5 V} (\esp{ver la fuente}, \esp{verificar la fecha}, \esp{valorar la intención}, \esp{contrastar la información}, \esp{dar un veredicto}) pour analyser des messages douteux ;
\item distinguer une \cle{information fiable}, une \cle{publicité déguisée} et un \cle{bulo} (fausse information virale) ;
\item rédiger en espagnol un court \cle{article de fact-checking} structuré (titre, présentation du cas, analyse, verdict) ;
\item présenter oralement un verdict argumenté en deux minutes, en soignant la prononciation et l'intonation.
\end{itemize}

\begin{definition}[Lexique thématique : liberté d'expression et médias]
\begin{itemize}
\item \esp{difamación} : action de nuire à la réputation de quelqu'un par des propos diffamatoires (diffamation).
\item \esp{calumnia} : accusation fausse et intentionnelle portée contre quelqu'un (calomnie).
\item \esp{censura} : suppression ou contrôle préventif de l'information par une autorité (censure).
\item \esp{bulo} : fausse nouvelle diffusée intentionnellement pour tromper (canular, \emph{fake news}).
\item \esp{desinformación} : information fausse ou trompeuse diffusée délibérément (désinformation).
\item \esp{publicidad} : communication commerciale pour promouvoir un produit ou service (publicité).
\item \esp{eslogan} : phrase courte et accrocheuse utilisée en publicité (slogan).
\item \esp{manipulación} : action de manipuler l'information ou l'opinion (manipulation).
\item \esp{publicidad encubierta} : publicité déguisée en information (publireportage, \emph{native advertising}).
\item \esp{fuente} : origine d'une information (source).
\item \esp{verificar} : vérifier, confirmer la véracité d'un fait.
\end{itemize}
\end{definition}

\courssub{Situation}

Tu es élève en Terminale A au lycée de Nkol-Eton, à Yaoundé. Ton professeur d'espagnol a créé avec ta classe une \emph{agence de vérification} appelée \esp{« Verificado »}. Chaque semaine, les élèves-détectives analysent des messages qui circulent sur WhatsApp dans les groupes de la communauté hispanophone du quartier (élèves, parents, commerçants du marché Mokolo, étudiants de Guinée équatoriale installés à Yaoundé).

Ce matin, trois messages ont été transmis à ton agence. À ton groupe de quatre détectives de déterminer lequel est une information fiable, lequel est une publicité déguisée et lequel est un \esp{bulo}.

\begin{textebox}[Mensaje 1 — reenviado muchas veces]
¡¡URGENTE!! ¡COMPARTE CON TODOS TUS CONTACTOS!\\
El Ministerio de Salud acaba de confirmar que el agua del grifo de Yaoundé contiene una bacteria mortal. NO bebas agua del grifo durante 15 días. Un médico del Hospital Central lo ha confirmado. ¡¡Salva vidas, reenvía este mensaje YA!! 100\% seguro.
\end{textebox}

\begin{textebox}[Mensaje 2 — publicado en un grupo de padres de alumnos]
¿Sabías que tu hijo puede hablar español como un madrileño en solo 3 meses? La Academia Estrella, en el barrio Bastos, garantiza resultados increíbles: el 98\% de nuestros alumnos aprueba el BAC con mención. Miles de familias ya confían en nosotros. Inscripciones abiertas: ¡plazas limitadas! Llama ya al 6XX XX XX XX. ¡No dejes pasar esta oportunidad única!
\end{textebox}

\begin{textebox}[Mensaje 3 — artículo de un periódico en línea conocido]
YAOUNDÉ — El Ministerio de Educación Secundaria anunció ayer, en un comunicado firmado por el ministro, que el curso escolar terminará el 13 de junio. El comunicado, publicado en el sitio web oficial del ministerio y confirmado por la radio nacional, precisa que los exámenes oficiales comenzarán el 16 de junio. Los directores de los establecimientos escolares recibirán las instrucciones detalladas esta semana.
\end{textebox}

\begin{aretenir}[Glossaire des messages (mots difficiles)]
\begin{center}
\begin{tabularx}{\linewidth}{|X|X|}
\hline
\rowcolor{popblueL} Español & Français \\
\hline
\esp{grifo} & robinet \\
\esp{bacteria mortal} & bactérie mortelle \\
\esp{reenvía} & transfère, renvoie (impératif) \\
\esp{madrileño} & madrilène (habitant de Madrid) \\
\esp{garantiza} & garantit \\
\esp{mención} & mention (au bac) \\
\esp{plazas limitadas} & places limitées \\
\esp{comunicado} & communiqué \\
\esp{precisa} & précise \\
\hline
\end{tabularx}
\end{center}
\end{aretenir}

\courssub{Consignes}

\textbf{Tarea 1 — Primera lectura (comprensión).} Lee los tres mensajes con tu grupo. Para cada uno, subraya dos o tres palabras o expresiones que te llaman la atención (mayúsculas, signos de exclamación, cifras, promesas). Anota tu primera impresión: ¿te parece fiable, sospechoso o publicitario? Justifica tu impresión en una frase en español.

\textbf{Tarea 2 — Aplicación de la méthode des 5 V.} Remplis la fiche d'analyse suivante pour chaque message (en espagnol) :

\begin{center}
\begin{tabularx}{\linewidth}{|l|X|X|X|}
\hline
\rowcolor{popblueL} Criterio & Mensaje 1 & Mensaje 2 & Mensaje 3 \\
\hline
Fuente (¿quién habla?) & & & \\
\hline
Fecha y lugar & & & \\
\hline
Intención (informar, vender, asustar) & & & \\
\hline
Verificación (¿se puede contrastar?) & & & \\
\hline
Veredicto (fiable / publicidad / bulo) & & & \\
\hline
\end{tabularx}
\end{center}

\textbf{Tarea 3 — Lexique en action.} Dans ta fiche, emploie obligatoirement les termes suivants au moins une fois : \esp{bulo}, \esp{fuente}, \esp{verificar}, \esp{manipulación}, \esp{publicidad encubierta}, \esp{eslogan}, \esp{desinformación}, \esp{difamación}, \esp{calumnia}, \esp{censura}. Souligne-les.

\textbf{Tarea 4 — Rédaction de l'article « Verificado ».} Choisis avec ton groupe le message qui vous semble le plus dangereux. Rédigez ensemble un article de fact-checking de 10 à 15 lignes intitulé \esp{« Verificado »}, avec la structure suivante :
\begin{enumerate}
\item un \cle{titre} qui annonce le verdict (\esp{Falso: ...} ou \esp{Verdadero: ...}) ;
\item la présentation du message analysé (où circule-t-il, que dit-il) ;
\item l'analyse : deux ou trois arguments fondés sur la méthode des 5 V ;
\item le verdict final et un conseil aux lecteurs.
\end{enumerate}

\textbf{Tarea 5 — Présentation orale.} Chaque groupe présente son verdict à la classe en deux minutes. Les autres groupes posent une question en espagnol (\esp{¿Por qué pensáis que es una manipulación?}). La classe vote ensuite pour le meilleur article de l'agence.

\courssub{Ressources à mobiliser}

\begin{aretenir}[La méthode des 5 V]
\begin{enumerate}
\item \esp{Ver la fuente} : qui parle ? Une source anonyme (\esp{« un médico »}, \esp{« un amigo »}) est suspecte ; une source officielle identifiable est plus fiable.
\item \esp{Verificar la fecha y el lugar} : le message donne-t-il une date précise, un lieu, un lien vers un site officiel ?
\item \esp{Valorar la intención} : informer, vendre, faire peur, faire rire ? Les majuscules, les exclamations et l'urgence (\esp{¡URGENTE!}, \esp{¡YA!}) signalent souvent une manipulation.
\item \esp{Contrastar la información} : chercher la même nouvelle dans deux médias sérieux ou sur un site officiel.
\item \esp{Dar un veredicto} : \esp{verdadero}, \esp{falso}, \esp{dudoso} ou \esp{publicidad encubierta}.
\end{enumerate}
\end{aretenir}

\begin{propriete}[Structures utiles pour l'article de fact-checking]
\begin{itemize}
\item Présenter le cas : \esp{El mensaje que analizamos circula por WhatsApp desde hace una semana.} ; \esp{Según el mensaje, ...}
\item Exprimer le doute : \esp{No hay ninguna prueba de que...} (subjonctif après une expression de doute) ; \esp{Dudamos que sea verdad.}
\item Argumenter : \esp{En primer lugar...} ; \esp{Además...} ; \esp{Por último...} ; \esp{Sin embargo...}
\item Conclure : \esp{Por lo tanto, nuestro veredicto es claro: es un bulo.} ; \esp{Recomendamos no compartir este mensaje.}
\end{itemize}
\end{propriete}

\begin{attention}[Pièges de langue]
\begin{itemize}
\item \esp{Falso} signifie « faux », pas « fausse » au sens grammatical : \esp{una noticia falsa} (accord en genre).
\item Ne traduis pas « vérifier » par \esp{*verificar} sans complément correct : \esp{verificar la información}, \esp{comprobar los hechos}.
\item \esp{Publicidad} = publicité ; « publicité » au sens de « notoriété » se dit \esp{publicidad} aussi, mais attention au faux ami \esp{actualidad} (= actualité, pas « actuality »).
\item Après \esp{es posible que}, \esp{no creemos que}, \esp{dudamos que} : subjonctif (\esp{sea}, \esp{contenga}, \esp{haya}).
\item \esp{Calumnia} et \esp{difamación} sont des délits de presse (\esp{delitos de prensa}) dans de nombreux codes pénaux hispanophones ; \esp{censura} désigne le contrôle préventif, pas l'autocensure.
\end{itemize}
\end{attention}

\courssub{Proposition de production et corrigé commenté}

\begin{exoresolu}[Article modèle de l'agence « Verificado »]
Rédige, avec ton groupe, un article de 10 à 15 lignes analysant le message 1 (l'eau de Yaoundé).
\tcblower
\textbf{Production modèle :}

\esp{Falso: el agua del grifo de Yaoundé no contiene ninguna bacteria mortal}

\esp{El mensaje que analizamos circula por WhatsApp desde hace varios días y pide a los lectores que lo reenvíen urgentemente. Según el texto, el agua del grifo de Yaoundé contiene una bacteria mortal confirmada por « un médico del Hospital Central ».}

\esp{Nuestra agencia ha aplicado el método de verificación. En primer lugar, la fuente es anónima: no se da el nombre del médico ni ningún comunicado oficial del Ministerio de Salud. Además, el mensaje utiliza mayúsculas, signos de exclamación y la palabra « URGENTE », técnicas típicas de la manipulación para asustar a la gente. Por último, hemos consultado el sitio web del Ministerio y la radio nacional: no hay ninguna prueba de que exista esta bacteria.}

\esp{Por lo tanto, nuestro veredicto es claro: es un bulo, es decir, una desinformación creada para provocar pánico. Recomendamos no compartir este mensaje y verificar siempre la fuente antes de reenviar una noticia.}

\textbf{Commentaire des choix de langue :}
\begin{itemize}
\item \esp{Falso:} en tête d'article : le verdict est annoncé dès le titre, comme dans les vraies rubriques de fact-checking.
\item \esp{pide a los lectores que lo reenvíen} : subjonctif obligatoire après \esp{pedir que}.
\item \esp{no hay ninguna prueba de que exista} : subjonctif après une négation d'existence — structure de Terminale bien employée.
\item \esp{En primer lugar / Además / Por último / Por lo tanto} : connecteurs logiques qui structurent l'argumentation.
\item Lexique de la leçon intégré naturellement : \esp{fuente}, \esp{manipulación}, \esp{bulo}, \esp{desinformación}, \esp{verificar}.
\item \esp{es decir} : reformulation qui définit le mot technique \esp{bulo} pour le lecteur — bonne pratique journalistique.
\end{itemize}
\end{exoresolu}

\courssub{Bilan de l'activité}

Cette activité a permis de réunir toutes les compétences de la leçon dans une tâche finale authentique. En jouant les détectives de l'information, tu as :
\begin{itemize}
\item \textbf{compris} trois documents de types différents (message viral, publicité, article de presse) et appris à repérer leurs marques formelles ;
\item \textbf{mobilisé le lexique} des médias et de la liberté d'expression dans un contexte de communication réel, celui des groupes WhatsApp camerounais ;
\item \textbf{appliqué une méthode} transférable : les 5 V te serviront pour tous les sujets d'examen portant sur les médias, mais aussi dans ta vie quotidienne de citoyen ;
\item \textbf{produit un texte argumenté} en espagnol, avec connecteurs, subjonctif et verdict explicite — exactement le type d'écrit attendu en expression écrite au BAC ;
\item \textbf{pris la parole} en public pour défendre un verdict, compétence orale essentielle de l'épreuve.
\end{itemize}

\begin{savaistu}[Le fact-checking dans le monde hispanophone]
En Espagne et en Amérique latine, de nombreux médias ont créé des rubriques de vérification inspirées du journalisme anglo-saxon : on y analyse les rumeurs qui circulent sur les réseaux sociaux, surtout pendant les élections ou les crises sanitaires. Le mot espagnol \esp{bulo} désigne une fausse nouvelle répandue volontairement ; il est devenu un mot-clé du débat public hispanophone, au même titre que l'anglais \emph{fake news}. Au Cameroun, où WhatsApp est un canal d'information majeur, savoir vérifier une source est une compétence de citoyenneté autant que de langue.
\end{savaistu}

\newpage

\coursec{Méthodes et exercices résolus}

\coursec{Leçon E26 : Lectura y comentario : libertad de expresión y delitos de prensa ; análisis de publicidad y fake news}

\courssub{Texte support : « Libertad de expresión: ¿un derecho sin límites? »}

\begin{textebox}[Artículo de opinión original — Contexto: España, América Latina y Guinea Ecuatorial]
La libertad de expresión es el pilar de toda democracia, pero su ejercicio conlleva responsabilidades. En la era digital, la línea entre el derecho a informar y el delito de prensa se vuelve difusa. Un periodista que denuncia la corrupción en un ministerio ejerce su derecho; quien acusa sin pruebas a un ciudadano de un delito comete calumnia. La difamación, la injuria y la propagación de bulos —las llamadas fake news— no son opiniones protegidas, sino ataques a la honra y a la convivencia.

Sin embargo, el remedio no puede ser la censura previa. Los gobiernos que silencian a la prensa crítica bajo pretexto de «seguridad nacional» o «lucha contra la desinformación» suelen buscar impunidad. En Guinea Ecuatorial, como en otros países, la ley de prensa debe garantizar el acceso a la información pública y proteger las fuentes, no criminalizar la disidencia. La solución pasa por una legislación clara, jueces independientes y, sobre todo, una ciudadanía educada en el pensamiento crítico: capaz de verificar una fuente, detectar un eslogan manipulador y no compartir un bulo por impulso. La libertad sin verdad es manipulación; la verdad sin libertad es propaganda.
\end{textebox}

\begin{definition}[Glossaire et traduction du texte support]
\textbf{Traduction :} La liberté d'expression est le pilier de toute démocratie, mais son exercice implique des responsabilités. À l'ère numérique, la frontière entre le droit d'informer et le délit de presse devient floue. Un journaliste qui dénonce la corruption dans un ministère exerce son droit ; qui accuse sans preuves un citoyen d'un délit commet un calomnie. La diffamation, l'injure et la propagation de canulars — les fameuses \emph{fake news} — ne sont pas des opinions protégées, mais des atteintes à l'honneur et à la coexistence.

Cependant, le remède ne peut être la censure préalable. Les gouvernements qui font taire la presse critique sous prétexte de « sécurité nationale » ou de « lutte contre la désinformation » cherchent souvent l'impunité. En Guinée équatoriale, comme dans d'autres pays, la loi sur la presse doit garantir l'accès à l'information publique et protéger les sources, non criminaliser la dissidence. La solution passe par une législation claire, des juges indépendants et, surtout, une citoyenneté éduquée à l'esprit critique : capable de vérifier une source, de détecter un slogan manipulateur et de ne pas partager un canular par impulsion. La liberté sans vérité est manipulation ; la vérité sans liberté est propagande.

\textbf{Mots difficiles :} \esp{conllevar} (impliquer, entraîner) ; \esp{difusa} (floue, imprécise) ; \esp{honra} (honneur, réputation) ; \esp{convivencia} (coexistence, vie en société) ; \esp{pretexto} (prétexte) ; \esp{impunidad} (impunité) ; \esp{disidencia} (dissidence) ; \esp{eslogan} (slogan) ; \esp{impulso} (impulsion, élan).
\end{definition}

\courssub{Compréhension du texte support}

\begin{exercice}[Questions sur l'article « Libertad de expresión: ¿un derecho sin límites? »]
Répondez en espagnol par des phrases complètes.
\begin{enumerate}
\item Según el texto, ¿cuál es la diferencia entre un periodista que denuncia la corrupción y quien acusa sin pruebas? (2 pts)
\item ¿Por qué el autor rechaza la censura previa como solución? (2 pts)
\item ¿Qué tres condiciones cita el texto para una buena legislación de prensa? (3 pts)
\item ¿Qué papel asigna el autor a la ciudadanía en la lucha contra la desinformación? (3 pts)
\item \textbf{Expresión personal (80-100 palabras) :} \esp{¿Crees que en tu país (Camerún, Guinea Ecuatorial, España...) la libertad de expresión goza de buena salud? Argumenta tu respuesta citando un ejemplo concreto.} (10 pts)
\end{enumerate}
\end{exercice}

\begin{corrige}[Exercice Compréhension texte support]
\begin{enumerate}
\item \esp{El periodista que denuncia la corrupción ejerce su derecho a informar, mientras que quien acusa sin pruebas comete calumnia, que es un delito de prensa.}
\item \esp{Porque los gobiernos la usan para silenciar a la prensa crítica y buscar impunidad, bajo pretextos como la «seguridad nacional».}
\item \esp{Legislación clara, jueces independientes y ciudadanía educada en el pensamiento crítico.}
\item \esp{La ciudadanía debe verificar las fuentes, detectar la manipulación publicitaria y no compartir bulos por impulso.}
\item \esp{(Exemple de réponse) En mi opinión, en Camerún la libertad de expresión avanza pero sigue frágil. Por un lado, existen periódicos independientes como «Le Messager» que critican al gobierno. Por otro, periodistas han sido detenidos por «noticias falsas» bajo la ley de 2019. Por ejemplo, en 2023, un reportero fue arrestado por denunciar la corrupción en una aduana. Eso muestra que la ley a veces sirve para callar la disidencia. Es necesario que los jueces sean independientes para proteger el derecho a informar.}
\end{enumerate}
\end{corrige}

\courssub{Fiche méthode 1 : Lire et commenter un texte sur la libertad de expresión}

\begin{methode}[Comprendre et commenter un texte argumentatif sur la libertad de expresión]
\begin{enumerate}
\item \textbf{Première lecture globale.} Lis le texte entier sans dictionnaire. Repère : le thème (\esp{¿De qué habla?}), le genre (\esp{artículo de opinión, editorial, crónica}), la thèse de l'auteur (\esp{¿Está a favor o en contra de los límites?}).
\item \textbf{Repérage des connecteurs.} Souligne les marqueurs logiques : \esp{sin embargo}, \esp{no obstante}, \esp{por lo tanto}, \esp{en primer lugar}, \esp{por un lado... por otro}. Ils révèlent la structure de l'argumentation.
\item \textbf{Lexique thématique.} Classe le vocabulaire en champs : el derecho (\esp{libertad de expresión, opinión, prensa, derecho fundamental}), los abusos (\esp{difamación, calumnia, injuria, bulo, desinformación}), la limitación (\esp{censura, ley, sanción, delito de prensa, responsabilidad}).
\item \textbf{Réponse aux questions.} Réponds toujours par une phrase complète en espagnol, en reprenant les termes de la question, puis justifie par une citation courte du texte entre guillemets (\esp{«...»}).
\item \textbf{Commentaire personnel.} Pour la question ouverte (\esp{¿Y tú, qué opinas?}), structure : \textbf{opinión + argumento + ejemplo (Camerún, Guinea Ecuatorial, España) + conclusión}.
\end{enumerate}
\textbf{Structures à retenir :} \esp{El autor defiende que...}, \esp{Según el texto, ...}, \esp{Esto significa que...}, \esp{En mi opinión, la libertad de expresión debe tener límites porque...}, \esp{Estoy de acuerdo / en desacuerdo porque...}
\end{methode}

\begin{exoresolu}[Compréhension de texte — Type Bac]
\textbf{Énoncé.} Lis le texte suivant, puis réponds en espagnol par des phrases complètes.

\emph{« La libertad de expresión es un derecho fundamental, pero no es absoluta. Cuando un periodista publica una calumnia contra un ciudadano, no informa: destruye una reputación. La difamación y los bulos no son opiniones, son delitos. Sin embargo, castigar la crítica legítima sería peor todavía: una sociedad sin prensa libre es una sociedad ciega. »}

\begin{enumerate}
\item ¿Cuál es la tesis del autor? (2 pts)
\item Según el texto, ¿qué diferencia hay entre informar y calumniar? (2 pts)
\item ¿Por qué es peligroso castigar la crítica legítima? (2 pts)
\item ¿Estás de acuerdo con el autor? Justifica tu respuesta en tres frases. (4 pts)
\end{enumerate}
\tcblower
\textbf{Raisonnement.} La thèse se trouve dans la première phrase (\esp{derecho fundamental pero no absoluto}). La question 2 exige de reformuler l'opposition \esp{informar}/\esp{destruir}. La question 4 appelle \esp{opinión + argumento + ejemplo}.

\textbf{Réponses rédigées :}
\begin{enumerate}
\item \esp{El autor defiende que la libertad de expresión es un derecho fundamental, pero que no es absoluta, es decir, que tiene límites.} (Reprise exacte de la thèse avec reformulation : \esp{es decir}.)
\item \esp{Según el texto, informar consiste en comunicar hechos, mientras que calumniar consiste en destruir la reputación de una persona con mentiras.} (Justification : \esp{mientras que} + indicatif pour opposer deux faits réels.)
\item \esp{Es peligroso castigar la crítica legítima porque una sociedad sin prensa libre sería una sociedad ciega, incapaz de denunciar los abusos.} (Attention : \esp{sería}, conditionnel, hypothèse.)
\item \esp{Estoy de acuerdo con el autor. En mi opinión, la libertad de expresión debe proteger la crítica, pero no las mentiras. Por ejemplo, en mi país, un periodista debe poder criticar al gobierno sin miedo, pero no debe inventar acusaciones falsas contra un ciudadano.} (Structure complète : acuerdo + opinión + exemple concret.)
\end{enumerate}
\textbf{Conclusion méthodologique :} Chaque réponse reprend les termes de la question, cite ou reformule le texte, et reste en espagnol correct avec accents (\esp{opinión, crítica, país}).
\end{exoresolu}

\courssub{Lexique thématique : délits de presse, publicité et fake news}

\begin{propriete}[Vocabulaire classé par champs sémantiques]
\begin{center}
\begin{tabularx}{\linewidth}{|X|X|X|}
\hline
\rowcolor{popblueL} \textbf{Français} & \textbf{Español} & \textbf{Contexte / Exemple} \\
\hline
La liberté d'expression & \esp{la libertad de expresión} & \esp{Es un derecho fundamental.} \\
La liberté de la presse & \esp{la libertad de prensa} & \esp{Garantiza el derecho a informar.} \\
L'opinion & \esp{la opinión} & \esp{Todos tienen derecho a su opinión.} \\
\hline
La diffamation & \esp{la difamación} & \esp{Cometer difamación: dañar la fama con hechos.} \\
La calomnie & \esp{la calumnia} & \esp{Acusar falsamente de un delito: calumnia.} \\
L'injure & \esp{la injuria} & \esp{Ofender la dignidad: injuria (insulto).} \\
Le délit de presse & \esp{el delito de prensa} & \esp{La calumnia y la injuria son delitos de prensa.} \\
La censure & \esp{la censura} & \esp{Ejercer la censura: prohibir publicar.} \\
\hline
Le canular / fake news & \esp{el bulo} & \esp{Propagar un bulo en WhatsApp.} \\
La désinformation & \esp{la desinformación} & \esp{Combatir la desinformación.} \\
La manipulation & \esp{la manipulación} & \esp{Manipulación de la opinión pública.} \\
\hline
La publicité & \esp{la publicidad} & \esp{Una campaña de publicidad.} \\
Le slogan & \esp{el eslogan} (pl. \esp{eslóganes}) & \esp{Un eslogan pegadizo.} \\
Le consommateur & \esp{el consumidor} & \esp{El consumidor tiene derechos.} \\
La cible & \esp{el público objetivo} & \esp{El anuncio se dirige al público joven.} \\
\hline
\end{tabularx}
\end{center}
\end{propriete}

\begin{definition}[Distinction clé : Difamación vs Calumnia vs Injuria]
\begin{itemize}
\item \esp{\textbf{Difamación}} : Imputation d'un fait précis qui porte atteinte à l'honneur/réputation (\esp{« El alcalde robó dinero »}). Atteinte à la \esp{fama}.
\item \esp{\textbf{Calumnia}} : Accusation fausse et \textbf{spécifique d'un délit/crime} devant la justice ou publiquement (\esp{« Lo acusé de asesinato »}). Atteinte à l'\esp{honra} + procédure judiciaire.
\item \esp{\textbf{Injuria}} : Expression méprisante, insulte, offense à la dignité, sans imputation de fait précis (\esp{« Es un ladrón, un sinvergüenza »}).
\end{itemize}
\textbf{Aide-mémoire :} \esp{Difamación} = faits (réputation) ; \esp{Calumnia} = délit imaginaire (justice) ; \esp{Injuria} = mots blessants (dignité).
\end{definition}

\courssub{Fiche méthode 2 : Maîtriser le lexique des délits de presse et des médias}

\begin{methode}[Employer sans confusion le vocabulaire juridique et médiatique]
\begin{enumerate}
\item \textbf{Distingue les faux amis et pièges.}
\begin{itemize}
\item \esp{Difamación} (diffamation) $\neq$ \esp{defamación} (anglicisme/erreur).
\item \esp{Calumnia} (calomnie : accusation fausse de crime) $\neq$ \esp{calumnia} (juste « diffamation » en français courant).
\item \esp{Injuria} (injure, insulte) $\neq$ \esp{injurias} (pluriel souvent utilisé pour « insultes »).
\item \esp{Censura} (censure) $\neq$ \esp{cesura} (césure, poésie).
\item \esp{Desinformación} (désinformation) : accent sur \esp{ó}, pas \esp{desinformation}.
\end{itemize}
\item \textbf{Mémorise les collocations (associations figées).}
\begin{itemize}
\item \esp{Cometer un delito de prensa} / \esp{propagar un bulo} / \esp{difundir desinformación}.
\item \esp{Ejercer la censura} / \esp{imponer la censura} / \esp{levantar la censura}.
\item \esp{Lanzar una campaña de publicidad} / \esp{un eslogan publicitario}.
\item \esp{Caer en la trampa} / \esp{verificar la fuente} / \esp{contrastar la información}.
\end{itemize}
\item \textbf{Travaille les familles de mots (dérivation).}
\begin{center}
\begin{tabular}{|l|l|l|}
\hline
\rowcolor{popgreenL} \textbf{Verbe} & \textbf{Nom} & \textbf{Agent / Adjectif} \\
\hline
\esp{censurar} & \esp{la censura} & \esp{el censor / censurable} \\
\esp{manipular} & \esp{la manipulación} & \esp{el manipulador / manipulable} \\
\esp{difamar} & \esp{la difamación} & \esp{el difamador / difamatorio} \\
\esp{calumniar} & \esp{la calumnia} & \esp{el calumniador / calumnioso} \\
\esp{informar} & \esp{la información} & \esp{el informador / informativo} \\
\esp{opinar} & \esp{la opinión} & \esp{el opinador / opinable} \\
\hline
\end{tabular}
\end{center}
\item \textbf{Vérifie l'accentuation (règle des mots en \esp{-ción / -sión} : accent sur le \esp{ó} / \esp{á}).}
\esp{Información, opinión, difamación, manipulación, censura, nación, constitución, televisión.} Exception : \esp{eslogan} (pas d'accent, mot adapté de l'anglais, \esp{eslóganes} au pluriel).
\end{enumerate}
\end{methode}

\begin{exoresolu}[Lexique — Complétion, dérivation et traduction]
\textbf{Énoncé.} A. Complète avec la forme correcte : \esp{bulo, censura, eslogan, difamación, manipulación, calumniar, informar}. (7 pts)

1. El periódico fue condenado por \dots{} después de acusar falsamente al alcalde de robo.
2. El gobierno impuso la \dots{} a tres revistas críticas antes de las elecciones.
3. « Un pueblo, un futuro » es el \dots{} de la campaña electoral.
4. La noticia de que el lago Nyos se había secado era un \dots{} que circuló en WhatsApp.
5. Las redes sociales facilitan la \dots{} de la opinión pública mediante algoritmos.
6. No se debe \dots{} a nadie sin pruebas; es un delito grave.
7. El deber del periodista es \dots{} con veracidad, no inventar.

B. Traduis en espagnol : « La désinformation est un danger pour la démocratie. » (2 pts)
C. Donne le nom de l'agent (celui qui fait l'action) pour : \esp{censurar, manipular, difamar}. (3 pts)
\tcblower
\textbf{Solution A.}
\begin{enumerate}
\item \esp{difamación} (nom) — accusation fausse portant sur un fait précis (vol) atteignant la réputation.
\item \esp{censura} (nom) — interdiction officielle de publication.
\item \esp{eslogan} (nom, sg, sans accent) — formule de campagne. Pluriel : \esp{eslóganes}.
\item \esp{bulo} (nom) — fausse nouvelle, canular (anglicisme \esp{fake news}).
\item \esp{manipulación} (nom en \esp{-ción}, accent sur \esp{ó}) — action d'influencer.
\item \esp{calumniar} (infinitif) — accuser faussement d'un crime/délit.
\item \esp{informar} (infinitif) — communiquer des faits vérifiés.
\end{enumerate}
\textbf{Solution B.} \esp{La desinformación es un peligro para la democracia.} (Faux ami : \esp{peligro} = danger ; \esp{desinformación} avec accent).
\textbf{Solution C.} \esp{El censor, el manipulador, el difamador.}


\end{exoresolu}

\courssub{Fiche méthode 3 : Analyser une publicité}

\begin{methode}[Analizar un anuncio publicitario : étapes et métalangage]
\begin{enumerate}
\item \textbf{Identifie le produit et la cible.} \esp{¿Qué se vende? ¿A quién se dirige?} (jóvenes, familias, deportistas, profesionales).
\item \textbf{Décode le message verbal (texte/son).} Repère : \esp{el eslogan} (formule courte, mémorable), les \textbf{impératifs} (\esp{¡Compra!, ¡Prueba!, ¡Llama ya!}), les \textbf{superlatifs} (\esp{el mejor, el más rápido, la número 1}), les \textbf{questions rhétoriques} (\esp{¿Quieres triunfar?}).
\item \textbf{Décode le message visuel / sonore.} Couleurs, musique, visages, décors : \esp{¿Qué sentimiento busca provocar?} (felicidad, prestigio, miedo, pertenencia, urgencia).
\item \textbf{Repère les procédés de persuasion (argumentation).}
\begin{itemize}
\item \esp{Testimonio / Endoso} : célébrité ou « utilisateur satisfait ».
\item \esp{Argumento de autoridad / Cifras} : \esp{9 de cada 10 dentistas recomiendan...} (souvent source non citée).
\item \esp{Promesa implícita / Asociación} : le produit = succès, beauté, liberté.
\item \esp{Apelación emocional} : peur (\esp{miedo}), joie, nostalgie, fierté.
\item \esp{Urgencia / Escasez} : \esp{Oferta limitada, ¡Solo hoy!}
\end{itemize}
\item \textbf{Juge de l'honnêteté (esprit critique).} \esp{¿Informa o manipula? ¿Qué información falta?} (prix, conditions, effets secondaires, source des chiffres).
\end{enumerate}
\textbf{Structures utiles pour l'oral/écrit :} \esp{El anuncio se dirige a... ; El eslogan utiliza la rima / la aliteración... ; La imagen sugiere que... ; Se trata de una técnica de [testimonio / urgencia / miedo] porque... ; El anuncio omite información esencial como el precio.}
\end{methode}

\begin{exoresolu}[Analyse d'une publicité — Type Bac]
\textbf{Énoncé.} Voici le texte d'une publicité radiophonique diffusée à Douala :

\emph{« ¡Joven! ¿Quieres triunfar en la vida? Bebe EnergyCam, la bebida de los campeones. El 90 \% de los futbolistas la prefieren. ¡EnergyCam: tu éxito empieza aquí! »}

Réponds en espagnol :
\begin{enumerate}
\item ¿Qué producto anuncia el texto y a quién se dirige? (2 pts)
\item Señala dos técnicas de persuasión utilizadas y explícalas. (4 pts)
\item ¿Te parece un anuncio honesto? Justifica tu respuesta. (4 pts)
\end{enumerate}
\tcblower
\textbf{Raisonnement.} Produit : boisson énergisante. Cible : jeunesse (\esp{¡Joven!}). Procédés : impératif + question directe (accroche), chiffre non sourcé (autorité fallacieuse), association produit/réussite (promesse implicite).

\textbf{Réponses rédigées :}
\begin{enumerate}
\item \esp{El texto anuncia una bebida energética llamada EnergyCam y se dirige sobre todo a los jóvenes, porque empieza con la palabra «¡Joven!» y habla de triunfar en la vida.}
\item \esp{Primera técnica: el uso de imperativos y preguntas directas («¿Quieres triunfar?», «¡Bebe!»), que crean una relación cercana y exigen una acción inmediata. Segunda técnica: el dato estadístico «el 90 \% de los futbolistas la prefieren», que da una apariencia científica sin citar ninguna fuente ni estudio (argumento de autoridad falso). Además, el eslogan asocia el producto con el éxito: «tu éxito empieza aquí» (promesa implícita).}
\item \esp{No me parece un anuncio completamente honesto. El porcentaje del 90 \% no está verificado y no sabemos qué futbolistas fueron consultados ni cuándo. Además, el anuncio sugiere que beber EnergyCam garantiza el éxito, lo cual es una promesa falsa y una manipulación de las aspiraciones de la juventud. En mi opinión, informa poco y manipula mucho.}
\end{enumerate}
\textbf{Conclusion :} Une bonne analyse cite toujours les mots exacts de la publicité entre guillemets (\esp{«...»}) et nomme le procédé technique.
\end{exoresolu}

\courssub{Fiche méthode 4 : Repérer une fausse information (fake news / bulo)}

\begin{methode}[Detectar un bulo en cinco pasos — Reflexe citoyen]
\begin{enumerate}
\item \textbf{Vérifie la source (\esp{la fuente}).} \esp{¿Quién publica?} Un medio conocido (\esp{El País, RFI, CRTV, La Nación}) ou une page anonyme (\esp{un blog, un forward de WhatsApp}) ? Cherche l'auteur (\esp{el autor}), la date (\esp{la fecha}) et les mentions légales (\esp{el aviso legal}).
\item \textbf{Lis au-delà du titre (\esp{el titular}).} Les titres \esp{clickbait} (\esp{¡Escándalo! ¡Increíble! ¡Urgente!}) cachent souvent un contenu vide, ancien ou sans rapport.
\item \textbf{Croise les sources (\esp{contrastar fuentes}).} Si aucun autre média sérieux ne reprend l'information, méfie-toi (\esp{desconfía}). Cherche sur Google / \esp{Google News} / sites de fact-checking (\esp{Maldita.es, Newtral, Africa Check, StopFake}).
\item \textbf{Examine les images / vidéos (\esp{imágenes / vídeos}).} Une photo peut être ancienne, prise dans un autre pays (inondation en Inde $\rightarrow$ présentée comme Garoua), ou générée par IA. Outils : \esp{Google Images (búsqueda inversa), TinEye, InVID}.
\item \textbf{Analyse le ton et l'appel à l'action.} Un bulo cherche l'\esp{emoción} (miedo, cólera, indignación) et presse de partager : \esp{¡Comparte antes de que lo borren! ¡Máxima difusión!} Une info vraie n'a pas besoin de suppliquer pour être partagée.
\end{enumerate}
\textbf{Réflexe final :} \esp{\textbf{Piensa antes de compartir.}} Une information vraie résiste à la vérification ; un bulo s'effondre dès qu'on le confronte aux faits.
\end{methode}

\begin{exoresolu}[Détection d'une fake news — Type Bac]
\textbf{Énoncé.} Un message circule sur WhatsApp à Yaoundé :

\emph{« ¡URGENTE! ¡COMPARTE CON TODOS TUS CONTACTOS! El gobierno va a cerrar todos los colegios durante un año por una nueva epidemia. Una fuente del ministerio lo confirma. ¡No dejes que te engañen, difunde la verdad! »}

\begin{enumerate}
\item Cita tres señales que indican que probablemente se trata de un bulo. (3 pts)
\item Explica qué deberías hacer antes de compartir este mensaje. (3 pts)
\item Traduis en français la première phrase du message. (2 pts)
\end{enumerate}
\tcblower
\textbf{Raisonnement.} Application des 5 pas : source anonyme (\esp{una fuente del ministerio}), ton alarmiste (majuscules, \esp{¡URGENTE!}), appel au partage viral, absence de date/lien vérifiable, contenu invraisemblable (fermeture 1 an).

\textbf{Réponses rédigées :}
\begin{enumerate}
\item \esp{Primera señal: la fuente es anónima («una fuente del ministerio»), sin nombre, cargo ni documento oficial. Segunda señal: el tono es alarmista, con mayúsculas y signos de exclamación («¡URGENTE!», «¡COMPARTE!»), típico de la manipulación emocional. Tercera señal: el mensaje exige compartir inmediatamente con todos los contactos, técnica viral para propagar el bulo sin verificación.}
\item \esp{Antes de compartir, debería verificar la información en medios de comunicación serios (radio nacional, prensa oficial, página web del Ministerio de Educación). También debería comprobar si otros medios fiables publican la misma noticia. Si ninguna fuente oficial lo confirma, no debo compartir el mensaje para no sembrar el pánico.}
\item « URGENT ! PARTAGE AVEC TOUS TES CONTACTS ! » (Remarque : \esp{comparte} = impératif affirmatif 2e pers. sg de \esp{compartir} ; \esp{todos tus contactos} = accord masculin pluriel).
\end{enumerate}
\textbf{Conclusion :} Repérer un bulo, c'est croiser trois critères invariables : \textbf{fuente + tono + verificación}.


\end{exoresolu}

\courssub{Fiche méthode 5 : Rédiger un commentaire personnel sur la libertad de expresión}

\begin{methode}[Escribir un comentario argumentado (80-120 palabras) — Épreuve de rédaction Bac]
\begin{enumerate}
\item \textbf{Introduction (1 phrase).} Présente le sujet et annonce ta thèse. \esp{La libertad de expresión es un derecho fundamental, pero su ejercicio plantea un dilema: ¿dónde están sus límites?}
\item \textbf{Paragraphe 1 : La valeur du droit (Arguments POUR).} \esp{En primer lugar,} / \esp{Por un lado,} + argument + \esp{Por ejemplo,} + exemple concret (Camerún, Guinea Ecuatorial, Espagne, réseau social). \esp{Una prensa libre permite denunciar la corrupción y defender los derechos humanos.}
\item \textbf{Paragraphe 2 : La nécessité des limites (Arguments CONTRE l'absolu).} \esp{Sin embargo,} / \esp{Por otro lado,} + argument + vocabulaire précis (\esp{difamación, calumnia, bulos, desinformación}). \esp{Este derecho no protege la mentira ni la injuria, que son delitos.}
\item \textbf{Conclusion (1 phrase).} Synthèse personnelle avec \esp{En conclusión,} / \esp{Por tanto,} + position nuancée. \esp{La libertad de expresión debe ir acompañada de responsabilidad y rigor.}
\item \textbf{Exigences grammaticales (points bonus).}
\begin{itemize}
\item Connecteurs variés : \esp{además, no obstante, por lo tanto, en definitiva}.
\item Subjonctif après expression de nécessité/valeur : \esp{Es fundamental que los periodistas \textbf{respeten} la verdad ; Es necesario que la ley \textbf{proteja} las fuentes.}
\item Conditionnel pour hypothèse : \esp{Una sociedad sin prensa libre \textbf{sería} ciega.}
\item Accords : \esp{la noticia falsa, el bulo viral, los medios de comunicación, las redes sociales}.
\end{itemize}
\end{enumerate}
\end{methode}

\begin{exoresolu}[Rédaction — Commentaire modèle Type Bac (≈ 110 mots)]
\textbf{Énoncé.} Rédige en espagnol un commentaire de 80 à 120 mots sur : \emph{« ¿Debe tener límites la libertad de expresión? »} (10 pts)
\tcblower
\textbf{Raisonnement.} Plan en 3 temps : Droit fondamental $\rightarrow$ Limites nécessaires (délits) $\rightarrow$ Synthèse responsable. Utilisation connecteurs + subjonctif + conditionnel.

\textbf{Production modèle :}

\esp{La libertad de expresión es un derecho fundamental de toda democracia. En primer lugar, gracias a una prensa libre, los ciudadanos pueden conocer la verdad y denunciar los abusos del poder. Por ejemplo, en muchos países africanos, los periodistas han revelado casos de corrupción que de otro modo quedarían ocultos.}

\esp{Sin embargo, este derecho no puede ser absoluto. La difamación, la calumnia y los bulos destruyen reputaciones y provocan conflictos sociales: no son opiniones, sino delitos de prensa. Es importante que los periodistas respeten la verdad y que las redes sociales no propaguen la desinformación.}

\esp{En conclusión, la libertad de expresión debe tener límites claros: proteger la crítica legítima, pero castigar la mentira. Un derecho sin responsabilidad se convierte en un peligro para todos.}

\textbf{Justifications grammaticales (pour l'auto-évaluation) :}
\begin{itemize}
\item \esp{pueden conocer} : \esp{poder} + infinitif (capacité).
\item \esp{han revelado} : Pretérito perfecto compuesto (faits passés avec répercussion présente).
\item \esp{Es importante que... respeten / propaguen} : \textbf{Subjonctif présent} (3e groupe \esp{-ar} $\rightarrow$ \esp{-e}, \esp{-ir/-er} $\rightarrow$ \esp{-a}) après expression impersonnelle de nécessité.
\item \esp{se convierte} : Verbe pronominal \esp{convertirse en} = « devenir ».
\item \esp{la crítica legítima} (fém.) / \esp{el bulo viral} (masc.) : accords corrects.
\end{itemize}
\textbf{Conclusion :} Le correcteur attend une position claire (\esp{debe tener límites}), deux arguments opposés, un exemple contextualisé, des connecteurs logiques et au moins un subjonctif bien employé.


\end{exoresolu}

\begin{attention}[Les 5 erreurs qui coûtent des points au Bac — Check-list finale]
\begin{enumerate}
\item \textbf{Fautes d'accent (orthographe lexicale).} \textbf{Écris systématiquement :} \esp{opinión, información, democracia, crítica, país, después, nación, constitución, televisión, manipulación, difamación, calumnia (pas d'accent), injuria (pas d'accent), eslogan (sg, pas d'accent), eslóganes (pl, accent sur ó).} \cle{Règle :} Mots en \esp{-ción / -sión} $\rightarrow$ accent sur \esp{ó/á} (agudas terminées en \esp{n/s}). Mots en \esp{-ía} (diphthongue) $\rightarrow$ pas d'accent (\esp{calumnia, injuria, publicidad}).
\item \textbf{Faux amis et calques lexicaux.}
\begin{itemize}
\item « Avertir (informer) » = \esp{avisar} ; « Avertir (mettre en garde) » = \esp{advertir} (\esp{advertir a alguien DE algo}). Ne pas confondre avec \esp{advertir} $\neq$ « advertiser » (publicité).
\item « Une réputation » = \esp{una reputación} (pas \esp{reputación} sans accent).
\item « Partager » = \esp{compartir} (jamais \esp{partir} = partir).
\item « Une censure » = \esp{la censura} (jamais \esp{la cesura} = césure poétique).
\item « Un reportage » = \esp{un reportaje} (pas \esp{reportaje} sans \esp{j} ? Non, s'écrit avec \esp{j} : \esp{reportaje}).
\end{itemize}
\item \textbf{Calques syntaxiques du français.}
\begin{itemize}
\item « Sans limites » $\neq$ \esp{sin límites} (correct) MAIS « est absolue sans limites » $\rightarrow$ \esp{es absoluta e ilimitada} ou \esp{no tiene límites}. Évite \esp{absoluta sin límites} (pléonasme maladroit).
\item « Il est important que les journalistes respectent » $\rightarrow$ \esp{Es importante que los periodistas \textbf{respeten}} (SUBJONCTIF, pas \esp{respetan}).
\item « Il faut que je vérifie » $\rightarrow$ \esp{Es necesario que \textbf{verifique}} (SUBJONCTIF, 1ère pers. sg).
\end{itemize}
\item \textbf{Accords en genre et nombre (pièges masculins en \esp{-a}).} \textbf{Masculins :} \esp{el problema, el tema, el sistema, el programa, el clima, el idioma, el mapa, el día, el planeta, el idioma.} \textbf{Féminins :} \esp{la noticia, la información, la opinión, la libertad, la prensa, la censura, la difamación, la calumnia, la injuria, la desinformación, la manipulación, la publicidad, la fuente, la verdad, la mentira, la responsabilidad.}
\item \textbf{Choix du temps et du mode (concordance).}
\begin{itemize}
\item Vérité générale / Habitude : \textbf{Indicatif présent} (\esp{La prensa informa, los bulos circulan}).
\item Hypothèse / Conséquence d'une condition : \textbf{Conditionnel} (\esp{Sería un error, una sociedad sería ciega}).
\item Après \esp{Es importante/necesario/urgente/fundamental que...} + \textbf{Subjonctif} (\esp{...que verifiquemos, ...que respeten, ...que no propaguen}).
\item Après \esp{Creo que / Pienso que / Es verdad que...} + \textbf{Indicatif} (certitude). Après \esp{No creo que / Dudo que / Es falso que...} + \textbf{Subjonctif} (doute/négation).
\end{itemize}
\end{enumerate}
\end{attention}

\begin{savaistu}[Le saviez-vous ? — Cameroun, Guinée équatoriale et liberté de presse]
\begin{itemize}
\item \textbf{Guinée équatoriale :} Seule pays africain hispanophone. La Constitution (art. 13) garantit la liberté d'expression, mais la loi de presse (1992, modifiée 2019) impose des peines de prison pour « diffamation » ou « injure » contre le Chef de l'État ou les institutions. Reporters Sans Frontières classe le pays régulièrement bas dans l'index mondial (109e/180 en 2023).
\item \textbf{Cameroun :} La loi n°2019/020 du 24 décembre 2019 sur la cybercriminalité punit la « propagation de fausses nouvelles » (art. 78) de 6 mois à 2 ans de prison et amende. Cette loi est critiquée par les organisations de défense des droits de l'homme (Amnesty International, RSF) comme outil de musellement de la presse en ligne et des lanceurs d'alerte.
\item \textbf{Espagne :} Modèle de transition démocratique. La loi organique 1/1982 protège le droit à l'honneur (\esp{derecho al honor}) et à l'intimité, base juridique des procès pour \esp{difamación} et \esp{injuria}. Le \esp{Tribunal Constitucional} veille à l'équilibre entre \esp{libertad de expresión} et \esp{derecho al honor}.
\item \textbf{Terminologie ONU/UNESCO :} \esp{Día Mundial de la Libertad de Prensa} = 3 mai. Thème 2024 : \esp{« Una prensa para el planeta: el periodismo frente a la crisis ambiental »}.
\end{itemize}
\end{savaistu}

\begin{experience}[Activité de classe : « Tribunal de la presse » — Expression orale guidée]
\textbf{Objectif :} Mettre en œuvre le lexique des délits de presse, l'argumentation (\esp{por/para}, subjonctif, conditionnel) et l'esprit critique.
\textbf{Déroulement (30 min) :}
\begin{enumerate}
\item \textbf{Groupes de 4-5 :} Rôles assignés : \esp{El Juez} (président), \esp{El Fiscal} (accusation), \esp{El Abogado defensor} (défense), \esp{El Periodista} (accusé), \esp{El Testigo} (optionnel).
\item \textbf{Cas pratiques (choisir un) :}
\begin{itemize}
\item \textit{Caso A :} Un journaliste publie sur Facebook : « Le Ministre X a détourné 500 millions FCFA » sans preuve. Le Ministre porte plainte pour \esp{calumnia}.
\item \textit{Caso B :} Un influenceur à Douala partage une vidéo : « L'eau de la SNEC contient du poison ! » (faux). La SNEC attaque pour \esp{difamación} et \esp{bulos}.
\item \textit{Caso C :} Un dessinateur de presse caricature le Président avec des traits exagérés. Le parquet ouvre une enquête pour \esp{injuria} à Chef d'État.
\end{itemize}
\item \textbf{Préparation (10 min) :} Chaque rôle prépare ses arguments avec structures : \esp{Mi defendido no cometió calumnia porque...}, \esp{El fiscal sostiene que...}, \esp{Es necesario que se proteja...}, \esp{Si condenamos esto, sería...}.
\item \textbf{Audience (15 min) :} Simulation. Le Juge rend un verdict motivé en espagnol.
\item \textbf{Débat (5 min) :} La classe vote : la décision respecte-t-elle l'équilibre \esp{libertad de expresión / derecho al honor} ?
\end{enumerate}
\textbf{Évaluation :} Grille : Vocabulaire précis (5 pts) + Structures complexes Subj/Cond (5 pts) + Argumentation logique (5 pts) + Prononciation/Intonation (5 pts).
\end{experience}

\newpage

\coursec{Exercices}

\coursec{Lectura y comentario: libertad de expresión y delitos de prensa; análisis de publicidad y fake news}

\courssub{Je vérifie mes connaissances}

\begin{exercice}[QCM : notions clés]
Chaque question ne comporte qu'une seule réponse correcte. Cochez la bonne case.
\begin{enumerate}
    \item Le terme \esp{difamación} désigne :
    \begin{itemize}
        \item[a)] \esp{Una información falsa difundida intencionalmente para dañar la reputación.}
        \item[b)] \esp{Una opinión crítica expresada en un editorial.}
        \item[c)] \esp{Una publicidad engañosa.}
    \end{itemize}
    \item La \esp{calumnia} se distingue de la diffamation parce qu'elle :
    \begin{itemize}
        \item[a)] \esp{Es siempre cometida por un periodista.}
        \item[b)] \esp{Imputa un hecho delictivo preciso y falso.}
        \item[c)] \esp{Solo concierne a las personas públicas.}
    \end{itemize}
    \item Un \esp{bulo} est :
    \begin{itemize}
        \item[a)] \esp{Una ley que limita la libertad de prensa.}
        \item[b)] \esp{Un rumor o falsa noticia propagada como verdadera, a menudo en redes.}
        \item[c)] \esp{Un género periodístico de investigación.}
    \end{itemize}
    \item L'\esp{eslogan} d'une publicité remplit principalement la fonction :
    \begin{itemize}
        \item[a)] \esp{Informar sobre las características técnicas del producto.}
        \item[b)] \esp{Fijar la marca en la memoria con una frase breve y pegadiza.}
        \item[c)] \esp{Dar las menciones legales obligatorias.}
    \end{itemize}
    \item La \esp{desinformación} se caractérise par :
    \begin{itemize}
        \item[a)] \esp{La ausencia total de información sobre un tema.}
        \item[b)] \esp{La difusión deliberada de informaciones falsas o manipuladas para engañar.}
        \item[c)] \esp{Un error de imprenta en un artículo de prensa.}
    \end{itemize}
\end{enumerate}
\end{exercice}

\begin{exercice}[Vrai ou Faux justifié sur le texte support]
Lisez (ou relisez) le texte support \og Libertad de expresión: ¿un derecho sin límites? \fg\ puis indiquez si les affirmations suivantes sont vraies (V) ou fausses (F). Justifiez chaque réponse en citant le numéro de ligne(s) correspondant(s).

\begin{textebox}[Texte support : Libertad de expresión: ¿un derecho sin límites?]
La libertad de expresión es un pilar fundamental de toda democracia, pero no es un derecho absoluto. La legislación internacional y las constituciones nacionales establecen límites claros: no se puede incitar al odio, a la violencia ni cometer delitos de difamación o calumnia. En la era digital, la propagación de bulos y la desinformación plantean nuevos retos: ¿cómo proteger el derecho a informar sin caer en la manipulación? La censura previa está prohibida en la mayoría de los Estados de derecho, pero la responsabilidad posterior —responder ante un juez por lo publicado— es la garantía del equilibrio. Los medios de comunicación tienen la obligación de verificar las fuentes antes de publicar; el periodismo riguroso es el mejor antídoto contra la mentira. En Camerún, como en España o México, los periodistas a veces sufren presiones, pero la ley debe proteger su labor, no silenciarla. En definitiva, la libertad termina donde comienza el derecho al honor, a la intimidad y a la seguridad de los demás.
\end{textebox}

\begin{enumerate}
    \item La libertad de expresión es un derecho absoluto sin ninguna restricción legal.
    \item La censura previa está permitida en los Estados de derecho cuando hay riesgo de desinformación.
    \item La difamación y la calumnia son delitos de prensa que limitan la libertad de expresión.
    \item Los periodistas no tienen la obligación de verificar las fuentes; la rapidez es lo más importante.
    \item El texto menciona que en Camerún los periodistas nunca sufren presiones.
    \item La responsabilidad posterior significa responder ante un juez por lo publicado.
\end{enumerate}
\end{exercice}

\begin{exercice}[Vocabulaire : article à trous]
Complétez l'article suivant avec les mots de la liste ci-dessous. Chaque mot s'utilise une seule fois. Mettez le verbe à la forme convenable si nécessaire.

\emph{Liste :} \esp{difamación, calumnia, censura, bulo, desinformación, publicidad, eslogan, manipulación, verificar, libertad de expresión}

\begin{textebox}[Artículo: Los retos de la información en la era digital]
En la sociedad actual, la \_\_\_\_\_\_\_\_\_\_\_\_ es un derecho fundamental, pero conlleva responsabilidades. Los códigos penales castigan la \_\_\_\_\_\_\_\_\_\_\_\_ (imputación de un delito falso) y la \_\_\_\_\_\_\_\_\_\_\_\_ (daño a la reputación sin imputar delito). Algunos gobiernos recurren a la \_\_\_\_\_\_\_\_\_\_\_\_ previa para controlar los medios, lo que vulnera el derecho a informar. En las redes sociales, un \_\_\_\_\_\_\_\_\_\_\_\_ puede volverse viral en minutos y causar daños irreparables. La \_\_\_\_\_\_\_\_\_\_\_\_ organizada busca polarizar a la ciudadanía. Ante esto, los periodistas deben \_\_\_\_\_\_\_\_\_\_\_\_ rigurosamente cada dato. La \_\_\_\_\_\_\_\_\_\_\_\_ comercial también usa técnicas de persuasión: un buen \_\_\_\_\_\_\_\_\_\_\_\_ vende más que mil argumentos. La educación mediática es clave para detectar la \_\_\_\_\_\_\_\_\_\_\_\_ y ejercer un pensamiento crítico.
\end{textebox}
\end{exercice}

\begin{aretenir}[Lexique clé : Libertad de expresión, delitos de prensa y medios]
\begin{center}
\begin{tabularx}{\linewidth}{|X|X|X|}
\hline
\rowcolor{popblueL} \textbf{Término} & \textbf{Definición / Contexto} & \textbf{Traduction / Note} \\
\hline
\esp{libertad de expresión} & Derecho fundamental (Art. 19 DUDH, Art. 13 CADH) & Liberté d'expression \\
\hline
\esp{censura previa} & Prohibición gubernamental de publicar antes de la difusión & Censure préalable (interdite en État de droit) \\
\hline
\esp{responsabilidad posterior} & Rendir cuentas ante la justicia tras la publicación & Responsabilité a posteriori \\
\hline
\esp{difamación} & Divulgación de informaciones falsas que dañan la fama (sin imputar delito) & Diffamation (atteinte à l'honneur) \\
\hline
\esp{calumnia} & Imputación falsa de un delito concreto & Calomnie (infraction plus grave) \\
\hline
\esp{injuria} & Expresiones que menoscaban la dignidad (no requiere falsedad) & Injure (atteinte à la dignité) \\
\hline
\esp{bulo} & Noticia falsa difundida como verdadera (redes sociales) & Rumeur, \emph{fake news}, intox \\
\hline
\esp{desinformación} & Difusión deliberada de información falsa para manipular & Désinformation (stratégie intentionnelle) \\
\hline
\esp{manipulación} & Uso sesgado de la información para influir & Manipulation \\
\hline
\esp{verificar} / \esp{contrastar} & Comprobar la veracidad de las fuentes & Vérifier, recouper \\
\hline
\esp{publicidad} & Comunicación comercial persuasive & Publicité \\
\hline
\esp{eslogan} & Frase breve y memorable de una campaña & Slogan \\
\hline
\esp{publicidad engañosa} & Publicidad con datos falsos que induce a error & Publicité trompeuse \\
\hline
\end{tabularx}
\end{center}
\end{aretenir}

\begin{exercice}[Conjugaison : modes et temps du commentaire]
Conjuguez les verbes entre parenthèses au temps et au mode appropriés (indicatif, subjonctif, conditionnel, impératif). Attention à la concordance des temps.

\begin{enumerate}
    \item Es fundamental que los periodistas \_\_\_\_\_\_\_\_\_\_\_\_ (verificar) la información antes de publicarla.
    \item Si el gobierno \_\_\_\_\_\_\_\_\_\_\_\_ (imponer) censura previa, violaría la Constitución.
    \item El juez ordenó que el medio \_\_\_\_\_\_\_\_\_\_\_\_ (rectificar) la noticia falsa en la misma portada.
    \item No creo que ese \_\_\_\_\_\_\_\_\_\_\_\_ (ser) un bulo; la fuente es muy fiable.
    \item \_\_\_\_\_\_\_\_\_\_\_\_ (No compartir) mensajes de odio en las redes sociales.
    \item Me gustaría que la ley \_\_\_\_\_\_\_\_\_\_\_\_ (proteger) mejor a los reporteros en zonas de conflicto.
    \item Aunque la noticia \_\_\_\_\_\_\_\_\_\_\_\_ (ser) cierta, el medio fue multado por difamación.
    \item Es urgente que las plataformas \_\_\_\_\_\_\_\_\_\_\_\_ (actuar) contra la desinformación coordinada.
\end{enumerate}
\end{exercice}

\courssub{Je m'entraîne}

\begin{methode}[Méthode : Analyser une publicité]
\textbf{Étapes pour une analyse structurée (commentaire écrit ou oral) :}
\begin{enumerate}
    \item \textbf{Identifier l'émetteur} (\emph{emisor}) : marque, organisme, association. Quelle est sa notoriété ? Son éthique ?
    \item \textbf{Identifier le destinataire} (\emph{destinatario / público objetivo}) : âge, CSP, centres d'intérêt, valeurs. Quel \emph{insight} (vérité consommateur) vise-t-il ?
    \item \textbf{Analyser le message} (\emph{mensaje}) : \textbf{Promesse} (bénéfice principal), \textbf{Preuve} (arguments, labels, ingrédients), \textbf{Ton} (humour, émotion, rationnel, engagement).
    \item \textbf{Étudier le slogan} (\emph{eslogan}) : brièveté, musicalité (rimes, allitérations), jeux de mots, impératif, possession (\emph{tu}, \emph{mi}).
    \item \textbf{Décrypter l'image et la mise en page} (\emph{imagen, composición, colores}) : symbolique des couleurs, cadrage, personnages, typographie.
    \item \textbf{Évaluer l'éthique} : \emph{Publicidad responsable} (vérifiable, respectueuse) vs \emph{Publicidad engañosa / subliminal / sexiste}. Y a-t-il \emph{manipulación} ? \emph{Estereotipos} ?
    \item \textbf{Conclure} sur l'efficacité probable et la portée culturelle/sociale.
\end{enumerate}
\end{methode}

\begin{methode}[Méthode : Détecter une fake news (bulo / desinformación)]
\begin{enumerate}
    \item \textbf{Source} : URL suspecte (ex: \emph{.co}, \emph{.xyz} imitant un grand média), auteur anonyme, site satirique non assumé.
    \item \textbf{Titre} : \emph{Clickbait} (sensationnaliste, majuscules, points d'exclamation), ne correspond pas au contenu.
    \item \textbf{Date et lieu} : Info ancienne recyclée (\emph{old news}), contexte modifié.
    \item \textbf{Preuves} : Absence de liens vers sources primaires, citations inventées, experts inexistants.
    \item \textbf{Images/Vidéos} : Hors contexte, retouchées (vérifier avec \emph{Google Images}, \emph{TinEye}, \emph{InVID}).
    \item \textbf{Biais cognitifs} : Joue sur la peur, la colère, la confirmation de vos opinions (\emph{sesgo de confirmación}).
    \item \textbf{Recoupement} : Chercher l'info sur 3 médias de référence (\emph{contrastar fuentes}). Sites de fact-checking : \emph{Maldita.es}, \emph{Newtral}, \emph{Africa Check}, \emph{Dubawa}.
\end{enumerate}
\end{methode}

\begin{exercice}[Traduction : thème et version]
Traduisez les phrases suivantes. Respectez le registre de langue (journalistique / formel).

\begin{enumerate}
    \item La liberté de la presse ne signifie pas le droit de calomnier impunément.
    \item Los bulos se propagan más rápido que las noticias verificadas.
    \item Il est indispensable que les citoyens apprennent à détecter la manipulation médiatique.
    \item La publicidad engañosa atenta contra los derechos del consumidor.
    \item Le journaliste a été poursuivi pour diffamation après avoir publié l'article.
    \item No todo lo que circula en WhatsApp es verdad; hay que contrastar las fuentes.
\end{enumerate}
\end{exercice}

\begin{exercice}[Transformation de phrases]
Transformez les phrases selon les indications. Conservez le sens original.

\begin{enumerate}
    \item \emph{Actif $\rightarrow$ Passif reflejo (se) :} Los medios difunden bulos constantemente.
    \item \emph{Style direct $\rightarrow$ Style indirect :} El director dijo: «Publicaremos la rectificación mañana».
    \item \emph{Indicatif $\rightarrow$ Subjonctif (nécessité) :} Es necesario que la ley protege a los periodistas.
    \item \emph{Phrase affirmative $\rightarrow$ Négation partielle (ningún) :} Algunos políticos usan la desinformación como arma.
    \item \emph{Verbe personnel $\rightarrow$ Impersonnel (se) :} La policía investiga el origen del bulo.
    \item \emph{Conditionnel $\rightarrow$ Si + imparfait du subjonctif :} Me gustaría que hubiera más transparencia.
\end{enumerate}
\end{exercice}

\begin{exercice}[Texte à trous : analyse publicitaire]
Complétez le texte avec les mots appropriés (un seul mot par trait d'union). Les premiers lettres sont données.

\begin{textebox}[Análisis de un anuncio de una marca de café camerounaise]
El anuncio muestra a un joven agricultor en una plantación de café en las montañas del Oeste de Camerún. El \_\_\_\_\_\_\_\_ (e\_\_\_\_\_) es la cooperativa \og Café Excelencia \fg. El \_\_\_\_\_\_\_\_ (d\_\_\_\_\_\_\_) son consumidores urbanos que valoran el producto local. El \_\_\_\_\_\_\_\_ (m\_\_\_\_\_\_) destaca la calidad, el comercio justo y el orgullo nacional. El \_\_\_\_\_\_\_\_ (e\_\_\_\_\_\_) \og Tu café, tu tierra, tu orgullo \fg usa la repetición del posesivo \_\_\_\_\_\_\_\_ (t\_\_) para crear cercanía. La \_\_\_\_\_\_\_\_ (i\_\_\_\_\_) combina colores verdes y dorados, símbolos de naturaleza y riqueza. No hay \_\_\_\_\_\_\_\_ (m\_\_\_\_\_\_\_\_\_) oculta: los datos sobre el precio justo al productor son \_\_\_\_\_\_\_\_ (v\_\_\_\_\_\_\_\_\_\_\_\_). Es publicidad \_\_\_\_\_\_\_\_ (r\_\_\_\_\_\_\_\_\_\_\_) que refuerza la identidad cultural.
\end{textebox}
\end{exercice}

\begin{exercice}[Appariements : concepts et définitions/exemples]
Reliez chaque concept de la colonne A à sa définition ou exemple correspondant dans la colonne B (une seule correspondance par ligne).

\begin{center}
\begin{tabularx}{\linewidth}{|X|X|}
\hline
\rowcolor{popblueL} \textbf{A. Concepto} & \textbf{B. Definición / Ejemplo} \\
\hline
1. Censura previa &  \\
\hline
2. Difamación &  \\
\hline
3. Calumnia &  \\
\hline
4. Bulo &  \\
\hline
5. Eslogan &  \\
\hline
6. Desinformación &  \\
\hline
\end{tabularx}
\end{center}
\emph{Options colonne B :}
\begin{itemize}
    \item A. \esp{¡Compra ya! Oferta limitada}
    \item B. \esp{Prohibición de publicar antes de salir a la luz}
    \item C. \esp{Dañar la fama de alguien con informaciones falsas (sin imputar delito)}
    \item D. \esp{Imputar falsamente un delito concreto a alguien}
    \item E. \esp{Falsa noticia viral: \og El agua cura el cáncer en 3 días \fg}
    \item F. \esp{Campaña coordinada para hacer creer que las elecciones fueron fraudulentas sin pruebas}
\end{itemize}
\end{exercice}

\courssub{Je me prépare au Bac}

\begin{methode}[Méthode : Commentaire de texte type Baccalauréat (Espagnol LV2)]
\textbf{Structure attendue (4 parties) :}
\begin{enumerate}
    \item \textbf{Introduction} (3-4 lignes) : \emph{Accroche} (sujet d'actualité), \emph{Présentation} (titre, auteur si connu, source, date, thème), \emph{Problématique} (question ouverte à laquelle le texte répond), \emph{Annonce du plan} (2-3 axes).
    \item \textbf{Développement} (2-3 paragraphes) : Un axe = une idée force. \textbf{Citation courte} (numéro de ligne) $\rightarrow$ \textbf{Analyse} (explication, vocabulaire, figures de style, argumentation de l'auteur) $\rightarrow$ \textbf{Lien/Transition}.
    \item \textbf{Analyse linguistique intégrée} : Relever des faits de langue (temps, modes, voix, connecteurs, champs lexicaux) \textbf{au service du sens} (pas de liste gratuite).
    \item \textbf{Conclusion} (3-4 lignes) : \emph{Bilan} (réponse à la problématique), \emph{Ouverture} (autre texte, actualité camerounaise/hispanique, avis personnel nuancé).
\end{enumerate}
\textbf{Consignes de langue :} Espagnol correct, registre soutenu, connecteurs logiques (\emph{por tanto, en consecuencia, sin embargo, además, en primer lugar}), pas de traduction mot à mot.
\end{methode}

\begin{exercice}[Commentaire de texte type Baccalauréat]
Lisez attentivement le texte suivant, puis répondez aux questions en espagnol. La qualité de la langue (correction, richesse lexicale, cohérence) sera évaluée.

\begin{textebox}[Redes sociales y libertad de expresión: el nuevo ágora digital]
Las redes sociales han transformado radicalmente el espacio público. Hoy, cualquier ciudadano con un smartphone puede convertirse en emisor de información, saltándose los filtros tradicionales de los medios de comunicación. Esta democratización de la palabra tiene luces y sombras. Por un lado, permite denunciar abusos, movilizar conciencias y dar voz a minorías silenciadas; pensemos en los movimientos \#YoTambien o en las protestas climáticas globales. Por otro, la ausencia de verificación previa facilita la propagación de bulos, discursos de odio y teorías de la conspiración. Los algoritmos, diseñados para maximizar la atención, premian lo sensacional y lo polarizante, no lo riguroso. \textbf{Se comparten millones de contenidos sin verificar cada minuto.} Así, la desinformación viaja seis veces más rápido que la verdad verificada, según estudios del MIT. Los Estados intentan regular: la Ley Europea de Servicios Digitales obliga a las plataformas a retirar contenidos ilegales y a transparentar sus algoritmos. Pero el riesgo de una censura privatizada —donde una corporación decide qué es verdad— es real. En Camerún, el cierre de internet en 2017-2018 en las regiones anglófonas mostró que el poder también apaga la red para silenciar. La solución no es menos libertad, sino más educación mediática: formar ciudadanos críticos que sepan verificar, contrastar y no compartir sin pensar. La libertad de expresión exige responsabilidad; el derecho a opinar no incluye el derecho a inventar hechos.
\end{textebox}

\begin{enumerate}
    \item \textbf{Compréhension globale :} Résumez en 3-4 phrases en espagnol l'idée principale du texte.
    \item \textbf{Vocabulaire en contexte :} Expliquez en espagnol le sens des expressions suivantes telles qu'elles sont utilisées dans le texte : \emph{a) saltarse los filtros tradicionales ; b) luces y sombras ; c) censura privatizada.}
    \item \textbf{Analyse linguistique :} Relevez dans le texte \textbf{deux verbes au subjonctif} et justifiez leur emploi (valeur). Relevez \textbf{une phrase à la voix passive réfléchie (se + verbe)} à l'indicatif et transformez-la à la voix active.
    \item \textbf{Argumentation du texte :} Selon l'auteur, quels sont les deux grands dangers des réseaux sociaux pour la qualité de l'information ? Citez les lignes.
    \item \textbf{Comparaison contextuelle :} Le texte mentionne le Cameroun et l'Europe. Quelle différence de régulation pointe l'auteur ? (2-3 lignes)
    \item \textbf{Expression personnelle (environ 15 lignes) :} \og ¿Crees que la educación mediática es suficiente para combatir la desinformación o hacen falta leyes más estrictas? Argumenta tu respuesta con ejemplos concretos. \fg
\end{enumerate}
\end{exercice}

\begin{exercice}[Thème et Version]
\textbf{A. Thème (Français $\rightarrow$ Espagnol) — Traduisez le paragraphe suivant en espagnol soigné (registre journalistique).}

La lutte contre les fausses nouvelles ne doit pas servir de prétexte pour museler la presse. Dans plusieurs pays africains, des lois sur la \og cybercriminalité \fg ont été votées, mais leur rédaction vague permet de poursuivre des journalistes pour de simples critiques du pouvoir. La communauté internationale rappelle que la diffamation ne doit pas être un délit pénal, mais une affaire civile. Au Cameroun, les organisations de la société civile réclament la dépénalisation des délits de presse et la création d'un organe de régulation indépendant. Seul un cadre juridique clair, respectueux des standards internationaux, garantira une information libre et responsable.

\textbf{B. Version (Espagnol $\rightarrow$ Français) — Traduisez le paragraphe suivant en français correct.}

\begin{textebox}[Texte à traduire]
La publicidad política en redes sociales plantea dilemas éticos inéditos. El microtargeting permite enviar mensajes distintos a chaque votante según su perfil psicológico, fragmentando el debate público en burbujas opacas. ¿Quién controla la veracidad de esos anuncios? Las plataformas se lavan las manos alegando que no son medios, pero su poder de amplificación es masivo. Urge una legislación que exija transparencia en la financiación y el contenido de la propaganda política digital, so pena de que la democracia se reduzca a una subasta de datos personales.
\end{textebox}
\end{exercice}

\begin{exercice}[Sujet de rédaction type Bac]
Traitez l'un des deux sujets suivants, au choix. Votre rédaction comportera \textbf{environ 15 lignes} (130-180 mots). Vous ferez un plan apparent (introduction, développement en 2-3 parties, conclusion). La correction de la langue, la richesse du vocabulaire thématique et la cohérence de l'argumentation seront notées.

\begin{enumerate}
    \item \og ¿Piensas que la libertad de expresión debe tener límites? Justifica tu respuesta con ejemplos. \fg
    \item \og Las redes sociales han democratizado la información, pero también han multiplicado los bulos. ¿Cómo conciliar libertad y responsabilidad en la era digital? \fg
\end{enumerate}
\end{exercice}



\newpage

\coursec{Corrigés des exercices}

\coursec{Corrigés détaillés — Leçon E26 : Liberté d'expression, presse, publicité et \emph{fake news}}

\begin{corrige}[Exercice 1 — QCM : notions clés]
\begin{enumerate}
    \item \textbf{a)} La \cle{difamación} est une information fausse diffusée intentionnellement pour nuire à la réputation, \textbf{sans} imputer un fait précis constitutif d'infraction. \textit{Point de vigilance :} ne pas confondre avec l'opinion critique, protégée par la liberté d'expression.
    \item \textbf{b)} La \cle{calumnia} impute faussement un \textbf{délit précis} (\esp{« Ha robado », « Ha matado »}). C'est la distinction juridique fondamentale avec la diffamation : \esp{calumnia = imputación falsa de un delito}.
    \item \textbf{b)} Le \cle{bulo} est une fausse nouvelle propagée comme vraie, surtout sur les réseaux sociaux. Synonymes : \esp{noticia falsa}, \esp{rumor}. \textit{Faux ami :} \esp{bulo} ne veut pas dire « boule » ni « bulletin ».
    \item \textbf{b)} L'\cle{eslogan} fixe la marque dans la mémoire grâce à une phrase brève et \esp{pegadiza} (accrocheuse). Les informations techniques relèvent du corps du texte publicitaire, pas du slogan.
    \item \textbf{b)} La \cle{desinformación} est \textbf{intentionnelle} : diffusion délibérée de fausses informations pour tromper. Une simple erreur (c) est une \esp{información errónea}, sans volonté de nuire.
\end{enumerate}
\textbf{Barème indicatif :} 1 point par réponse correcte (5 points).
\end{corrige}

\begin{corrige}[Exercice 2 — Vrai ou Faux justifié]
\begin{enumerate}
    \item \textbf{F.} Ligne 2 : \esp{« no es un derecho absoluto »}. Le texte affirme exactement le contraire.
    \item \textbf{F.} Ligne 6 : \esp{« La censura previa está prohibida en la mayoría de los Estados de derecho »}. Elle est interdite, pas permise.
    \item \textbf{V.} Lignes 3-4 : \esp{« límites claros: no se puede incitar al odio... ni cometer delitos de difamación o calumnia »}.
    \item \textbf{F.} Ligne 8 : \esp{« Los medios de comunicación tienen la obligación de verificar las fuentes antes de publicar »}. La vérification est une obligation, la rapidité ne prime pas.
    \item \textbf{F.} Ligne 9 : \esp{« los periodistas a veces sufren presiones »}. Le texte dit « parfois », l'affirmation dit « jamais » : attention aux quantificateurs dans les énoncés V/F.
    \item \textbf{V.} Ligne 7 : \esp{« la responsabilidad posterior —responder ante un juez por lo publicado— »}. La définition est littérale dans le texte.
\end{enumerate}
\textbf{Barème indicatif :} 0,5 point par réponse V/F + 0,5 point par justification citée (6 points). \textbf{Point de vigilance :} une réponse juste sans citation de ligne ne vaut que la moitié des points.
\end{corrige}

\begin{corrige}[Exercice 3 — Article à trous]
Texte complété :
\begin{textebox}[Article complété]
En la sociedad actual, la \textbf{libertad de expresión} es un derecho fundamental, pero conlleva responsabilidades. Los códigos penales castigan la \textbf{calumnia} (imputación de un delito falso) y la \textbf{difamación} (daño a la reputación sin imputar delito). Algunos gobiernos recurren a la \textbf{censura} previa para controlar los medios, lo que vulnera el derecho a informar. En las redes sociales, un \textbf{bulo} puede volverse viral en minutos y causar daños irreparables. La \textbf{desinformación} organizada busca polarizar a la ciudadanía. Ante esto, los periodistas deben \textbf{verificar} rigurosamente cada dato. La \textbf{publicidad} comercial también usa técnicas de persuasión: un buen \textbf{eslogan} vende más que mil argumentos. La educación mediática es clave para detectar la \textbf{manipulación} y ejercer un pensamiento crítico.
\end{textebox}

\textbf{Justifications :}
\begin{itemize}
    \item \esp{calumnia} avant \esp{difamación} : l'ordre est imposé par les définitions entre parenthèses (imputation d'un délit = calomnie ; atteinte à la réputation = diffamation).
    \item \esp{un bulo} : masculin singulier, article indéfini \esp{un}.
    \item \esp{deben verificar} : infinitif obligatoire après le verbe modal \esp{deber}.
    \item \esp{un buen eslogan} : apocope de \esp{bueno} devant nom masculin singulier (\esp{buen}).
\end{itemize}
\textbf{Barème indicatif :} 1 point par mot correctement placé (10 points).
\end{corrige}

\begin{corrige}[Exercice 4 — Conjugaison : modes et temps]
\begin{enumerate}
    \item \esp{Es fundamental que los periodistas \textbf{verifiquen} la información antes de publicarla.} — Subjonctif présent : expression impersonnelle de nécessité (\esp{Es fundamental que} + subjonctif). Verbe \esp{verificar} (régulier -ar).
    \item \esp{Si el gobierno \textbf{impusiera} censura previa, violaría la Constitución.} — Subjonctif imparfait : structure hypothétique \esp{Si} + imparfait du subjonctif $\rightarrow$ conditionnel. \esp{Impusiera} (ou \esp{impusiese}), irrégulier : radical \esp{impus-} (de \esp{imponer}).
    \item \esp{El juez ordenó que el medio \textbf{rectificara} la noticia falsa en la misma portada.} — Subjonctif imparfait : verbe de volonté (\esp{ordenar}) au passé simple $\rightarrow$ concordance des temps avec l'imparfait du subjonctif. \esp{Rectificara} (ou \esp{rectificase}), régulier -ar.
    \item \esp{No creo que ese \textbf{sea} un bulo; la fuente es muy fiable.} — Subjonctif présent : \esp{no creer que} exprime le doute. \esp{Sea} : subjonctif irrégulier de \esp{ser}.
    \item \esp{\textbf{No compartan} mensajes de odio en las redes sociales.} — Impératif négatif (\esp{ustedes}) = forme du subjonctif présent. Accepté aussi : \esp{No compartas} (\esp{tú}). \textit{Piège :} jamais \esp{no compartir} (infinitif) ni \esp{no comparten} (indicatif).
    \item \esp{Me gustaría que la ley \textbf{protegiera} mejor a los reporteros en zonas de conflicto.} — Subjonctif imparfait après \esp{me gustaría que} (souhait non réalisé). Attention à l'orthographe : \esp{protegiera} garde le \esp{g} du radical \esp{proteg-} (songe \esp{g} dur devant \esp{ie}).
    \item \esp{Aunque la noticia \textbf{sea} cierta, el medio fue multado por difamación.} — Subjonctif présent : \esp{aunque} + subjonctif = concession hypothétique (« même si »). L'indicatif \esp{es} serait possible si le fait était présenté comme avéré et connu.
    \item \esp{Es urgente que las plataformas \textbf{actúen} contra la desinformación coordinada.} — Subjonctif présent : nécessité. Orthographe : \esp{actúen} avec accent (radical accentué \esp{actú-} du présent de l'indicatif \esp{actúan} $\rightarrow$ subjonctif \esp{actúen}).
\end{enumerate}
\textbf{Barème indicatif :} 1 point par verbe (forme + mode corrects) ; 0,5 si le mode est juste mais la forme fautive (8 points). \textbf{Point de vigilance :} la concordance des temps (passé au principal $\rightarrow$ imparfait du subjonctif) est le critère majeur.
\end{corrige}

\begin{corrige}[Exercice 5 — Traduction : thème et version]
\begin{enumerate}
    \item \esp{La libertad de prensa no significa el derecho a calumniar impunemente.} — \textit{Points de vigilance :} \esp{libertad de prensa} (pas \esp{de la prensa}) ; « calomnier » = \esp{calumniar} ; « impunément » = \esp{impunemente} (adverbe en \esp{-mente}).
    \item « Les fausses nouvelles se propagent plus vite que les nouvelles vérifiées. » \esp{Los bulos se propagan más rápido que las noticias verificadas.} — \esp{bulos} = fausses nouvelles/rumeurs ; \esp{se propagan} = pronominal, 3\up{e} personne du pluriel ; \esp{verificadas} = participe passé accordé avec \esp{noticias} (fém. plur.).
    \item \esp{Es indispensable que los ciudadanos aprendan a detectar la manipulación mediática.} — Subjonctif \esp{aprendan} obligatoire après \esp{Es indispensable que} ; \esp{aprender a} + infinitif ; \esp{mediática} avec accent.
    \item « La publicité trompeuse porte atteinte aux droits du consommateur. » \esp{La publicidad engañosa atenta contra los derechos del consumidor.} — \esp{atentar contra} = porter atteinte à ; \esp{engañosa} = trompeuse ; \esp{del consumidor} = contraction \esp{de + el}.
    \item \esp{El periodista fue procesado por difamación tras publicar el artículo.} — Passif avec \esp{ser} + participe accordé (\esp{fue procesado}) ; « après avoir publié » = \esp{tras publicar} (préposition + infinitif, tournure idiomatique espagnole). Accepté : \esp{fue denunciado por difamación}.
    \item « Tout ce qui circule sur WhatsApp n'est pas vrai ; il faut recouper les sources. » \esp{No todo lo que circula por WhatsApp es verdad; hay que contrastar las fuentes.} — \esp{No todo lo que... es verdad} : négation partielle, l'ordre des mots compte ; \esp{hay que} + infinitif = « il faut » ; \esp{contrastar las fuentes} = recouper/vérifier les sources.
\end{enumerate}
\textbf{Barème indicatif :} 2 points par phrase (1 point lexique, 1 point grammaire) (12 points).
\end{corrige}

\begin{corrige}[Exercice 6 — Transformation de phrases]
\begin{enumerate}
    \item \esp{\textbf{Se difunden} bulos constantemente (por los medios).} — Passif réfléchi : \esp{se} + verbe accordé avec le sujet patient \esp{bulos} (pluriel $\rightarrow$ \esp{difunden}).
    \item \esp{El director dijo que \textbf{publicarían} la rectificación \textbf{al día siguiente}.} — Discours indirect : futur \esp{publicaremos} $\rightarrow$ conditionnel \esp{publicarían} (concordance avec \esp{dijo}) ; \esp{mañana} $\rightarrow$ \esp{al día siguiente}.
    \item \esp{Es necesario que la ley \textbf{proteja} a los periodistas.} — \esp{Es necesario que} + subjonctif présent : \esp{proteja} (radical \esp{protej-}, g$\rightarrow$j devant \esp{a}).
    \item \esp{\textbf{Ningún} político usa la desinformación como arma.} — \esp{Ningún} (apocope de \esp{ninguno} devant nom masculin) + verbe à la forme affirmative, sans \esp{no} quand le mot négatif précède le verbe.
    \item \esp{\textbf{Se investiga} el origen del bulo (por la policía).} — Passif réfléchi : verbe au singulier, accordé avec \esp{el origen}.
    \item \esp{\textbf{Si hubiera} más transparencia, me gustaría.} — \esp{Si} + imparfait du subjonctif (\esp{hubiera}, de \esp{haber}) + conditionnel dans la principale. Variante complète : \esp{Si hubiera más transparencia, la situación sería mejor.}
\end{enumerate}
\textbf{Barème indicatif :} 1 point par transformation réussie (6 points). \textbf{Point de vigilance :} au passif réfléchi, l'accord du verbe avec le sujet patient est l'erreur la plus fréquente des francophones (\esp{se difunden bulos}, pas \esp{se difunde bulos}).
\end{corrige}

\begin{corrige}[Exercice 7 — Texte à trous : analyse publicitaire]
\begin{enumerate}
    \item \esp{emisor} — la coopérative « Café Excelencia » est l'émetteur du message.
    \item \esp{destinatario} — les consommateurs urbains sont le public cible (\esp{público objetivo}).
    \item \esp{mensaje} — le contenu : qualité, commerce équitable, fierté nationale.
    \item \esp{eslogan} — \esp{« Tu café, tu tierra, tu orgullo »}.
    \item \esp{tu} — adjectif possessif répété trois fois (anaphore) pour créer la proximité.
    \item \esp{imagen} — couleurs vert et or : nature et richesse.
    \item \esp{manipulación} — \esp{no hay manipulación oculta}.
    \item \esp{verificables} — adjectif au pluriel, accordé avec \esp{datos}.
    \item \esp{responsable} — \esp{publicidad responsable} : publicité éthique et vérifiable.
\end{enumerate}
\textbf{Barème indicatif :} 1 point par mot (9 points). \textbf{Point de vigilance :} les lettres données imposent la forme exacte ; \esp{verificables} sans \esp{s} finale est une demi-erreur.
\end{corrige}

\begin{corrige}[Exercice 8 — Appariements]
\begin{center}
\begin{tabularx}{\linewidth}{|l|X|}
\hline
\rowcolor{popblueL} \textbf{Concepto} & \textbf{Correspondance} \\
\hline
1. Censura previa & \textbf{B} — \esp{Prohibición de publicar antes de salir a la luz} \\
\hline
2. Difamación & \textbf{C} — \esp{Dañar la fama con informaciones falsas (sin imputar delito)} \\
\hline
3. Calumnia & \textbf{D} — \esp{Imputar falsamente un delito concreto} \\
\hline
4. Bulo & \textbf{E} — \esp{Falsa noticia viral: « El agua cura el cáncer en 3 días »} \\
\hline
5. Eslogan & \textbf{A} — \esp{¡Compra ya! Oferta limitada} \\
\hline
6. Desinformación & \textbf{F} — \esp{Campaña coordinada... sin pruebas} \\
\hline
\end{tabularx}
\end{center}
\textbf{Justification :} la paire 2/3 repose sur la distinction juridique clé (imputation d'un délit = calomnie) ; l'exemple E est un \esp{bulo} sanitaire typique (fausse promesse médicale virale) ; F décrit une campagne \textbf{coordonnée}, marque de la désinformation.
\textbf{Barème indicatif :} 1 point par paire (6 points).
\end{corrige}

\begin{corrige}[Exercice 9 — Commentaire de texte type Bac]
\textbf{Barème indicatif (20 points) :} compréhension globale 3 pts ; vocabulaire en contexte 3 pts ; analyse linguistique 4 pts ; argumentation 3 pts ; comparaison contextuelle 2 pts ; expression personnelle 5 pts (dont 2 pour la correction de la langue).

\begin{enumerate}
    \item \textbf{Compréhension globale (modèle) :} \esp{El texto analiza el doble efecto de las redes sociales sobre la libertad de expresión: democratizan la palabra pública y permiten movilizaciones ciudadanas, pero también facilitan la difusión de bulos y discursos de odio, amplificados por algoritmos que premian lo sensacional. Frente a la regulación europea y a los cortes de internet como el de Camerún, el autor defiende la educación mediática como solución, porque la libertad de expresión exige responsabilidad.}
    \item \textbf{Vocabulaire en contexte :}
    \begin{itemize}
        \item[a)] \esp{Saltarse los filtros tradicionales} : publicar directamente sin pasar por el control profesional de los periodistas y redactores que verificaban la información antes de difundirla.
        \item[b)] \esp{Luces y sombras} : aspectos positivos y negativos a la vez; ventajas e inconvenientes.
        \item[c)] \esp{Censura privatizada} : situación en la que una empresa privada (una plataforma digital) decide por sí misma qué contenidos se eliminan, actuando como juez de la verdad sin control judicial.
    \end{itemize}
    \item \textbf{Analyse linguistique :}
    \begin{itemize}
        \item \textbf{Subjonctifs :} \esp{sepan} (\esp{ciudadanos críticos que sepan verificar}) : subordonnée relative à antécédent indéfini/visée $\rightarrow$ valeur de but et de qualité recherchée. \esp{se reduzca} (\esp{so pena de que la democracia se reduzca}) : locution \esp{so pena de que} (crainte d'une conséquence) $\rightarrow$ subjonctif obligatoire.
        \item \textbf{Passif réfléchi :} \esp{Se comparten millones de contenidos sin verificar cada minuto.} $\rightarrow$ Voix active : \esp{Los usuarios comparten millones de contenidos sin verificar cada minuto.} (Accord : \esp{comparten} au pluriel, sujet patient \esp{contenidos}.)
    \end{itemize}
    \item \textbf{Argumentation :} Premier danger : la propagation des \esp{bulos}, \esp{discursos de odio y teorías de la conspiración} due à l'\esp{ausencia de verificación previa} (l. 5-6). Deuxième danger : les algorithmes qui \esp{premian lo sensacional y lo polarizante, no lo riguroso} (l. 6-7), si bien que \esp{la desinformación viaja seis veces más rápido que la verdad verificada} (l. 7-8).
    \item \textbf{Comparaison contextuelle (modèle) :} \esp{En Europa, el Estado regula mediante leyes como la Ley de Servicios Digitales, que obliga a las plataformas a retirar contenidos ilegales, aunque con el riesgo de una censura privatizada. En Camerún, en cambio, el poder cortó directamente internet en las regiones anglófonas en 2017-2018 para silenciar a la población: es una forma de represión, no de regulación.}
    \item \textbf{Expression personnelle (modèle) :} \esp{Creo que la educación mediática es necesaria pero no suficiente. A largo plazo, es la única vacuna sólida: un ciudadano formado verifica las fuentes, detecta los titulares sensacionalistas y no comparte sin leer. En Camerún, los talleres de verificación y las iniciativas de fact-checking como Dubawa muestran el camino. Sin embargo, la velocidad y la escala de la desinformación organizada —bots, cuentas falsas, microtargeting— superan la capacidad del individuo aislado. Por eso hacen falta también leyes estrictas pero precisas, como la Ley Europea de Servicios Digitales, que obliguen a las plataformas a la transparencia algorítmica y a retirar rápidamente los contenidos ilegales. El peligro es que leyes mal redactadas sirvan de pretexto para censurar: la definición de lo ilegal debe ser clara y controlada por jueces independientes. En síntesis, educación para los ciudadanos y regulación inteligente para las plataformas: las dos piernas de una democracia sana.}
\end{enumerate}
\textbf{Points de vigilance :} citer les lignes pour les questions 4 et 5 ; au résumé, ne pas recopier le texte ; dans l'expression personnelle, annoncer clairement sa position (\esp{Creo que...}, \esp{En mi opinión...}) et conclure.
\end{corrige}

\begin{corrige}[Exercice 10 — Thème et Version]
\textbf{A. Thème (proposition de traduction) :}

\begin{textebox}[Version espagnole]
La lucha contra las noticias falsas no debe servir de pretexto para amordazar a la prensa. En varios países africanos se han aprobado leyes sobre «cibercriminalidad», pero su redacción vaga permite procesar a periodistas por simples críticas al poder. La comunidad internacional recuerda que la difamación no debe ser un delito penal, sino una cuestión civil. En Camerún, las organizaciones de la sociedad civil exigen la despenalización de los delitos de prensa y la creación de un órgano de regulación independiente. Solo un marco jurídico claro, respetuoso de los estándares internacionales, garantizará una información libre y responsable.
\end{textebox}

\textit{Points de vigilance :} « museler » = \esp{amordazar} ; « poursuivre (judiciairement) » = \esp{procesar} (faux ami : \esp{perseguir} = poursuivre physiquement) ; « dépénalisation » = \esp{despenalización} ; « organe » = \esp{órgano} (accent !) ; \esp{solo... garantizará} : futur pour une vérité d'engagement.

\textbf{B. Version (proposition de traduction) :}

\begin{textebox}[Texte source espagnol]
La publicidad política en redes sociales plantea dilemas éticos inéditos. El microtargeting permite enviar mensajes diferentes a cada votante según su perfil psicológico, fragmentando el debate público en burbujas opacas. ¿Quién controla la veracidad de estos anuncios? Las plataformas se lavan las manos alegando que no son medios, pero su poder de amplificación es masivo. Urge adoptar una legislación que exija transparencia sobre la financiación y el contenido de la propaganda política digital, so pena de que la democracia se reduzca a una subasta de datos personales.
\end{textebox}

\textbf{Traduction française :}
La publicité politique sur les réseaux sociaux soulève des dilemmes éthiques inédits. Le microciblage permet d'envoyer des messages différents à chaque électeur selon son profil psychologique, fragmentant le débat public en bulles opaques. Qui contrôle la véracité de ces annonces ? Les plateformes s'en lavent les mains en prétendant qu'elles ne sont pas des médias, mais leur pouvoir d'amplification est massif. Il est urgent d'adopter une législation qui exige la transparence sur le financement et le contenu de la propagande politique numérique, sous peine de voir la démocratie se réduire à une vente aux enchères de données personnelles.

\textit{Points de vigilance :} \esp{microtargeting} = microciblage ; \esp{lavarse las manos} = s'en laver les mains (expression figée à traduire par son équivalent, pas mot à mot) ; \esp{urge} = « il est urgent » (verbe défectif à la 3\up{e} personne) ; \esp{so pena de que} = « sous peine de » ; \esp{subasta} = vente aux enchères.

\textbf{Barème indicatif :} 10 points par partie ; en thème, toute faute de mode (subjonctif après \esp{exigen que}, \esp{permite}) coûte 0,5 point.
\end{corrige}

\begin{corrige}[Exercice 11 — Sujets de rédaction type Bac]
\textbf{Barème indicatif (20 points) :} respect du sujet et du plan 4 pts ; richesse de l'argumentation et exemples 6 pts ; vocabulaire thématique 4 pts ; correction grammaticale (modes, accords, orthographe) 4 pts ; cohérence et connecteurs 2 pts.

\textbf{Sujet 1 — Plan attendu :}
\begin{enumerate}
    \item Introduction : la liberté d'expression, pilier démocratique, est-elle absolue ? Problématique annoncée.
    \item Axe 1 : les limites légitimes (incitation à la haine, à la violence, \esp{difamación}, \esp{calumnia}) protègent d'autres droits : l'honneur, la sécurité.
    \item Axe 2 : le danger des limites abusives (lois vagues sur la « cybercriminalité », censure préalable, journalistes poursuivis).
    \item Conclusion : des limites oui, mais exceptionnelles, précises, contrôlées par des juges indépendants.
\end{enumerate}

\textbf{Sujet 2 — Plan attendu :}
\begin{enumerate}
    \item Introduction : les réseaux, nouvel agora démocratique mais aussi vecteur de \esp{bulos}.
    \item Axe 1 : responsabilité des plateformes (transparence des algorithmes, modération).
    \item Axe 2 : responsabilité des citoyens (éducation aux médias, vérification avant de partager).
    \item Axe 3 : un cadre juridique clair et garant des libertés.
    \item Conclusion : ni jungle numérique ni censure : liberté responsable.
\end{enumerate}

\textbf{Modèle de rédaction (sujet 2, environ 160 mots) :}

\begin{textebox}[Redacción modelo]
Las redes sociales han democratizado el derecho a informar: hoy cualquier ciudadano puede denunciar un abuso o movilizar a la sociedad, como mostraron los movimientos ciudadanos de los últimos años. Sin embargo, esa misma libertad sin filtros ha multiplicado los bulos, el discurso de odio y la manipulación, porque los algoritmos premian lo sensacional y no lo riguroso. ¿Cómo conciliar libertad y responsabilidad? En primer lugar, las plataformas deben asumir su poder: transparencia algorítmica, moderación eficaz y retirada rápida de contenidos ilegales, como exige la legislación europea. En segundo lugar, los ciudadanos necesitan educación mediática desde la escuela: aprender a verificar las fuentes, detectar los titulares engañosos y no compartir sin leer; en Camerún, los programas de fact-checking son un ejemplo a seguir. Por último, el Estado debe legislar con precisión, castigando los delitos graves sin usar leyes vagas que criminalicen la crítica. En definitiva, no hay libertad duradera sin responsabilidad compartida: la solución no es menos libertad, sino más ciudadanía crítica y leyes que protejan el debate público y no al poder.
\end{textebox}

\textbf{Points de vigilance :} éviter les calques du français (\esp{las redes sociales han democratizado}, pas \esp{han democratisé}) ; soigner les connecteurs (\esp{sin embargo}, \esp{en primer lugar}, \esp{por último}, \esp{en definitiva}) ; utiliser au moins un subjonctif (\esp{leyes que protejan}) et un exemple concret (Cameroun, Europe) ; respecter le volume demandé (130-180 mots).
\end{corrige}

\newpage

\coursec{Fiche bilan}

\begin{popfigure}
\begin{tikzpicture}[mindmap, grow cyclic, every node/.style={concept, text=white, font=\small\bfseries, align=center},
  root concept/.append style={concept color=popblue, font=\small\bfseries, minimum size=3.1cm},
  level 1/.append style={level distance=3.5cm, sibling angle=51.4},
  level 1 concept/.append style={minimum size=2.3cm}]
\node[root concept]{Libertad de expresión y delitos de prensa}
  child[concept color=poporange]{ node[align=center]{Texto support\\ ¿derecho sin\\ límites?} }
  child[concept color=popgreen]{ node[align=center]{Lexique\\ difamación\\ calumnia, bulo} }
  child[concept color=poppurple]{ node[align=center]{Publicidad\\ eslogan\\ persuasión} }
  child[concept color=poppink]{ node[align=center]{Fake news\\ verificar\\ la fuente} }
  child[concept color=popgold]{ node[align=center]{Gramática\\ subjuntivo\\ estilo indirecto} }
  child[concept color=popblue!70]{ node[align=center]{Comentario\\ de texto\\ método} }
  child[concept color=popdark]{ node[align=center]{Cultura\\ de paz\\ medios} };
\end{tikzpicture}
\legende{Carte mentale de la leçon : liberté d'expression, délits de presse, publicité et fausses informations.}
\end{popfigure}

\begin{bilan}
\textbf{1. Lexique clé (à connaître par cœur)}
\begin{itemize}
  \item \esp{la libertad de expresión} : la liberté d'expression ; \esp{la libertad de prensa} : la liberté de la presse.
  \item \esp{la difamación} : la diffamation ; \esp{la calumnia} : la calomnie ; \esp{el delito de prensa} : le délit de presse.
  \item \esp{la censura} : la censure ; \esp{censurar} : censurer ; \esp{el periodista} : le journaliste.
  \item \esp{el bulo} : la fausse information, le canular ; \esp{la desinformación} : la désinformation.
  \item \esp{la publicidad} : la publicité ; \esp{el eslogan} : le slogan ; \esp{el anuncio} : l'annonce publicitaire.
  \item \esp{la manipulación} : la manipulation ; \esp{manipular} : manipuler ; \esp{la fuente} : la source.
  \item \esp{verificar / comprobar} : vérifier ; \esp{la veracidad} : la véracité ; \esp{el titular} : le titre (d'un article).
  \item \esp{el derecho} : le droit ; \esp{el límite} : la limite ; \esp{la responsabilidad} : la responsabilité.
\end{itemize}

\textbf{2. Grammaire à revoir pour cette leçon}
\begin{center}
\begin{tabularx}{\linewidth}{|X|X|X|}
\hline
\rowcolor{popblueL} Point de grammaire & Emploi dans la leçon & Exemple \\
\hline
Subjonctif après \esp{es importante que}, \esp{no creo que} & Exprimer un jugement sur les médias & \esp{Es importante que verifiques la fuente.} (Il est important que tu vérifies la source.) \\
\hline
Style indirect (\esp{dijo que...}) & Rapporter les propos d'un article & \esp{El periodista dijo que la noticia era falsa.} \\
\hline
\esp{por / para} & \esp{luchar por la libertad} (cause) ; \esp{ley para proteger} (but) & \esp{una ley para limitar los bulos} \\
\hline
\esp{ser / estar} & \esp{la noticia es falsa} (caractéristique) ; \esp{está censurada} (résultat) & \esp{El periódico está cerrado.} \\
\hline
Connecteurs d'argumentation & Structurer le commentaire & \esp{sin embargo, por lo tanto, en conclusión} \\
\hline
\end{tabularx}
\end{center}

\textbf{3. Méthodes express}
\begin{itemize}
  \item \textbf{Commentaire de texte} : 1) lire deux fois et souligner les mots clés ; 2) présenter (auteur, source, thème) ; 3) dégager la thèse de l'auteur ; 4) analyser les arguments et les procédés ; 5) donner son avis argumenté en conclusion.
  \item \textbf{Analyser une publicité} : identifier le produit, la cible (\esp{el público objetivo}), le slogan, l'image, les procédés de persuasion (émotion, humour, promesse, témoignage).
  \item \textbf{Repérer une fake news} : vérifier la source et l'auteur, contrôler la date, comparer avec d'autres médias, se méfier des titres sensationnalistes (\esp{¡Increíble!}, \esp{¡Urgente!}), vérifier les images.
\end{itemize}
\end{bilan}

\begin{aretenir}[Checklist avant le Bac]
\begin{itemize}
  \item $\square$ Je connais le lexique : \esp{difamación}, \esp{calumnia}, \esp{censura}, \esp{bulo}, \esp{desinformación}, \esp{eslogan}, \esp{manipulación}.
  \item $\square$ Je sais distinguer \esp{difamación} (attaque à l'honneur) et \esp{calumnia} (accusation d'un délit faux).
  \item $\square$ Je sais présenter un texte de presse : type, source, thème, thèse de l'auteur.
  \item $\square$ Je maîtrise le subjonctif après les verbes d'opinion niés (\esp{no creo que} + subjonctif).
  \item $\square$ Je sais rapporter une parole au style indirect avec la concordance des temps.
  \item $\square$ Je distingue \esp{por} (cause) et \esp{para} (but) dans mes productions.
  \item $\square$ Je sais analyser une publicité : produit, cible, slogan, procédés de persuasion.
  \item $\square$ Je sais appliquer les cinq réflexes pour détecter une fake news.
  \item $\square$ J'utilise des connecteurs : \esp{sin embargo}, \esp{además}, \esp{por lo tanto}, \esp{en conclusión}.
  \item $\square$ Je peux donner mon avis argumenté sur la liberté d'expression et ses limites.
  \item $\square$ Je fais attention aux faux amis : \esp{la publicidad} = la publicité (pas « publicité » au sens de notoriété) ; \esp{un delito} = un délit (pas « délit » au sens juridique français uniquement).
  \item $\square$ Je relis ma copie : accents (\esp{expresión}, \esp{información}), signes ¿ ¡, accords.
\end{itemize}
\end{aretenir}

\findecours

\end{document}
